\documentclass[11pt,oneside,openany]{book}
\usepackage[T1]{fontenc}
\usepackage{lmodern}
\usepackage{microtype}
\usepackage[a4paper,margin=28mm,headheight=14pt]{geometry}
\usepackage{amsmath,amssymb,amsthm,mathtools}
\usepackage[hidelinks,bookmarksopen,bookmarksnumbered]{hyperref}
\hypersetup{pdftitle={Classical and Weak Linear Partial Differential Equations},
  pdfauthor={MathGenome},pdfsubject={Unified notes, P2-01}}
\newcommand{\R}{\mathbb R}
\newcommand{\C}{\mathbb C}
\newcommand{\N}{\mathbb N}
\newcommand{\norm}[1]{\lVert #1\rVert}
\newcommand{\pair}[2]{\langle #1,#2\rangle}
\newcommand{\dashint}{\mathop{\overline{\int}}\nolimits}
\DeclareMathOperator{\Div}{div}
\DeclareMathOperator{\supp}{supp}
\DeclareMathOperator{\Tr}{Tr}
\allowdisplaybreaks[1]
\setlength{\emergencystretch}{2em}
\theoremstyle{plain}
\newtheorem{theorem}{Theorem}[chapter]
\newtheorem{lemma}[theorem]{Lemma}
\newtheorem{proposition}[theorem]{Proposition}
\newtheorem{corollary}[theorem]{Corollary}
\theoremstyle{definition}
\newtheorem{definition}[theorem]{Definition}
\newtheorem{example}[theorem]{Example}
\theoremstyle{remark}
\newtheorem{remark}[theorem]{Remark}

\begin{document}
\frontmatter
\hypersetup{pageanchor=false}
\begin{titlepage}
\thispagestyle{empty}
\vspace*{31mm}
{\large MathGenome \enspace / \enspace P2-01\par}
\vspace{7mm}
\hrule height 0.8pt
\vspace{11mm}
{\Huge\bfseries Classical and Weak Linear\\[2mm]
Partial Differential Equations\par}
\vspace{13mm}
{\large Unified notes\par}
\vspace{4mm}
{\normalsize Potential theory, elliptic estimates, and spectral realization\\
Heat, transport, wave, and integral equations\par}
\vfill
{\large 8 November 2026\par}
\end{titlepage}
\hypersetup{pageanchor=true}
\tableofcontents

\chapter*{Scope and conventions}
\addcontentsline{toc}{chapter}{Scope and conventions}
The notes develop classical and weak solution methods for linear
partial differential equations. After the common framework, we treat
harmonic and subharmonic functions, the weak and classical elliptic
Dirichlet theories, heat and parabolic evolution, transport and waves,
and two classical integral equations. The source map tracks 42 writing
groups and the 40 original mathematical obligations.

Spatial domains are open subsets of \(\mathbb R^n\), \(n\ge 1\).
Time intervals used for trace theorems have positive length.
The fundamental solution of the Laplacian is normalized by
\(-\Delta\Phi=\delta_0\), and the normal on an outer boundary is
outward. The Fourier convention is
\(\widehat f(\xi)=\int e^{-2\pi i x\cdot \xi}f(x)\,dx\).
We use \(H^{-1}(\Omega)=(H^1_0(\Omega))'\) with the continuous dual
and the pivot \(L^2\) pairing. On a Lipschitz boundary the
\(H^{-1/2}\) pairing is with the trace space \(H^{1/2}\).
Real forms suffice for the weak PDE arguments. Complexification
is available where an operator spectrum is used, with conjugation
in the second test factor of a sesquilinear energy form. Statements
about a real maximum use real solutions.

We use standard measure and integration theory, including Bochner
integration, and standard functional analysis, including uniform
boundedness, weak compactness in reflexive spaces, and Hilbert and
Radon measure representation. We also use classical ordinary
differential equations, Lebesgue-space duality, density and
mollification, the \(L^p\) Fourier and Riesz-transform tools stated at
their uses, and the Lipschitz trace, lifting, and compact embedding
theorems from \cite{P102}. The general Hilbert form theory is in
\cite{P106}; the compact operator and generator results used below
are recorded in \cite{D16}. The flow arguments use \cite{D13}.
All applications of these prerequisites specify the relevant domains,
norms, data, and coefficients.

Evans provides the main route through the classical models and
weak methods. Armitage--Gardiner and Gilbarg--Trudinger guide
potential theory and elliptic estimates. Grisvard supplies the
\(C^{1,1}\) boundary setting. Taylor, Renardy--Rogers, Brezis and
Lieberman complement the weak and classical evolution arguments,
and Engel--Nagel guides the semigroup formulation. John and Strauss
guide the characteristic and wave representations, and Kress and
Brunner the integral equations.

\mainmatter
\chapter{The Common PDE Framework}
\label{chap:pde-common}

\section{Order, principal symbols, and evolution variables}
\label{pde:C01}

For a multi-index \(\alpha=(\alpha_1,\ldots,\alpha_n)\), write
\(|\alpha|=\alpha_1+\cdots+\alpha_n\) and
\(\partial^\alpha=\partial_1^{\alpha_1}\cdots\partial_n^{\alpha_n}\).
A scalar linear differential operator of order \(m\) has the form
\(Lu=\sum_{|\alpha|\le m}a_\alpha(x)\partial^\alpha u\), with at least
one order-\(m\) coefficient nonzero somewhere. We use the principal
symbol
\[
 \sigma_m(L)(x,\xi)=\sum_{|\alpha|=m}a_\alpha(x)(i\xi)^\alpha .
\]
Indeed substituting \(u(x)=e^{ix\cdot \xi}\) gives
\(e^{-ix\cdot\xi}Lu(x)=\sum a_\alpha(x)(i\xi)^\alpha\), whose
homogeneous term of degree \(m\) is the symbol. Under the Fourier
convention of these notes the actual multiplier in the constant
coefficient case is \(\sigma_m(L)(2\pi\xi)\) at principal order.
For a scalar equation ellipticity means \(\sigma_m(L)(x,\xi)\ne0\)
for every \(\xi\ne0\) at each point \(x\).

A linear equation \(Lu=f\) has linear operator \(L\) and prescribed
source \(f\). In a semilinear equation the top-order linear
coefficients depend on the independent variables and lower-order
terms may depend nonlinearly on the unknown jet. In a quasilinear
equation the top-order derivatives enter linearly, with coefficients
allowed to depend on lower-order jets. In a fully nonlinear equation
dependence on the highest-order derivatives may itself be nonlinear.
For instance, \(\Delta u+u^3=0\), \((1+u^2)\Delta u=0\), and
\(\det D^2u=1\) for \(n\ge2\) illustrate the respective allowed
dependences.

\begin{proposition}[Real quadratic principal parts]
\label{prop:pde-symbol-types}
For a real symmetric second-order principal matrix \(a\), scalar
ellipticity is equivalent to definiteness of \(q(\xi)=\xi^{\mathsf T}a\xi\).
For an evolution operator with one positive and all remaining negative
directions in its quadratic form, a covector \(e\) with \(q(e)>0\)
defines a time direction in which \(q(\eta+\tau e)=0\) has two
distinct real roots for every covector \(\eta\) not proportional to \(e\).
\end{proposition}
\begin{proof}
The second-order symbol is \(-q\). A zero eigenvalue violates
ellipticity. Eigenvalues of opposite signs give a nonzero null
vector by taking an appropriate linear combination of their
orthogonal eigenvectors. Definiteness plainly excludes such vectors.
For the evolution assertion, write \(\eta=\eta_\perp+\beta e\),
where \(\eta_\perp\) is orthogonal to \(e\) for the bilinear form
of \(q\). On that complement the form is negative definite,
as follows by a diagonalization of the restriction and the
assumption on the index. Thus
\(q(\eta+\tau e)=q(\eta_\perp)+(\beta+\tau)^2q(e)\)
has exactly the asserted roots.
\end{proof}

For our positive sign convention, uniform ellipticity of the
principal matrix means \(q_x(\xi)\ge\theta|\xi|^2\) for a single
\(\theta>0\) throughout the region. A uniformly negative definite
matrix is put in this convention by changing the overall sign.
The constant-speed wave equation
has quadratic form \(\tau^2-c^2|\xi|^2\). The heat operator has a
first-order time derivative and an elliptic second-order spatial
part. Its parabolic character uses that choice of time derivative
and the sign of spatial diffusion. A zero eigenvalue of the total
quadratic part in space-time alone does not give parabolic evolution.
For classical classification and the basic examples see
\cite[Sections 1.1--1.3]{Evans} and \cite[Chapter 1]{Strauss}.

\section{Distributions, Green identities, and the normal trace}
\label{pde:C02}

For locally integrable \(F\) the distributional convention is
\(\langle-\Div F,\varphi\rangle=\int F\cdot\nabla\varphi\) for
\(\varphi\in C_c^\infty(\Omega)\). The dualities below are real
in this section. We write \(\Tr\) for the usual trace on a bounded
Lipschitz open set. Its map \(H^1(\Omega)\to H^{1/2}(\partial\Omega)\)
is bounded and onto with a bounded right inverse, and its kernel
is \(H^1_0(\Omega)\) \cite[Unit U5]{P102}. In dimension one the
boundary consists of finitely many endpoints, and we take
\(H^{1/2}(\partial\Omega)=\ell^2(\partial\Omega)\) with counting
measure. The one-dimensional endpoint trace has the same kernel,
and affine interpolation on each interval is a bounded linear
right inverse.

\begin{proposition}[The interior equation and its test functions]
\label{prop:pde-interior-testing}
If \(A\in C^1(\Omega;\mathbb R^{n\times n})\),
\(u\in C^2(\Omega)\) and \(f\in C(\Omega)\), then
\(-\Div(A\nabla u)=f\) pointwise if and only if
\[
 \int_\Omega A\nabla u\cdot\nabla\varphi\,dx
     =\int_\Omega f\varphi\,dx\qquad(\varphi\in C_c^\infty(\Omega)).
\]
If instead \(F\in L^2(\Omega;\mathbb R^n)\), \(f\in H^{-1}(\Omega)\)
and \(-\Div F=f\) distributionally, the equality
\(\int F\cdot\nabla v=\langle f,v\rangle\) holds for every
\(v\in H^1_0(\Omega)\).
\end{proposition}
\begin{proof}
Integration by parts on the support of the test gives the forward
implication. For the reverse, the continuous function
\(-\Div(A\nabla u)-f\) induces the zero distribution.
If it were nonzero at a point, one sign would be strict on a
neighborhood and a nonnegative test supported there would give
a nonzero integral. The last assertion follows by density in
\(H^1_0\) and continuity of both functionals.
\end{proof}

The restriction to interior tests imposes no boundary value.
For example, on a ball both the zero function and any constant
function satisfy the homogeneous Laplace equation in these tests.

\begin{theorem}[Normal trace and weak Green identities]
\label{thm:pde-normal-trace}
On a bounded Lipschitz open set, if \(F\in L^2(\Omega;\mathbb R^n)\)
and \(\Div F\in L^2(\Omega)\), there is a unique
\(\gamma_\nu F\in H^{-1/2}(\partial\Omega)\) with
\[
 \langle\gamma_\nu F,\Tr v\rangle
 =\int_\Omega\{F\cdot\nabla v+(\Div F)v\}\,dx\qquad(v\in H^1).
\]
The norm of this trace is bounded by a domain constant times
\(\|F\|_2+\|\Div F\|_2\). It agrees with the outward flux
for fields smooth up to the boundary. If \(u,v\in H^1(\Omega)\) and both Laplacians
are in \(L^2\), then
\[
 \int_\Omega(u\Delta v-v\Delta u)\,dx
   =\langle\gamma_\nu\nabla v,\Tr u\rangle
      -\langle\gamma_\nu\nabla u,\Tr v\rangle .
\]
\end{theorem}
\begin{proof}
The functional on the right of the first identity is bounded on
\(H^1\) by Cauchy--Schwarz. It vanishes on compactly supported smooth
functions by the definition of divergence, hence on \(H^1_0\).
Its value on a trace \(g\) is therefore independent of the chosen
lift, and the bounded lifting gives the asserted dual norm.
Surjectivity of the trace gives uniqueness. For a field and a test
smooth up to the boundary, we compute the integral of \(\Div(vF)\)
in local graph charts. On a Lipschitz epigraph \(x_n>\psi(x')\), the
Lipschitz graph area theorem of \cite[Theorem 4.16]{P102} gives
\(dS=\sqrt{1+|\nabla\psi|^2}\,dx'\) on the graph, with
Euclidean-normalized Hausdorff surface measure. The outward
normal is \((\nabla\psi,-1)/\sqrt{1+|\nabla\psi|^2}\) at
differentiability points of \(\psi\), since the inward side
is above the graph. Thus the geometric normal times surface
measure is exactly \((\nabla\psi,-1)\,dx'\) in these coordinates.
Write \(W\) for the product \(vF\) localized by a cutoff supported
in the chart. The vertical integral of \(\partial_nW_n\) is
the negative boundary value. For \(j<n\), the function
\(A_j(x')=\int_{\psi(x')}^\infty W_j(x',t)\,dt\) is absolutely
continuous on coordinate lines, with
\[
 \partial_j A_j(x')=\int_{\psi(x')}^\infty\partial_j W_j(x',t)\,dt
             - W_j(x',\psi(x'))\partial_j\psi(x').
\]
The integral of \(\partial_j A_j\) is zero by compact support, leaving
the graph term \(W_j(x',\psi(x'))\partial_j\psi(x')\).
These terms combine to the stated geometric flux. A partition
of unity combines the graph computations, with interior terms
integrating to zero and partition-derivative terms summing to zero.
In dimension one the surface measure counts endpoints, and the
same assertion is the fundamental theorem of calculus on intervals.
Approximation in \(H^1\)
extends the identity to all tests. Taking \(F=\nabla v\) gives
Green's first identity with the indicated outward normal.
Subtracting the two first identities gives the second one.
\end{proof}

See \cite[Sections 5.2 and 6.1]{Evans} and \cite[Chapter 5, Section 7]{TaylorI}.
The normal trace is a dual boundary functional and need not have
a pointwise representative.

\section{Data, traces, and well-posedness}
\label{pde:C03}

A Cauchy problem prescribes values on an initial hypersurface,
including the needed initial time derivatives. A boundary-value
problem prescribes data on a spatial boundary. An initial-boundary
problem specifies both. For a classical problem, continuity and
differentiability are required wherever its pointwise equations or
boundary conditions are evaluated. In an energy problem the spatial
boundary data may instead be imposed by the Sobolev trace and
the initial data by a continuous time representative. In the
semigroup setting the corresponding definition of a mild solution
is given in Section~\ref{pde:C07}.

Fix topological data and solution spaces. Hadamard well-posedness
means existence for the stated data, uniqueness in the stated class,
and continuity of the solution map for those topologies.
These three claims concern the specified problem and function spaces
\cite[Section 1.3]{Evans}\cite[Section 1.1.5]{RR}.

\begin{proposition}[Energy control of differences]
\label{prop:pde-energy-stability}
Suppose two solutions \(u,v\) have difference energy \(E(t)\ge0\).
Assume \(E(t)\) controls \(\|u(t)-v(t)\|^2\) in the solution norm
and \(E(t)\le C E(0)\) on the time interval. If identical data
give \(E(0)=0\), the two solutions with identical data agree.
If \(E(0)\) is the squared data norm of the difference and
\(E(t)=\|u(t)-v(t)\|^2\), the stability constant is \(\sqrt C\).
\end{proposition}
\begin{proof}
Identical data force all the nonnegative energies to be zero,
so the solution norm of the difference is zero. Taking the
square root gives the second conclusion. If the energy seminorm
does not determine the whole state, an additional control is
needed before using this proposition. For instance, the wave
energy controls the gradient and velocity, and Section~\ref{pde:W03}
also estimates the \(L^2\) displacement.
\end{proof}

\begin{example}[Backward heat instability]
On the unit circle with periodic coordinate \(x\), the backward
heat equation \(u_t+\Delta u=0\) admits the smooth solutions
\[
 u_j(t,x)=e^{-4\pi^2j^2(T-t)}\cos(2\pi jx).
\]
The initial data tend to zero in every fixed Sobolev norm,
since every polynomial in \(j\) is dominated by the exponential.
At \(t=T\), distinct positive frequencies are orthogonal with \(L^2\)
norm \(1/\sqrt2\) each, so their mutual distance is \(1\). The terminal
values cannot converge. Thus no \(L^2\)-valued extension of the solution
map to the initial-data set with zero adjoined is continuous there in any
of these topologies. If the class already contains the zero solution,
its solution map to time \(T\) is discontinuous at zero.
\end{example}

\section{Localization and difference quotients}
\label{pde:C04}

\begin{proposition}[Distributional cutoff identity]
\label{prop:pde-cutoff}
For \(u\in\mathcal D'(\Omega)\), \(\chi\in C^\infty(\Omega)\),
\[
 \Delta(\chi u)=\chi\Delta u+2\nabla\chi\cdot\nabla u+(\Delta\chi)u
 \quad\text{in }\mathcal D'(\Omega).
\]
For compactly supported \(\chi\) the left-hand input \(\chi u\)
has compact support in \(\Omega\).
\end{proposition}
\begin{proof}
For a test function, differentiating the product in
\(\langle\partial_j(\chi u),\varphi\rangle
=-\langle u,\chi\partial_j\varphi\rangle\)
gives the one-derivative rule. Apply it again, sum over \(j\),
and note that the product distribution vanishes on tests supported
outside \(\supp\chi\).
\end{proof}

Fix a coordinate vector \(e_j\), and write
\(D_h w(x)=(w(x+h e_j)-w(x))/h\) for \(h\ne0\).
On a domain we evaluate this only where all the arguments belong
to the domain.

\begin{lemma}[Difference quotient estimates]
\label{lem:pde-difference-quotients}
On the whole space, or for translated tests supported away from
the boundary, discrete integration by parts is
\(\int (D_h w)\psi=-\int w D_{-h}\psi\).
If \(U\Subset U'\Subset\Omega\), \(w\in W^{1,p}(U')\),
\(1\le p<\infty\), and the segments from \(x\) to \(x+h e_j\) for
\(x\in U\) remain in \(U'\), then
\(\|D_h w\|_{L^p(U)}\le\|\partial_j w\|_{L^p(U')}\).
Conversely, for \(1<p<\infty\), if \(w\in L^p(U')\) and these
difference quotients are uniformly bounded in \(L^p(U)\)
for all sufficiently small \(h\), then \(\partial_j w\in L^p(U)\)
with that bound.
\end{lemma}
\begin{proof}
The change of variables \(y=x+h e_j\) in the translated integral
gives the discrete identity. For a smooth function,
\(D_hw(x)=\int_0^1 \partial_jw(x+\theta h e_j)\,d\theta\).
Minkowski's inequality and translation invariance give the bound.
For a general \(w\), approximate it in \(W^{1,p}(U')\) by smooth
functions, using interior Sobolev density \cite[Unit U4]{P102}.
For fixed \(h\) the two restriction and translation maps on \(L^p\)
have norm at most one since \(U\) and \(U+h e_j\) lie in \(U'\).
Thus the bound passes to the limit at each allowed \(h\).
For the converse take \(h_k\to0\) and a weakly convergent subsequence
of \(D_{h_k}w\) in \(L^p(U)\). The discrete identity on
\(C_c^\infty(U)\) and uniform convergence of the test difference
quotients identify the weak limit with the distributional derivative.
Weak lower semicontinuity gives the norm bound.
\end{proof}

This is the local version used in Section~\ref{pde:E05}.
Compare \cite[Section 5.8.2]{Evans}.

\section{Galerkin approximation and weak-limit testing}
\label{pde:C05}

\begin{theorem}[Coercive Galerkin approximation]
\label{thm:pde-cea}
Let \(V\) be a real Hilbert space, \(a\) a bilinear form with
\(|a(u,v)|\le M\|u\|\|v\|\) and \(a(v,v)\ge\alpha\|v\|^2\)
for \(\alpha>0\), and \(\ell\in V'\).
For any nested finite-dimensional \(V_m\subset V\) of dense union,
the equations \(a(u_m,v)=\ell(v)\) on \(V_m\) have unique solutions.
The full weak solution \(u\) satisfies
\[
 \|u-u_m\|\le (M/\alpha)\inf_{v\in V_m}\|u-v\|.
\]
In particular \(u_m\to u\) strongly.
\end{theorem}
\begin{proof}
Apply real Lax--Milgram on \(V\) and on each \(V_m\)
\cite[Chapter 3]{P106}. Subtraction gives
\(a(u-u_m,w)=0\) for \(w\in V_m\). With \(e=u-u_m\),
for any \(v\in V_m\) one has
\(\alpha\|e\|^2\le a(e,e)=a(e,u-v)\le M\|e\|\|u-v\|\).
If \(e=0\) the estimate is immediate, and otherwise division and
infimization prove it. Nestedness and dense union make the
best approximation error tend to zero.
\end{proof}

\begin{proposition}[Passing a linear equation to a weak limit]
\label{prop:pde-weak-galerkin}
Let \(V\) be a real reflexive Banach space with nested finite-dimensional
subspaces \(V_m\) of dense union, \(a\) a bounded bilinear form and
\(\ell\in V'\). Suppose \((u_m)\) is bounded in \(V\) and
\(a(u_m,v)=\ell(v)\) for every \(v\in V_m\).
A subsequence converges weakly to a solution of
\(a(u,v)=\ell(v)\) for all \(v\in V\). If \(a\) is coercive,
that solution is unique and is the only possible weak cluster point.
\end{proposition}
\begin{proof}
Reflexivity gives the weak subsequence
\cite[Chapter 3]{P106}. For a fixed \(v\in V_j\) all sufficiently
late equations include that test, and \(a(\,\cdot\,,v)\) is
a continuous linear functional. Pass to the limit and then use
the norm density of the test union. The same argument applies
to any weak cluster point of tails. Under coercivity, the difference
of two limits can be used as its own test and must vanish.
\end{proof}

\section{Hilbert triples, time traces, and energy rules}
\label{pde:C06}

Let \(V\hookrightarrow H\) be a dense continuous embedding of
real Hilbert spaces. The map \(h\mapsto(v\mapsto(h,v)_H)\)
identifies \(H\) with a subspace of \(V'\), with continuous
injective inclusion. A distributional time derivative in \(V'\)
is understood via Bochner integration against smooth scalar tests.

\begin{theorem}[Time trace and product rule]
\label{thm:pde-time-trace}
Suppose \(T>0\), \(u\in L^2(0,T;V)\), and its distributional
derivative \(u'\) belongs to \(L^2(0,T;V')\).
Then \(u\) has a unique \(C([0,T];H)\) representative. For that
representative,
\[
 \|u(t)\|_H^2-\|u(s)\|_H^2
       =2\int_s^t\langle u'(r),u(r)\rangle\,dr
       \qquad(0\le s\le t\le T).
\]
The trace map into \(C([0,T];H)\) is continuous for the sum of the
two \(L^2\) norms. If \(v\) has the same regularity, then
\((u,v)_H\) is absolutely continuous with derivative
\(\langle u',v\rangle+\langle v',u\rangle\) almost everywhere.
\end{theorem}
\begin{proof}
For a smooth \(V\)-valued function the first identity is
differentiation in \(H\). Choosing \(s\) with
\(\|u(s)\|_H^2\le T^{-1}\|u\|_{L^2(H)}^2\) yields the estimate
\[
 \sup_t\|u(t)\|_H^2
 \le T^{-1}\|u\|_{L^2(H)}^2+
       2\|u'\|_{L^2(V')}\|u\|_{L^2(V)} .
\]
For the general case first regard \(u\) as a \(V'\)-valued function.
The Bochner primitive of \(u'\) is absolutely continuous in \(V'\);
subtracting it leaves a function with distributional derivative
zero. To identify that difference, convolve it locally in time.
The convolutions have derivative zero and so are constant on
each smaller interval. Their \(L^2(V')\) convergence makes the
original difference constant there, with compatible constants on
overlaps. Thus \(u\) has an absolutely continuous \(V'\) representative.
Reflect it across each endpoint onto a slightly larger interval.
There is no derivative jump because the two traces at each
reflection agree in \(V'\). Reflection preserves \(L^2(V)\)
and \(L^2(V')\) derivative bounds locally. Convolution in time
now gives smooth \(V\)-valued \(u_\epsilon\) on \([0,T]\),
converging to \(u\) in \(L^2(V)\) with derivatives converging
in \(L^2(V')\). Apply the displayed estimate to their differences.
The smooth approximations converge in \(C(H)\), and the pairings
of their derivatives with themselves converge in \(L^1\)
by Cauchy--Schwarz. Pass to the limit for every \(s,t\).
The estimate gives trace continuity, and two continuous
representatives that agree almost everywhere agree everywhere
because \(T>0\). Polarize using \(u+v\) and \(u-v\) to obtain
the product formula. The integrability of its right-hand side
gives absolute continuity.
\end{proof}

Compare \cite[Section 5.9.2]{Evans}, \cite[Chapter 11]{RR}
and the corrected Hilbert-triple result in \cite{D16}.

\section{Semigroups and variation of constants}
\label{pde:C07}

Let \(S(t)\) be a strongly continuous semigroup on a real or
complex Banach space \(X\), with generator \(A\). We use its
local operator norm bound, dense generator domain and closedness.
These are the standard generator facts of
\cite[Chapter II]{EngelNagel}\cite{D16}.
Write \(D(A)\) with its complete graph norm \(\|x\|+\|Ax\|\).
For each \(h\ge0\), define
\[
 J(h)x=\int_0^h S(r)x\,dr.
\]
The difference quotient of \(S(k)J(h)x\) at \(k=0+\) is the
difference of the two averages over \([h,h+k]\) and \([0,k]\).
Hence \(J(h)x\in D(A)\) and \(AJ(h)x=S(h)x-x\).
Both \(J(h)x\) and \(AJ(h)x\) are continuous in \(h\);
on a fixed compact time interval the maps
\(J(h):X\to D(A)\) are uniformly bounded.

\begin{theorem}[Duhamel formula and the integrated equation]
\label{thm:pde-duhamel}
Let \(x_0\in X\) and \(f\in L^1(0,T;X)\).
There is a unique continuous function \(u:[0,T]\to X\) for which
\(\int_0^t u(s)\,ds\in D(A)\) and
\[
 u(t)=x_0+A\int_0^t u(s)\,ds+\int_0^t f(s)\,ds.
\]
It is the mild solution, given by
\[
 u(t)=S(t)x_0+\int_0^t S(t-s)f(s)\,ds .
\]
If \(x_0\in D(A)\) and \(f\in C^1([0,T];X)\), then
\(u\in C^1([0,T];X)\cap C([0,T];D(A))\) and
\(u'=Au+f,\ u(0)=x_0\). Every such classical solution obeys the formula.
\end{theorem}
\begin{proof}
The formula defines a continuous function. In the convolution use
strong continuity pointwise in \(s\) away from the moving endpoint,
and dominate by the local bound for \(S\) times \(\|f(s)\|\).
Bochner Fubini gives the time primitive of the displayed formula as
\[
 J(t)x_0+\int_0^t J(t-r)f(r)\,dr .
\]
The integral here exists in \(D(A)\) with its graph norm:
measurability follows by simple approximation of \(f\) and the
strong continuity of \(J(\,\cdot\,)\) into that norm for each
fixed vector. The uniform graph bound yields integrability.
Applying \(A\) now gives the integrated equation by the identity
for \(AJ\).

For uniqueness, the primitive \(w\) of the difference of two
integrated solutions has \(w'=Aw\), \(w(0)=0\), and is continuous
in the graph norm since both \(w\) and \(Aw=w'\) are continuous.
For fixed \(t>0\), differentiate \(S(t-r)w(r)\) for \(0<r<t\).
The differentiation rule on \(D(A)\) follows from the generator
limit and the semigroup law, so the derivative is
\(S(t-r)(w'(r)-Aw(r))=0\).
Continuity to the endpoints implies \(w(t)=S(t)w(0)=0\).

For classical forcing, put \(v(t)=u(t)-S(t)x_0\).
The Bochner fundamental theorem for \(f\) and Fubini give
\[
 v(t)=J(t)f(0)+\int_0^t J(t-r)f'(r)\,dr .
\]
This is graph-continuous. Its norm derivative equals
\(S(t)f(0)+\int_0^t S(t-r)f'(r)\,dr\), as follows from
\(J'(t)x=S(t)x\), dominated convergence and \(J(0)=0\).
The identity for \(AJ\) gives \(Av=v'-f\).
The homogeneous term is graph-continuous and differentiable
with derivative \(S(t)Ax_0\). Conversely any classical solution
can be integrated in the graph norm to give the integrated
equation, so uniqueness gives the formula.
\end{proof}

This proves the agreement between the formula definition and the
integrated definition of mild solution. See
\cite[Chapter VI, Section 7]{EngelNagel}.
\chapter{Laplace's Equation and Potential Theory}
\label{chap:pde-potential}

Write \(\omega_{n-1}=|S^{n-1}|\) for surface measure of the
unit sphere. The volume of \(B(0,r)\) is
\(\omega_{n-1}r^n/n\). In one dimension the sphere consists of
two endpoints, each with weight one. The notation \(\dashint\)
is the integral divided by the measure of the set of integration.

\section{Fundamental solutions and Newtonian potentials}
\label{pde:L01}

\begin{theorem}[Normalized fundamental solutions]
\label{thm:pde-laplace-fundamental}
The following locally integrable functions satisfy
\(-\Delta\Phi=\delta_0\) in distributions:
\[
 \Phi(x)=
 \begin{cases}
  -|x|/2, & n=1,\\
  -\dfrac1{2\pi}\log|x|, & n=2,\\[1mm]
  \dfrac{|x|^{2-n}}{(n-2)\omega_{n-1}}, & n\ge3.
 \end{cases}
\]
If \(f\in C_c^\infty(\mathbb R^n)\), then \(u=\Phi*f\) is smooth
and satisfies \(-\Delta u=f\). The distribution identity
\(-\Delta(\Phi*F)=F\) also holds for a compactly supported
distribution \(F\).
\end{theorem}
\begin{proof}
In dimension one the slopes of \(-|x|/2\) are \(1/2\) and
\(-1/2\) on the two sides of zero. Their jump is \(-1\);
integration by parts on the two half-lines gives the assertion.
For \(n\ge2\), direct differentiation of the radial functions
gives \(\Delta\Phi=0\) off zero and
\(\partial_r\Phi(r)=-1/(\omega_{n-1}r^{n-1})\).
Local integrability follows in polar coordinates, including
the logarithm at zero. Take a test \(\varphi\) supported inside
a large ball and apply Green's identity on the annulus
\(\epsilon<|x|<R\). The outer term is zero. On the inner sphere
the outward normal of the annulus is \(-\partial_r\), hence
\[
 \int_{|x|>\epsilon}\Phi(-\Delta\varphi)\,dx
 =\dashint_{|x|=\epsilon}\varphi\,dS
      -\int_{|x|=\epsilon}\Phi\,\partial_\nu\varphi\,dS.
\]
The first term tends to \(\varphi(0)\). The second is
\(O(\epsilon)\) in \(n\ge3\) and \(O(\epsilon|\log\epsilon|)\)
in \(n=2\), and the omitted volume term tends to zero by local
integrability. This proves the sign and normalization.

For smooth compact \(f\), in
\(\int\Phi(y)f(x-y)\,dy\) all derivatives can fall on \(f\).
On any compact \(x\)-set the integrals are over one bounded
\(y\)-set, and local integrability gives differentiation
by dominated convergence. For general compact \(F\),
convolution with a distribution is defined by the test pairing
with a cutoff equal to one around the relevant compact support,
and derivatives commute with it \cite[Unit U4]{P102}.
Consequently \(-\Delta(\Phi*F)=\delta_0*F=F\).
\end{proof}

In particular the two-dimensional normalization has a logarithmic
growth at infinity, and the identity is local in distributions.
See \cite[Section 2.2.1]{Evans}, \cite[Section 4.1]{GT},
and \cite[Section 4.3]{AG}.

\section{Mean values and Weyl's lemma}
\label{pde:L02}

A function is harmonic on an open set if it is \(C^2\) and
\(\Delta u=0\) there. All statements below are local on each
connected component.

\begin{theorem}[Mean values and distributional harmonicity]
\label{thm:pde-mean-weyl}
For a harmonic function on a neighborhood of
\(\overline B(a,r)\),
\[
 u(a)=\dashint_{\partial B(a,r)}u\,dS
     =\dashint_{B(a,r)}u\,dx .
\]
Conversely a continuous function satisfying either mean-value
identity whenever \(\overline B(a,r)\subset\Omega\) is harmonic.
Every distribution \(U\) with \(\Delta U=0\) is represented
by a smooth harmonic function.
\end{theorem}
\begin{proof}
Set \(M(s)=\omega_{n-1}^{-1}\int_{S^{n-1}}u(a+s\theta)\,dS_\theta\).
For \(s>0\), differentiation and the divergence theorem give
\[
 M'(s)=\frac{1}{\omega_{n-1}s^{n-1}}
                 \int_{B(a,s)}\Delta u\,dx .
\]
This also holds for \(n=1\) by the ordinary fundamental theorem of
calculus. If \(u\) is harmonic, \(M(s)\) is constant and its
limit at zero is \(u(a)\). Integrating the spherical mean in
the radial variable gives the ball mean.

A continuous function with the ball identity has the sphere
identity by differentiating the equality
\(\int_{B(a,s)}u=|B(a,s)|u(a)\) in \(s>0\). Thus it suffices to
start with sphere identities. Convolving with a smooth
radial unit-mass kernel of sufficiently small support preserves
the function locally, because the convolution is a radial
weighted average of the sphere means. It follows that the function
is smooth. Taylor expansion at each center shows that the
spherical mean has increment
\((s^2/(2n))\Delta u(a)+o(s^2)\): the linear and mixed
terms average to zero and the average of \(\theta_i^2\) is \(1/n\).
Hence \(\Delta u=0\).

For a distribution, first mollify at scale \(\epsilon\) on
relatively compact subregions. The functions \(U_\epsilon\)
are smooth harmonic where the convolution is defined, since
derivatives commute. On a smaller subregion fix \(s>0\) such
that all translations by a ball of radius \(2s\) stay inside
the original region. A fixed smooth radial unit kernel \(\rho_s\)
supported in the radius-\(s\) ball satisfies
\(U_\epsilon*\rho_s=U_\epsilon\) there for \(\epsilon<s\).
Pass to distributions, moving the convolution to test functions
whose translated supports remain inside. Thus \(U=U*\rho_s\)
on the smaller region. The convolution with this test kernel
is smooth, as follows by differentiating the translated kernel
in the distribution pairing. The local representatives agree on
overlaps because they induce the same distribution. Their
Laplacians vanish pointwise.
\end{proof}

This is Weyl's lemma in the form used below for harmonic remainders.
Compare \cite[Sections 1.2--1.3]{AG}, \cite[Sections 2.2 and 2.7]{GT},
and \cite[Section 2.2]{Evans}.

\section{Derivative bounds, Liouville, and analyticity}
\label{pde:L03}

\begin{proposition}[Interior derivative estimate and Liouville]
\label{prop:pde-harmonic-derivatives}
For harmonic \(u\) on a neighborhood of \(\overline B(a,r)\),
\[
 |\partial^\alpha u(a)|
 \le C_{n,\alpha}r^{-n-|\alpha|}
                      \int_{B(a,r)}|u(y)|\,dy .
\]
A real harmonic function on all of \(\mathbb R^n\) bounded
above or below is constant.
\end{proposition}
\begin{proof}
Choose a fixed smooth radial unit kernel supported in the unit
ball and scale it by \(r\). On a neighborhood of \(a\) the
convolution represents \(u\) by the mean-value identity.
Move the derivatives to the kernel, whose sup norm after
\(\alpha\) derivatives is bounded by
\(C_{n,\alpha}r^{-n-|\alpha|}\).
For Liouville, if \(u\le M\) apply the estimate with
\(|\alpha|=1\) to the nonnegative harmonic function \(M-u\).
The integral of that function on \(B(a,r)\) is its value at
\(a\) times the ball volume. Let \(r\to\infty\) to obtain
zero gradient everywhere. Apply the upper-bound case to \(-u\)
for a lower bound.
\end{proof}

\begin{corollary}[Analyticity and continuation]
\label{cor:pde-harmonic-analytic}
Harmonic functions are real analytic. Two harmonic functions
on a connected open set that agree on a nonempty open subset
agree on the whole set.
\end{corollary}
\begin{proof}
In dimension one a harmonic function is locally affine.
For \(n\ge2\), use the Poisson representation proved in
Section~\ref{pde:L05} on a sphere of radius \(r\) centered
at any interior point \(a\). For \(h=x-a\), the denominator
factor apart from \(r^n\) is
\[
 \bigl(1-2(\zeta-a)\cdot h/r^2+|h|^2/r^2\bigr)^{-n/2}.
\]
In a sufficiently small coordinate polydisc, the sum of the
absolute values of the coefficients of the nonconstant
polynomial inside the parentheses, evaluated at the radii,
is bounded by \(2\|h\|_1/r+\|h\|_1^2/r^2<1\).
The binomial series therefore converges absolutely as a
power series in \(h\), uniformly in \(\zeta\) on the sphere
on a smaller polydisc. Multiplying by the polynomial
numerator and integrating term by term proves analyticity.
For the continuation assertion consider the difference. The
set of points where all its derivatives, including the function,
are zero is relatively closed by continuity. It is open by
the convergent local Taylor series, and contains the given open
subset. Connectedness completes the proof.
\end{proof}

\section{Subharmonicity and the maximum principles}
\label{pde:L04}

\begin{definition}
For \(n\ge2\), a subharmonic function on \(\Omega\) is an upper
semicontinuous \(u:\Omega\to[-\infty,\infty)\) which is not
identically \(-\infty\) on any component and satisfies
\[
 u(a)\le \dashint_{\partial B(a,r)}u\,dS
 \qquad(\overline B(a,r)\subset\Omega).
\]
The extended integral is defined since \(u\) is bounded above
on each compact sphere.
\end{definition}

\begin{theorem}[Comparison and maximum principles]
\label{thm:pde-subharmonic-comparison}
The sphere submean condition can equivalently be replaced by
the ball submean condition on every such ball, or by the
following comparison condition: a harmonic function in a ball
\(D\Subset\Omega\), continuous on its closure, that dominates
the interior boundary upper limits of \(u\), dominates \(u\)
throughout \(D\).
If a subharmonic function attains its finite maximum at an
interior point of a component, it is constant on that component.
If \(\Omega\) is bounded and
\(\limsup_{x\to\zeta,\ x\in\Omega}u(x)\le M\) on \(\partial\Omega\),
then \(u\le M\). For real \(C^2\) functions the submean condition
is equivalent to \(\Delta u\ge0\).
\end{theorem}
\begin{proof}
The sphere condition implies the ball condition by polar
integration. On a compact ball, apply the nonnegative integral
form of Fubini to a common upper bound minus \(u\);
this justifies the assertion for extended values.
Suppose the ball condition holds and \(u(a)=m\) is a finite
maximum on a connected component. On every sufficiently
small ball about \(a\) the integral of \(m-u\) is zero.
If \(u<m\) somewhere on such a ball, upper semicontinuity
would give an open set where \(m-u\) is bounded below by
a positive constant. Thus the maximum set is open, and it is
relatively closed by upper semicontinuity. It is the entire
component.

On a bounded domain with the stated boundary control, any
value strictly greater than \(M\) would yield a finite interior
maximum of this kind. To see this, extend \(u\) to the boundary
using its upper limits, so the extension is upper semicontinuous
on the compact closure and has no value \(+\infty\).
An interior maximum would force constancy on its component,
and approaching the boundary of that bounded component
would contradict the boundary upper bound. The same
argument applied to \(u-h\) gives ball comparison, since
a harmonic \(h\) satisfies the mean-value equality.

For the reverse, fix a compact sphere in \(\Omega\).
Its upper semicontinuous values are decreasing limits of
continuous majorants. For example on that compact metric
sphere one can use
\(\max(-j,\sup_\eta(u(\eta)-j|\zeta-\eta|))\).
These functions are Lipschitz, decrease with \(j\), and
converge pointwise to \(u\) by upper semicontinuity and
the finite common upper bound. If all values on the sphere
are \(-\infty\), use the constants \(-j\).
Let \(h_j\) be the continuous-data Poisson solutions of
Section~\ref{pde:L05}. Upper semicontinuity of \(u\) in the
larger region gives the boundary upper-limit comparison with
\(h_j\). At the center, \(u(a)\le h_j(a)\) is bounded by
the sphere average of each majorant. Monotone convergence
applied to a common upper bound minus the majorants
gives the sphere condition.

Finally, for \(u\in C^2\) the derivative formula for the
spherical means used in Theorem~\ref{thm:pde-mean-weyl}
shows that \(\Delta u\ge0\) implies the submean inequality.
If \(\Delta u<0\) at a point, continuity gives strictly
decreasing means for sufficiently small radii, a contradiction.
\end{proof}

The extended-real formulation is useful for logarithmic
potentials. It also makes explicit that the strong principle
concerns a maximum on the component. See
\cite[Chapter 3]{AG} and \cite[Section 2.8]{GT}.

\section{The Poisson kernel and Harnack's inequalities}
\label{pde:L05}

\begin{theorem}[The ball Poisson problem]
\label{thm:pde-ball-poisson}
Let \(n\ge2\), \(B=B(a,r)\), \(r>0\), and \(g\in C(\partial B)\).
The kernel
\[
 P_B(x,\zeta)=\frac{r^2-|x-a|^2}
                  {r\omega_{n-1}|x-\zeta|^n}
 \qquad (x\in B,\ \zeta\in\partial B)
\]
is positive, harmonic in \(x\) and has surface integral \(1\)
for every \(x\). The Poisson integral
\(Pg(x)=\int_{\partial B}P_B(x,\zeta)g(\zeta)\,dS_\zeta\)
is the unique harmonic function in \(B\) continuous on
\(\overline B\) with boundary value \(g\).
Every harmonic function on a neighborhood of \(\overline B\)
is given there by this formula using its trace.
\end{theorem}
\begin{proof}
Translate to \(a=0\). Write \(d=|x-\zeta|\).
The Laplacian of \((r^2-|x|^2)d^{-n}\) is
\[
 2n d^{-n-2}
 \bigl(r^2-|x|^2+2x\cdot(x-\zeta)-|x-\zeta|^2\bigr)=0
\]
since \(|\zeta|=r\). The factor is positive in the ball.
Differentiation under the surface integral is valid on every
compact subball. The integral of the kernel is therefore
a smooth rotationally invariant harmonic function. The
sphere mean identity makes it constant, and evaluation at
zero gives \(1\). For a boundary point \(\zeta_0\), the
mass on the complement of any fixed cap around \(\zeta_0\)
tends to zero as \(x\to\zeta_0\): the denominator there
stays bounded below and \(r^2-|x|^2\to0\). Split the integral
into that complement and the cap. On a small cap
\(|g(\zeta)-g(\zeta_0)|\) is uniformly small, and the
kernel has unit mass. This gives the boundary limit along
every interior approach. For uniqueness a real harmonic
difference with zero boundary values cannot have a positive
maximum inside the ball: the sphere mean identity makes its
maximizer set open, hence the difference would be constant.
Apply this to both signs and to both real and imaginary parts
as needed. The final representation follows by uniqueness.
\end{proof}

\begin{corollary}[Harnack bounds and monotone convergence]
\label{cor:pde-harmonic-harnack}
Let \(u\ge0\) be harmonic on an open connected \(\Omega\).
On a closed ball \(\overline B(a,r)\subset\Omega\), for
\(|x-a|\le \rho r\), \(0\le\rho<1\), one has
\[
 \frac{1-\rho}{(1+\rho)^{n-1}}\,u(a)
 \le u(x)\le
 \frac{1+\rho}{(1-\rho)^{n-1}}\,u(a).
\]
On each compact \(K\subset\Omega\) its values are bounded above
and below by positive constant multiples of \(u(a_0)\) for
any fixed \(a_0\in\Omega\). An increasing sequence of real
harmonic functions on \(\Omega\) either tends to \(+\infty\)
everywhere or converges uniformly on compact sets to a harmonic
function. The latter alternative holds if the sequence is
bounded above at one point.
\end{corollary}
\begin{proof}
Put \(s=|x-a|/r\). The ratio
\(P_B(x,\zeta)/P_B(a,\zeta)\) lies between
\((1-s^2)/(1+s)^n\) and \((1-s^2)/(1-s)^n\) by the
triangle inequality. The lower expression is decreasing
and the upper one increasing in \(s\). Integrating against
the nonnegative boundary values gives the bounds.
Any point of \(\Omega\) can be joined to \(a_0\) by
a polygonal path in \(\Omega\). A finite chain of
overlapping strictly interior balls along the path compares
the value at \(a_0\) to values on a neighborhood of the
point. Finitely many such neighborhoods cover \(K\);
multiplying the constants along each chain gives the
compact-set assertion.

For an increasing sequence \(u_j\) the differences
\(u_m-u_j\), \(m\ge j\), are nonnegative harmonic.
If \(u_j(a_0)\) is bounded above, its scalar differences tend
to zero, and the compact Harnack bound makes the sequence
uniformly Cauchy on every compact set. The limit is continuous
and retains all local mean values by uniform convergence on
each closed ball. It is harmonic by
Theorem~\ref{thm:pde-mean-weyl}. If instead the values
at \(a_0\) diverge, apply the lower bound to \(u_j-u_1\)
to obtain divergence at every point.
\end{proof}

The statements in one dimension follow directly for affine
harmonic functions on intervals when needed.
See \cite[Sections 1.3--1.5]{AG} and
\cite[Sections 2.3 and 2.5]{GT}. The disk specialization
is also treated in \cite[Chapter 13]{D12}.

\section{The ball Green kernel and the Poisson--Dirichlet problem}
\label{pde:L06}

\begin{theorem}[Poisson--Dirichlet on a ball]
\label{thm:pde-poisson-dirichlet-ball}
For \(n\ge2\), a ball \(B\), \(f\in C_c^\infty(B)\), and
\(g\in C(\partial B)\), there is a unique
\(u\in C^2(B)\cap C(\overline B)\) with
\(-\Delta u=f\) and boundary data \(g\). It has representation
\[
 u(x)=\int_B G_B(x,y)f(y)\,dy+
                   \int_{\partial B}P_B(x,\zeta)g(\zeta)\,dS_\zeta ,
\]
where \(G_B\) is the Dirichlet Green kernel. In the second
variable \(-\partial_\nu G_B(x,\zeta)=P_B(x,\zeta)\).
\end{theorem}
\begin{proof}
Translate to the radius-\(r\) ball about zero. For \(x\ne0\) set
\(x^*=r^2x/|x|^2\). Define
\[
 G_B(x,y)=
 \begin{cases}
 \displaystyle\Phi(y-x)-(r/|x|)^{n-2}\Phi(y-x^*),
       & n\ge3,\\[1mm]
 \displaystyle\frac{1}{2\pi}
       \log\frac{|x|\,|y-x^*|}{r\,|y-x|},& n=2 .
 \end{cases}
\]
For \(x=0\) use \(\Phi(y)-\Phi(r)\), with \(\Phi(r)\)
denoting the radial value. The image pole \(x^*\) lies
outside the ball, so the correction is harmonic in \(y\)
and \(-\Delta_yG_B(x,\cdot)=\delta_x\).
The equality
\(|\zeta-x^*|=(r/|x|)|\zeta-x|\) on the boundary gives
zero boundary value in both dimensions. Symmetry follows from
\[
 |x|^2|y-x^*|^2
    =r^4-2r^2 x\cdot y+|x|^2|y|^2 .
\]
The resulting formula in terms of this symmetric polynomial
extends to \(x=0\) or \(y=0\) off the diagonal.
For the derivative, use
\(\nabla\Phi(v)=-v/(\omega_{n-1}|v|^n)\)
for all \(n\ge2\). On the boundary the correction's
normal derivative has the factor
\[
 (r/|x|)^{n-2}|\zeta-x^*|^{-n}
      =(|x|^2/r^2)|\zeta-x|^{-n},
\]
with exponent \(n-2=0\) in dimension two. In the
difference of the two fluxes the numerators
\(r^2-\zeta\cdot x\) and \(|x|^2-\zeta\cdot x\)
give \(r^2-|x|^2\). The sign of the outward derivative
is negative, yielding the asserted Poisson kernel.
At \(x=0\) it follows directly from the radial derivative.

For \(y\) in the compact support of \(f\), the harmonic
correction as a function of \(x\) is smooth up to the
sphere with uniformly bounded derivatives there. Indeed
the symmetric polynomial above is at least
\((r^2-r|y|)^2>0\) on the closed ball, with uniform positive
lower bound for such \(y\). Thus the volume integral
is a Newton potential minus a smooth harmonic correction.
It solves \(-\Delta u=f\) and has zero trace, the latter
by the kernel's boundary identity and continuity near the
sphere away from the source support. Adding the harmonic
Poisson integral proves existence. Uniqueness is the
maximum principle on the difference.
\end{proof}

Compare \cite[Section 2.2.4]{Evans} and
\cite[Sections 2.4--2.5 and 4.2]{GT}.

\section{Regular domains, harmonic lifting, and Green representation}
\label{pde:L07}

\begin{lemma}[Continuous harmonic Dirichlet data on a regular domain]
\label{lem:pde-harmonic-lift}
Let \(\Omega\) be a bounded \(C^2\) open set, \(n\ge2\).
Every \(g\in C(\partial\Omega)\) has a unique harmonic extension
in \(C(\overline\Omega)\). For boundary data which are the trace
of a smooth function \(G_0\) on \(\overline\Omega\), the extension
agrees with the weak harmonic solution in
\(G_0+H^1_0(\Omega)\) and lies in \(H^2(\Omega)\).
\end{lemma}
\begin{proof}
We first treat real data. At each \(z\in\partial\Omega\)
there is an exterior tangent closed ball \(\overline B(y,r)\)
meeting \(\overline\Omega\) only at \(z\). Locally, in graph
coordinates, the \(C^2\) bound makes the graph lie on the
interior side of a sufficiently small such sphere by the
quadratic Taylor bound. Choosing its radius small makes the
whole ball lie in that coordinate neighborhood, so the contact
condition holds globally. The harmonic function
\[
 b_z(x)=
 \begin{cases}
  r^{2-n}-|x-y|^{2-n},&n\ge3,\\
  \log(|x-y|/r),&n=2
 \end{cases}
\]
is continuous on \(\overline\Omega\), nonnegative there,
and vanishes precisely at \(z\). For each \(\epsilon>0\),
the harmonic functions
\(g(z)-\epsilon-M b_z\) and \(g(z)+\epsilon+M b_z\)
bound \(g\) on the boundary when \(M\) is large: close to
\(z\) use continuity of \(g\), and on the compact rest of
the boundary use the positive minimum of \(b_z\).

For the existence argument take the supremum \(U\) of
continuous real subharmonic functions on \(\overline\Omega\)
whose boundary values are at most \(g\), with subharmonicity
required only in \(\Omega\). The class contains a constant
minorant, and all its members are bounded above by
\(\max_{\partial\Omega}g\). It is closed under finite
maxima, by harmonic comparison. It is also closed under
harmonic replacement on any ball \(D\Subset\Omega\):
use the Poisson solution with the function's boundary data
on that ball. The replacement majorizes the old function
and joins it continuously. For comparison on a test ball
crossing \(\partial D\), compare the old function first
throughout the test ball, then compare the replacement's
harmonic part on its intersection with \(D\).
The boundary inequalities on the two parts of that
intersection follow respectively from the old function and
the test boundary. The bounded-domain maximum principle
gives comparison there as well, proving subharmonicity.

Fix \(a\) in a ball \(D\Subset\Omega\). Take a sequence in
the class whose values at \(a\) approach \(U(a)\).
Successively take maxima with the previous member and then
replace harmonically on \(D\). This gives an increasing
sequence of class members harmonic on \(D\).
Its limit \(h\) is harmonic on \(D\) by Harnack convergence,
with \(h(a)=U(a)\) and \(h\le U\). For any other class
member \(v\), perform harmonic replacement on the maxima
of \(v\) with this sequence. The new limit \(k\) is harmonic,
satisfies \(k\ge h,v\) on \(D\), and \(k(a)=h(a)\).
The nonnegative harmonic difference \(k-h\) vanishes at an
interior point, so it is zero by the strong principle.
This gives \(h\ge v\) for every member and therefore
\(U=h\) on \(D\). Thus \(U\) is harmonic in \(\Omega\).
The barrier minorants belong to the class, and the barrier
majorants bound every member by comparison. Letting
\(x\to z\) and then \(\epsilon\to0\) proves continuity with
boundary value \(g\). Uniqueness follows by comparison.
Real and imaginary parts give the complex statement.

For the smooth-lift assertion, solve weakly for
\(w\in H^1_0(\Omega)\) with
\(\int\nabla w\cdot\nabla v=-\int\nabla G_0\cdot\nabla v\).
The coercive Dirichlet theorem and the \(C^{1,1}\) boundary
estimate of Sections~\ref{pde:E01} and \ref{pde:E06}
apply, since \(C^2\) graph charts on a compact boundary
are \(C^{1,1}\) on smaller patches and \(\Delta G_0\in L^2\).
Thus \(\widetilde h=G_0+w\) belongs to \(H^2\) and is
harmonic in the interior by Weyl's lemma.
The same barrier bounds hold for \(\widetilde h\) by weak
comparison in Section~\ref{pde:E04}: the positive part of
the difference with a majorant has zero trace and is an
admissible test, and the negative part treats the minorant.
The bound is pointwise in the interior by harmonic
continuity. Hence \(\widetilde h\) has the continuous
boundary data, and uniqueness identifies the two lifts.
\end{proof}

\begin{theorem}[Green representation on bounded \(C^2\) open sets]
\label{thm:pde-green-representation}
For \(x\in\Omega\) let \(G(x,\cdot)\) be the Dirichlet Green
function of \(-\Delta\), normalized by
\(-\Delta_y G(x,y)=\delta_x(y)\). Such a function exists and is
unique with continuous zero boundary values away from the pole.
Off the diagonal it is symmetric and is positive for two
points in the same component, with value zero across
different components. For \(u\in C^2(\overline\Omega)\),
\[
 u(x)=\int_\Omega G(x,y)(-\Delta u(y))\,dy
       -\int_{\partial\Omega}u(y)\,\partial_{\nu_y}G(x,y)\,dS_y.
\]
The outward derivative can be interpreted as the weak normal
trace on a collar of the boundary.
\end{theorem}
\begin{proof}
The difference of two such Green functions is distributionally
harmonic across the pole. By Weyl's lemma it extends harmonically
through the pole, and the maximum principle with zero boundary
values makes it zero.
Assume first \(n\ge2\).
Fix \(x\) and take the harmonic lift \(h_x\) of
\(\Phi(x-\cdot)\) from the boundary. A smooth lift
of this data is obtained by multiplying the fundamental
solution by a smooth cutoff zero near \(x\) and one
near \(\partial\Omega\). The preceding lemma gives
\(h_x\in H^2(\Omega)\). Set
\(G(x,y)=\Phi(y-x)-h_x(y)\).
This gives its distributional equation, harmonicity away
from the pole, and zero continuous boundary value.
On a boundary collar \(G\) has \(H^2\) regularity. Hence
its normal trace agrees with the \(H^1\) trace of its
gradient paired with the outward normal, by smooth
approximation in the Green identity of
Theorem~\ref{thm:pde-normal-trace}. It is thus possible
to read the displayed integral either as a boundary
integral of traces or as a trace duality.

The fundamental solution diverges positively as \(y\to x\).
Comparison on the punctured region and then shrinking
the puncture gives nonnegativity. On the component of
\(x\), deleting that one point leaves a connected open
set for \(n\ge2\): polygonal paths can be detoured in a
small ball around the point. The strong principle and
the positive values near the pole give strict positivity
there. On the other components the function is harmonic
with zero continuous boundary and hence zero.

Apply the weak second Green identity on
\(\Omega\setminus\overline B(x,\epsilon)\), with
\(\epsilon\) small enough that the closed ball lies in
\(\Omega\). The functions are \(H^2\) there, so the
identity applies; the derivative on the inner sphere is
classical. Its volume side is
\(\int G(x,y)(-\Delta u(y))\,dy\) on the punctured region.
The outer contribution is
\(\int_{\partial\Omega}u\,\partial_\nu G\) since \(G=0\)
there. On the inner sphere, the outward normal of the
punctured region points towards \(x\), so
\(\partial_\nu\Phi=1/(\omega_{n-1}\epsilon^{n-1})\).
The term \(u\,\partial_\nu G\) tends in integral to \(u(x)\);
the derivative of the smooth harmonic correction contributes
\(O(\epsilon^{n-1})\). The term \(G\,\partial_\nu u\)
has integral tending to zero, of order \(O(\epsilon)\) or
\(O(\epsilon|\log\epsilon|)\). Local integrability at the
pole gives the asserted formula. For symmetry apply the same
second identity to \(G(x,\cdot)\) and \(G(z,\cdot)\) with
two small disjoint balls removed. Both are harmonic in
the punctured domain and zero on the outer boundary.
On the inner sphere at \(x\) the limiting contribution is
\(-G(z,x)\), and on the one at \(z\) it is \(G(x,z)\),
by the flux computation just made. Their sum is zero.

In dimension one the component of \(x\) is an interval
\((a,b)\). There set
\[
 G(x,y)=\frac{(\min(x,y)-a)(b-\max(x,y))}{b-a},
\]
and set the kernel equal to zero on other intervals.
The function is symmetric, continuous, and piecewise affine
in \(y\), with slope jump \(-1\) at \(x\) and zero endpoint
values. It is strictly positive in the open interval.
One integration by parts on each subinterval gives
the representation, with \(\partial_\nu=-\partial_y\)
at \(a\) and \(\partial_\nu=\partial_y\) at \(b\).
This also verifies the normalization.
\end{proof}

For comparison see \cite[Sections 4.1, 6.6 and 6.8]{AG},
\cite[Sections 2.8 and 4.2]{GT}, and the planar results
of \cite[Chapter 14]{D12}.
\chapter{Linear Elliptic Equations}\label{pde:elliptic-chapter}

Except in the spectral sections, the functions and coefficients in this
chapter are real. For a matrix field \(A\), ellipticity with constant
\(\lambda>0\) means
\[
 \xi^{\mathsf T}A(x)\xi\geq \lambda|\xi|^2
 \qquad(\xi\in\mathbb R^n)
\]
at almost every \(x\) for measurable coefficients, and at every \(x\)
for continuous coefficients. In particular the antisymmetric part does
not enter this inequality. The positive sign convention for a
divergence-form realization is \(-\operatorname{div}(A\nabla)\). The
classical maximum principle below uses a nondivergence operator with
positive second-order coefficients and the opposite sign convention.

\section{Coercive weak Dirichlet problems}\label{pde:E01}

Write \(H^{-1}(\Omega)=(H_0^1(\Omega))'\), with duality denoted by
\(\langle f,v\rangle\). The norm on \(H_0^1\) is the full \(H^1\) norm
unless otherwise specified.

\begin{theorem}[Coercive Dirichlet problem]\label{pde:elliptic-weak-dirichlet}
Let \(\Omega\subset\mathbb R^n\) be bounded and open, \(n\geq1\).
Let \(A\in L^\infty(\Omega;\mathbb R^{n\times n})\) be elliptic
with constant \(\lambda>0\). For \(f\in H^{-1}(\Omega)\) there is
exactly one \(u\in H_0^1(\Omega)\) satisfying
\[
 a(u,v):=\int_\Omega A\nabla u\cdot\nabla v
       =\langle f,v\rangle \qquad(v\in H_0^1(\Omega)).
\]
It satisfies \(\|u\|_{H^1}\leq C_\Omega\lambda^{-1}
\|f\|_{H^{-1}}\). If \(\Omega\) is Lipschitz and
\(g\in H^{1/2}(\partial\Omega)\), there is exactly one \(H^1\)
solution of the same equation with trace \(g\), and its \(H^1\)
norm is bounded by
\[
 C_\Omega\bigl(\lambda^{-1}\|f\|_{H^{-1}}+
 (1+\lambda^{-1}\|A\|_\infty)\|g\|_{H^{1/2}(\partial\Omega)}\bigr).
\]
\end{theorem}

\begin{proof}
Take a slab \(a<x_1<b\) containing \(\Omega\). Extend a compactly
supported smooth \(v\) by zero to the slab. On each coordinate line
\(v(x_1,x')=\int_a^{x_1}\partial_1v(s,x')\,ds\).
Cauchy--Schwarz followed by integration gives
\(\|v\|_2\leq(b-a)\|\partial_1v\|_2\), hence the gradient norm is
equivalent to the full norm on the closure of the tests.
Completeness is inherited from the closed subspace \(H_0^1\) of
\(H^1\). In the gradient Hilbert norm the form has bound
\(\|A\|_\infty\) and coercivity constant \(\lambda\).
The nonsymmetric Lax--Milgram theorem gives the solution. Testing
with \(u\) proves the estimate after applying the norm equivalence.

On a Lipschitz domain the trace map has a bounded right inverse and
kernel \(H_0^1\). If \(w\) lifts \(g\), solve
\(a(z,v)=\langle f,v\rangle-a(w,v)\) for \(z\in H_0^1\) and put
\(u=w+z\). The last bound follows from the bound for the homogeneous
problem and that for the lift. Two choices of lift give solutions
whose difference has zero trace and solves the homogeneous equation,
so that difference vanishes.
\end{proof}

The trace and lifting theorems are the results in Unit U5 of
\cite{P102}. The Hilbert theorem and the zero-boundary compactness
framework are in V03.07 and V08.01 of \cite{P106}. Compare
\cite[Sections 6.2.1--6.2.2]{Evans},
\cite[Sections 9.2.2--9.2.4]{RR}, and \cite[Section 8.2]{GT}.
If \(A\) is symmetric, fix a lift \(w\) of the boundary data, with
\(w=0\) in the homogeneous case. The functional on the affine trace
class is
\[
 J(w+z)=\tfrac12a(w+z,w+z)-\langle f,z\rangle,
                  \qquad z\in H_0^1 .
\]
This pairing is defined even when \(f\) acts only on \(H_0^1\).
Changing the lift adds a constant to \(J\). For any variation
\(v\in H_0^1\), its increment at the solution is
\(\tfrac12a(v,v)\) by the variational equation, so the solution is
its unique minimizer. For nonsymmetric \(A\) the quadratic
derivative uses the symmetric part of \(A\), so we use the
variational equation directly.

\section{Weak Neumann data and constants}\label{pde:E02}

\begin{theorem}[Neumann compatibility]\label{pde:elliptic-neumann}
Let \(\Omega\) be a nonempty bounded connected Lipschitz open subset
of \(\mathbb R^n\), let \(f\in L^2(\Omega)\), and let
\(g\in H^{-1/2}(\partial\Omega)\). A weak solution of
\(-\Delta u=f,\ \partial_\nu u=g\) means \(u\in H^1(\Omega)\) with
\[
 \int_\Omega\nabla u\cdot\nabla v
  =\int_\Omega fv+\langle g,\gamma v\rangle
   \qquad(v\in H^1(\Omega)),
\]
where \(\gamma\) is the trace. A solution exists exactly when
\(\int_\Omega f+\langle g,1\rangle=0\).
It is unique modulo constants and its representative of mean zero
satisfies
\[
 \|u\|_{H^1}\leq C_\Omega
 \bigl(\|f\|_{2}+\|g\|_{H^{-1/2}(\partial\Omega)}\bigr).
\]
For the quotient norm on \(H^1/\mathbb R\) the same estimate holds.
\end{theorem}

\begin{proof}
The constant test \(v=1\) proves necessity. On the closed mean-zero
subspace of \(H^1\), the gradient norm is equivalent to the full
norm. Indeed, otherwise there would be a sequence of mean zero,
unit \(L^2\) norm, and gradients tending to zero. Rellich compactness
gives a strongly \(L^2\) convergent subsequence. Its limit has mean
zero and norm one, and its distributional gradient is zero.
A function with zero distributional gradient is constant on a
connected open set: it is constant on each ball after local
mollification, and overlapping balls identify the constants.
This contradicts the two conditions on the limit.
The trace bound and Cauchy--Schwarz make the displayed right-hand
side a bounded functional on the mean-zero Hilbert space.
Lax--Milgram supplies a unique solution in that space and the
estimate. Subtracting the mean from any \(v\in H^1\) changes
neither side by compatibility, giving the full test class. The
difference of any two solutions has gradient zero by testing with
itself and is constant.

Compactly supported tests give \(\operatorname{div}\nabla u=-f\).
The weak normal trace in Section~\ref{pde:C02}, applied to
\(\nabla u\), is the functional on traces given by
\(\int_\Omega\nabla u\cdot\nabla v-\int_\Omega fv\).
Surjectivity of the trace then identifies that functional with \(g\).
\end{proof}

Connectedness is used to identify the zero-gradient space.
On a finite disjoint union with componentwise Lipschitz closures
and separate traces the same argument applies on each component,
with one compatibility condition and one free constant per
component. See \cite[Section 9.4.1]{RR},
\cite[Chapter 5, Section 7]{TaylorI}, and the mean-zero Poincaré
and trace theorems in Unit U5 of \cite{P102}.

\section{The classical maximum and Hopf principles}\label{pde:E03}

\begin{theorem}[Classical comparison]\label{pde:elliptic-max}
Let \(\Omega\) be a bounded connected open set and
\[
 Lu=a^{ij}u_{ij}+b^iu_i+cu ,
\]
where repeated indices are summed, \(a^{ij}=a^{ji}\), the
coefficients are continuous in \(\Omega\), the principal matrix
is uniformly positive definite, and \(c\leq0\).
If \(u\in C^2(\Omega)\cap C(\bar\Omega)\) and \(Lu\geq0\), then
\[
 \sup_\Omega u\leq\max\{0,\sup_{\partial\Omega}u\}.
\]
If \(u\) attains a global maximum over \(\Omega\) at an interior
point and the maximum is nonnegative, \(u\) is constant.
For \(c=0\) the boundary bound is \(\sup_{\partial\Omega}u\)
and the strong conclusion has no sign restriction.
Consequently \(Lu\geq Lv\) and \(u\leq v\) at the boundary
imply \(u\leq v\) in the domain.
\end{theorem}

\begin{proof}
First suppose the coefficients are bounded. Choose \(k>0\) so large
that \(\lambda k^2-\|b\|_\infty k-\|c\|_\infty>0\).
Then \(L(e^{kx_1})>0\). For every \(\epsilon>0\) the perturbed
function \(u+\epsilon e^{kx_1}\) has strictly positive \(L\), so it
cannot have a positive interior maximum. At such a maximum the
gradient is zero, the Hessian is negative semidefinite, and its
contraction with the positive principal matrix, plus the zero-order
term, is nonpositive. Applying this to the perturbed function and
letting \(\epsilon\downarrow0\) gives the weak bound. For \(c=0\)
no interior maximum of the strict perturbation is possible at any
height, which proves the sharper bound.

Here is a local barrier argument for the strong conclusion. Suppose
\(M=\max_\Omega u\) is attained and put \(w=M-u\geq0\).
Then \(Lw=cM-Lu\leq0\), also for an arbitrary \(M\) if \(c=0\).
If both the zero set of \(w\) and the set \(w>0\) are nonempty,
there is a ball \(B_R(y)\) on which \(w>0\), whose closure lies in
\(\Omega\) and whose sphere contains a zero \(z\) of \(w\).
To construct it, take a boundary point of \(\{w>0\}\) in
\(\Omega\), choose \(y\) in that set sufficiently close to the
point, and take its distance to the zero set as radius. The
distance is attained on a compact ball inside \(\Omega\).
On the annulus \(R/2<|x-y|<R\), set
\(\psi=e^{-k|x-y|^2}-e^{-kR^2}\). Its \(L\) is positive for \(k\)
large. Indeed, after division by \(e^{-k|x-y|^2}\), it is bounded
below by
\[
 4k^2\lambda(R/2)^2-2k\|\operatorname{tr}a\|_\infty
 -2kR\|b\|_\infty-\|c\|_\infty
\]
on this annulus. Take \(\epsilon>0\) with \(w\geq\epsilon\psi\)
on the inner sphere, using compactness and \(w>0\) there.
This comparison also holds on the outer sphere since \(\psi=0\).
The bounded-coefficient weak principle applied to
\(\epsilon\psi-w\) proves it in the annulus. At \(z\) the estimate
gives a strictly positive one-sided inward slope for \(w\),
whereas its gradient at its interior minimum vanishes.
This contradiction proves the strong principle using only local
coefficient bounds.

Finally, when the coefficients are merely continuous in the open
domain, any violation of the desired weak bound yields a global
maximum on the compact closure attained in the interior, with the
sign required above. Constancy contradicts the boundary values.
The last assertion follows by applying the bound to \(u-v\).
\end{proof}

\begin{corollary}[Hopf boundary inequality]\label{pde:elliptic-hopf}
Suppose \(Lu\geq0\), \(u\in C^2(\Omega)\cap C(\bar\Omega)\),
and \(z\in\partial\Omega\) has an interior tangent ball. Assume the
coefficients are bounded near \(z\) in that ball, with uniform
ellipticity, and \(c\leq0\). If \(u(z)>u(x)\) in the ball and
\(u(z)\geq0\), then
\[
 \liminf_{h\downarrow0}\frac{u(z)-u(z-h\nu)}{h}>0,
\]
where \(\nu\) is the outward radial direction of the tangent ball
at \(z\). For \(c=0\) the sign restriction on \(u(z)\) is absent.
In particular the outward directional derivative is positive if
it exists. At a differentiable boundary point with a tangent ball
this direction is the outward unit normal.
\end{corollary}

\begin{proof}
Use \(w=u(z)-u\), positive in the ball, and the same annular
barrier. One may first reduce to a smaller interior ball tangent
at \(z\), to have bounded coefficients on the annulus. The weak
principle uses continuity on its closure. On the inward normal
segment \(\psi(z-h\nu)/h\to2kR e^{-kR^2}>0\).
\end{proof}

Compare \cite[Sections 3.1--3.2]{GT} and
\cite[Sections 4.1.1--4.1.2]{RR}. The usual bounded connected
\(C^1\) domain version is a special case of the maximum principle.
The geometric hypothesis in the Hopf assertion is an additional
interior sphere condition.

\section{Weak comparison and a local energy estimate}\label{pde:E04}

For \(v\in H^1\) the function \(v^+=\max\{v,0\}\) is in \(H^1\),
with \(\nabla v^+=1_{\{v>0\}}\nabla v\).
Indeed, approximate the positive-part function by smooth
functions with derivatives in \([0,1]\), equal to zero on the
negative half-line, whose derivatives tend to \(1_{\{s>0\}}\).
The chain rule for a smooth scalar function with bounded derivative
follows by approximation of \(v\) by smooth functions, using the
bounded derivative and dominated convergence on the limiting
gradient. Passing to the limit proves the formula. Applying it
also to \(-v\) shows \(\nabla v=0\) almost everywhere on \(\{v=0\}\).
These formulas give continuity of positive part in \(H^1\):
for a strongly convergent sequence, take an almost everywhere
convergent subsequence, use the gradient formula on the limit
gradient, and use that gradient's vanishing on the zero set.
The subsequence criterion gives continuity for the sequence.
On a Lipschitz domain approximation and trace continuity give
\(\gamma(v^+)=(\gamma v)^+\), since positive part is continuous
in \(L^2\) of the boundary.

\begin{theorem}[Weak divergence-form comparison]\label{pde:elliptic-weak-max}
Let \(\Omega\) be bounded and open, with bounded elliptic \(A\)
and \(0\leq c\in L^\infty\). Suppose \(u\in H^1(\Omega)\),
\(u^+\in H_0^1(\Omega)\), and
\[
 \int_\Omega A\nabla u\cdot\nabla\phi+cu\phi\leq0
 \qquad(0\leq\phi\in C_c^\infty(\Omega)).
\]
Then \(u\leq0\) almost everywhere. In particular weak sub- and
supersolutions for \(-\operatorname{div}(A\nabla)+c\) compare
if their difference has positive part in \(H_0^1\), a condition
guaranteed by ordered traces on a Lipschitz boundary.
\end{theorem}

\begin{proof}
Nonnegative smooth compactly supported tests are dense in the
nonnegative cone of \(H_0^1\). To see this, take compactly
supported smooth approximants, pass to their positive parts,
and mollify each positive part within its distance from the
boundary. Thus \(u^+\) is an admissible test. The chain rule gives
\(\lambda\|\nabla u^+\|_2^2+\int c(u^+)^2\leq0\).
The slab inequality from Section~\ref{pde:E01} gives \(u^+=0\).
Apply the same reasoning to the difference for comparison.
\end{proof}

\begin{proposition}[Caccioppoli inequality]\label{pde:elliptic-caccioppoli}
If \(u\in H^1_{\mathrm{loc}}(\Omega)\),
\(-\operatorname{div}(A\nabla u)=f\) weakly,
\(A\) is bounded and elliptic on a neighborhood of the support
of \(\eta\in C_c^\infty(\Omega)\), and \(f\in L^2_{\mathrm{loc}}\),
then
\[
 \int_\Omega \eta^2|\nabla u|^2
 \leq C\int_\Omega|\nabla\eta|^2 |u|^2
     +\frac{2}{\lambda}\int_\Omega |\eta^2 f u|.
\]
Here \(C\) depends only on the local bound for \(A\) and \(\lambda\).
\end{proposition}

\begin{proof}
The function \(\eta^2u\) is an admissible test by compact support
and \(H^1\) approximation. Expand its gradient. The term
\(2\int \eta u A\nabla u\cdot\nabla\eta\) is bounded in absolute
value by \(\frac{\lambda}{2}\int \eta^2|\nabla u|^2\)
plus \(\frac{2\|A\|_\infty^2}{\lambda}
\int|\nabla\eta|^2|u|^2\). Move the first term to the left and
bound the source pairing by its absolute integral.
\end{proof}

See \cite[Sections 8.1 and 8.6]{GT}. This argument needs the
Sobolev truncation and no pointwise derivative at an extremum.

\section{Interior second derivatives in \texorpdfstring{\(L^2\)}{L2}}\label{pde:E05}

\begin{theorem}[Interior \(H^2\) estimate]\label{pde:elliptic-interior-h2}
Let \(A\) be locally Lipschitz and locally uniformly elliptic in
\(\Omega\), and let \(u\in H^1_{\mathrm{loc}}(\Omega)\) solve
\(-\operatorname{div}(A\nabla u)=f\) with
\(f\in L^2_{\mathrm{loc}}(\Omega)\). Then \(u\in
H^2_{\mathrm{loc}}(\Omega)\). If \(\Omega'\Subset\Omega''
\Subset\Omega\) are open, then
\[
 \|u\|_{H^2(\Omega')}\leq C
   \bigl(\|u\|_{H^1(\Omega'')}+\|f\|_{L^2(\Omega'')}\bigr).
\]
The constant depends on the nested sets and the corresponding
coefficient and ellipticity bounds. Symmetry is unnecessary.
\end{theorem}

\begin{proof}
Choose \(0\leq\eta\leq1\) smooth, equal to one on a neighborhood
of \(\bar\Omega'\) and supported in \(\Omega''\). All translates
below are taken small enough to stay in \(\Omega''\) on the
supports of the tests. For one coordinate \(e_k\) write
\[
 D_h v(x)=\frac{v(x+he_k)-v(x)}{h}.
\]
Use \(-D_{-h}(\eta^2 D_hu)\) as a weak test. It is a compactly
supported \(H^1\) function for each nonzero \(h\). Discrete
integration by parts from Section~\ref{pde:C04} gives
\[
 \int D_h(A\nabla u)\cdot\nabla(\eta^2D_hu)
       =-\int f D_{-h}(\eta^2D_hu).
\]
Put \(X_h=\|\eta D_h\nabla u\|_2\) and
\(Y=\|u\|_{H^1(\Omega'')}\). In the product difference,
\[
 D_h(A\nabla u)=A(x+he_k)D_h\nabla u+(D_h A)\nabla u.
\]
The first part against \(\eta^2D_h\nabla u\) gives at least
\(\lambda X_h^2\). The remaining three products after expanding
\(\nabla(\eta^2D_hu)\) are bounded by \(C(X_hY+Y^2)\).
Here one uses the Lipschitz bound for \(D_h A\) and the
difference bound for \(D_hu\) in terms of \(\nabla u\).
For the right-hand side the same difference bound, now on the
extended compactly supported function \(\eta^2D_hu\), gives
\[
 \bigl|\textstyle\int f D_{-h}(\eta^2D_hu)\bigr|
       \leq C\|f\|_2(X_h+Y).
\]
Absorbing terms linear in \(X_h\) yields
\(X_h\leq C(Y+\|f\|_2)\) uniformly for small \(h\).
Thus the coordinate difference quotients of each component of
\(\nabla u\) are bounded in \(L^2(\Omega')\). The local inverse
difference-quotient criterion identifies their weak limits as
the corresponding distributional derivatives, with the same
bounds. Including the original \(H^1\) norm gives the estimate.
\end{proof}

This is the interior argument in \cite[Section 6.3.1]{Evans}.
Compare \cite[Section 9.5.1]{RR} and \cite[Section 8.3]{GT}.
The estimate uses no derivative of \(f\).

\section{The \texorpdfstring{\(C^{1,1}\)}{C(1,1)} Dirichlet boundary}\label{pde:E06}

\begin{theorem}[Global Dirichlet \(H^2\) estimate]\label{pde:elliptic-global-h2}
Let \(\Omega\) be bounded with \(C^{1,1}\) boundary, let \(A\)
be real symmetric, bounded, Lipschitz, and uniformly elliptic
on \(\Omega\). Suppose \(f\in L^2(\Omega)\) and \(u\in H_0^1(\Omega)\)
satisfies \(-\operatorname{div}(A\nabla u)=f\) weakly. Then
\[
 u\in H^2(\Omega),\qquad
 \|u\|_{H^2(\Omega)}\leq C
    \bigl(\|f\|_{L^2(\Omega)}+\|u\|_{H^1(\Omega)}\bigr).
\]
The constant depends on the domain and the coefficient bounds,
including the ellipticity constant.
\end{theorem}

\begin{proof}
Near each boundary point rotate coordinates and write the domain as
\(x_n>\varphi(x')\) with \(\varphi\in C^{1,1}\).
On a small half-cylinder use \(x=F(y)=(y',y_n+\varphi(y'))\).
Both \(F,F^{-1}\) are \(C^{1,1}\), the Jacobian determinant is
one, and both differential matrices are bounded. The weak
equation for \(w=u\circ F\) has matrix
\[
 B(y)=DF(y)^{-1}A(F(y))\,DF(y)^{-\mathsf T}
\]
and source \(f\circ F\). This follows by the \(H^1\) chain rule
and change of variables, first for smooth approximations and then
in \(H^1\). The transformed matrix is Lipschitz, bounded and
elliptic with controlled constants. The bottom trace of \(w\)
is zero, by the trace transformation or approximation from
\(C_c^\infty(\Omega)\).

Take a cutoff in the half-cylinder, allowed to reach the flat
bottom and supported away from all the other faces.
For each tangential direction the test
\(-D_{-h}(\eta^2 D_hw)\) has zero bottom trace and is supported
away from the other faces. Its pushforward, extended by zero,
belongs to \(H_0^1(\Omega)\). The difference integration by parts
holds over the half-cylinder because the shifts are tangential
and all side translates stay inside. The estimate of
Section~\ref{pde:E05} therefore applies verbatim in these directions,
bounding \(D_iD_jw\) on a smaller half-cylinder whenever at least
one of \(i,j\) is tangential. Mixed distributional derivatives
commute. The weak equation now gives
\[
 \partial_n(B_{nn}\partial_n w)
   =-f\circ F-\sum_{(i,j)\ne(n,n)}\partial_i(B_{ij}\partial_jw)
       \quad\hbox{in }L^2 .
\]
Every term on the right is \(L^2\) by the tangential result and the
Lipschitz product rule. Since \(B_{nn}\geq\lambda_0>0\), its
reciprocal is Lipschitz. Multiplication by this reciprocal in the
one-direction Sobolev product rule gives \(\partial_{nn}w\in L^2\),
with the claimed bound. The product rule here follows by applying
the one-dimensional Sobolev rule on almost every normal line
and integrating, or by smooth approximation of the multiplier.
For \(n=1\) the sum is empty and this is the complete local
boundary argument.

Pullback and pushforward preserve \(H^2\) locally under a
\(C^{1,1}\) coordinate change: the first derivatives transform
by the \(H^1\) chain rule and the resulting coefficient matrices
have bounded weak first derivatives, giving the second-order
product rule. Thus each small boundary chart gives an \(H^2\)
bound in the original coordinates. A finite boundary cover and
an interior compact region, where Section~\ref{pde:E05} applies,
give the global estimate. All second derivatives obtained locally
are restrictions of the corresponding distributions, so they
agree on the overlaps.
\end{proof}

\begin{corollary}[A lifted boundary value]\label{pde:elliptic-h2-lift}
Under the domain and coefficient hypotheses of the global estimate,
if \(w\in H^2(\Omega)\) is prescribed as a lift and \(u-w\in H_0^1\),
then any weak solution
of \(-\operatorname{div}(A\nabla u)=f\in L^2\) is in \(H^2\).
Its estimate has the additional term \(C\|w\|_{H^2}\).
In particular the gradients have \(H^{1/2}\) boundary traces.
\end{corollary}

\begin{proof}
Apply the homogeneous estimate to \(u-w\) with source
\(f+\operatorname{div}(A\nabla w)\in L^2\) and add \(w\).
The trace assertion is the \(H^1\) trace theorem applied to
each weak first derivative.
\end{proof}

Grisvard's Theorem 2.4.2.5 gives the \(C^{1,1}\) comparison
with the zero lower coefficients, \(p=2\), and zero boundary
data \cite[pp. 124--125]{Grisvard}. The \(C^2\) boundary argument
in \cite[Section 6.3.2]{Evans} and the corresponding arguments
in \cite[Section 9.5.2]{RR} and \cite[Section 8.4]{GT} are
useful comparisons. The Lipschitz control of the transformed
coefficients is exactly what the difference estimate uses.

\section{The interior Poisson Schauder estimate}\label{pde:E07}

Write \([g]_{\alpha;U}=\sup_{x\ne y\in U}
 |g(x)-g(y)|/|x-y|^\alpha\), and use the ordinary Hölder
norms obtained by adding the relevant supremum norms.
All the Hölder exponents below satisfy \(0<\alpha<1\).

\begin{lemma}[Cancellation on Hölder data]\label{pde:elliptic-holder-kernel}
In \(\mathbb R^m\), let \(K\) be smooth off the origin,
homogeneous of degree \(-m\), with zero spherical mean.
For \(g\in C_c^{0,\alpha}(\mathbb R^m)\), the principal value
\(Tg(x)=\operatorname{pv}\int K(x-y)g(y)\,dy\) exists and is
\(C^{0,\alpha}\). If the support is in a fixed ball, its
\(C^{0,\alpha}\) norm is bounded by a constant times that of \(g\).
The constant depends on \(\alpha\), the ball, and bounds of
\(K,\nabla K\) on the unit sphere.
\end{lemma}

\begin{proof}
On \(|x-y|<1\) subtract \(g(x)\); integration in radius is then
bounded by \(C[g]_\alpha\int_0^1 r^{\alpha-1}\,dr\).
On the complement the absolute integral is at most \(C\|g\|_1\).
To compare \(x,x'\), put \(h=|x-x'|\leq1\) and split at
\(B_{2h}(x)\). In this ball, subtract \(g(x)\) for the first kernel
and \(g(x')\) for the second. The two integrals so obtained are
bounded by \(C[g]_\alpha h^\alpha\), since \(B_{2h}(x)\subset
B_{3h}(x')\). One may make these subtractions over the entire
space in the following sense: take radial large cutoffs, then
let their radii tend to infinity. Each subtracted constant has
integral zero for concentric annuli, and shifting the outer
center by \(h\) changes its integral by \(O(h/R)\) at radius \(R\).
Thus the difference on the complement of \(B_{2h}(x)\) is the
sum of the integrals of
\[
 \bigl(K(x-y)-K(x'-y)\bigr)(g(y)-g(x))
 \quad\hbox{and}\quad
 (g(x')-g(x))K(x'-y).
\]
The first is bounded by
\(C[g]_\alpha h\int_{2h}^\infty r^{\alpha-2}\,dr\).
The integral of \(K(x'-y)\) on this complement, in the stated
limit, is bounded independently of \(h\): subtract the zero
radial integral of \(K(x-y)\), and bound the difference by
\(Ch\int_{2h}^\infty r^{-2}\,dr\).
This proves the Hölder bound for nearby pairs. For \(h>1\) use
the supremum bound.
\end{proof}

\begin{lemma}[The Newtonian Hessian]\label{pde:elliptic-newton-hessian}
For \(n\geq2\), let \(\Phi\) be the fundamental solution of
\(-\Delta\) in Section~\ref{pde:L01} and put
\(\sigma=|S^{n-1}|\). Then
\[
 D_iD_j\Phi=\operatorname{pv}K_{ij}-\frac{\delta_{ij}}{n}\delta_0,
 \qquad
 K_{ij}(x)=\frac1{\sigma}
      \left(\frac{n x_ix_j}{|x|^{n+2}}-
                       \frac{\delta_{ij}}{|x|^n}\right).
\]
If \(F\in C_c^{0,\alpha}\), its Newtonian potential \(\Phi*F\)
is \(C^{2,\alpha}_{\mathrm{loc}}\).
On a fixed bounded set its \(C^{2,\alpha}\) norm is bounded
in terms of the set, a bound for the support, and \(\|F\|_{C^{0,\alpha}}\).
\end{lemma}

\begin{proof}
Away from the origin \(\partial_i\Phi=-x_i/(\sigma|x|^n)\).
Its distributional derivative \(D_j\) is found by integration by
parts outside \(B_\epsilon(0)\). The boundary term is minus
\[
 \frac1{\sigma}\int_{S^{n-1}}\theta_i\theta_j
                                      \phi(\epsilon\theta)\,dS_\theta
\]
against a test \(\phi\), since the normal to the punctured region
on the inner sphere is \(-\theta\).
It tends to \(-\delta_{ij}\phi(0)/n\). The ordinary derivative
away from zero is \(K_{ij}\), whose spherical mean vanishes
by \(\int\theta_i\theta_j\,dS=\sigma\delta_{ij}/n\).
This also gives existence of the distributional principal value.
The identity follows. The trace of this Hessian is \(-\delta_0\),
consistent with the sign of the fundamental solution.

For the potential, local integrability of \(\Phi\) and \(\nabla\Phi\)
on bounded sets, including the logarithm when \(n=2\), bounds its
size and first derivatives there. The Hessian convolution is
\(T_{ij}F-\delta_{ij}F/n\), where \(T_{ij}\) is the operator of
the preceding lemma. To see the equality, convolve the punctured
kernels first, and use the uniform \(O(\epsilon^\alpha)\) bound
on the omitted small-ball integrals after subtracting \(F(x)\).
This also identifies the distributional derivative by testing
and taking the limit. The convolution and its first derivatives
are continuous, by continuity of translations in \(L^1\) of
the locally truncated kernels. Continuous distributional
derivatives agree with the classical derivatives: locally mollify,
use uniform convergence of the function and of the specified
derivative, and pass to the limit in the fundamental theorem of
calculus on coordinate segments. Applying this argument twice
gives the assertion. When \(n=1\), \(D^2\Phi=-\delta_0\),
and the corresponding bound follows directly by integration.
\end{proof}

\begin{theorem}[Interior Schauder estimate for Poisson]\label{pde:elliptic-poisson-schauder}
Suppose \(u\in C^2(\Omega)\), \(f\in C^{0,\alpha}_{\mathrm{loc}}(\Omega)\),
and \(-\Delta u=f\). Then \(u\in C^{2,\alpha}_{\mathrm{loc}}(\Omega)\).
For bounded open \(\Omega'\Subset\Omega''\subset\Omega\),
\[
 \|u\|_{C^{2,\alpha}(\Omega')}
   \leq C\bigl(\|u\|_{C^0(\Omega'')}+
                      \|f\|_{C^{0,\alpha}(\Omega'')}\bigr)
\]
whenever the right-hand side is finite. In particular one may use
the norms on \(\Omega\) when they are finite.
\end{theorem}

\begin{proof}
Choose a smooth compact cutoff \(\chi\) in \(\Omega''\) equal to
one on an intermediate neighborhood of \(\bar\Omega'\).
The zero extension \(F=\chi f\) is globally Hölder, with norm
controlled by the cutoff and the norm of \(f\) on \(\Omega''\).
The Newtonian lemma gives the \(C^{2,\alpha}\) estimate of its
potential \(w\) on bounded sets. Since \(u-w\) is harmonic
where \(\chi=1\), the interior estimates in Section~\ref{pde:L03}
bound its derivatives through order three on a neighborhood
of \(\bar\Omega'\) by its supremum on a slightly larger
neighborhood. Nearby pairs can be joined by segments staying
in that smaller neighborhood, yielding the required Hölder
bound for second derivatives there. For pairs farther apart
use the supremum bound on the derivatives. The theorem follows
by addition.
\end{proof}

On concentric balls the unit-scale proof rescales to a bound on
\(B_r\) from \(B_{2r}\) with the left-hand terms
\[
 \|u\|_\infty+r\|Du\|_\infty+
 r^2\|D^2u\|_\infty+r^{2+\alpha}[D^2u]_\alpha
\]
and right-hand terms
\(\|u\|_\infty+r^2\|f\|_\infty+r^{2+\alpha}[f]_\alpha\).
In two dimensions one scales the equation and applies the unit
estimate, so no additive logarithmic constant is required.
Compare \cite[Section 4.3]{GT}.

\section{The Poisson \texorpdfstring{\(W^{2,p}\)}{W(2,p)} model}\label{pde:E08}

Use the Riesz operators with multiplier \(-i\xi_j/|\xi|\).
Their \(L^p\) bounds for \(1<p<\infty\) are proved in Unit U3
of \cite{P102}. The choice of Fourier normalization leaves
these multipliers unchanged.

\begin{theorem}[Whole-space and local \(W^{2,p}\) estimates]
\label{pde:elliptic-poisson-w2p}
For \(1<p<\infty\) and \(u\in C_c^\infty(\mathbb R^n)\),
\[
 \|D^2u\|_{L^p(\mathbb R^n)}
     \leq C_{n,p}\|\Delta u\|_{L^p(\mathbb R^n)} .
\]
If \(u,f\in L^p_{\mathrm{loc}}(\Omega)\) and
\(-\Delta u=f\) as distributions, then \(u\in W^{2,p}_{\mathrm{loc}}(\Omega)\).
For \(\Omega'\Subset\Omega''\Subset\Omega\),
\[
 \|u\|_{W^{2,p}(\Omega')}
 \leq C\bigl(\|f\|_{L^p(\Omega'')}+\|u\|_{L^p(\Omega'')}\bigr).
\]
\end{theorem}

\begin{proof}
For the whole-space estimate the Fourier multiplier gives
\(D_iD_j u=R_iR_j(-\Delta u)\). This identity holds first in
\(L^2\) by Plancherel and then with both sides interpreted as
\(L^p\) functions, using the same Riesz operators on \(L^2\cap L^p\).
For \(n=1\) the estimate is immediate.

For the local assertion, we first identify the Hessian of the
potential on a compactly supported smooth source. Away from the
frequency origin the relation \(-\Delta\Phi=\delta_0\) gives
the multiplier \(-\xi_i\xi_j/|\xi|^2\) for \(D_iD_j\Phi\).
There is no additional distribution supported at the frequency
origin. Indeed, \(D_iD_j\Phi\) is homogeneous of degree \(-n\)
as a distribution in physical space. This follows by scaling
\(\Phi\) and differentiating twice. In dimension two the
additive constant in the scaling of \(\log|x|\) disappears.
The inverse transform of the bounded homogeneous degree-zero
multiplier is also homogeneous of degree \(-n\).
Their difference is a polynomial, since its Fourier transform
is supported at the single point zero. For this last fact, a
distribution supported at one point has finite order on a fixed
compact neighborhood. Subtract the Taylor jet of that order
from a test and multiply the remainder by cutoffs shrinking
to the point. The seminorms of the remainder through that order
then tend to zero, so only the Taylor jet contributes. This
expresses the distribution as a finite sum of delta derivatives.
A polynomial homogeneous of the negative degree \(-n\) is zero.
Thus for smooth compactly supported \(F\)
\[
 D_iD_j(\Phi*F)=R_iR_jF .
\]
The contact term of the preceding Hessian lemma is part of this
distributional identity.

Now choose a compact cutoff in \(\Omega''\) equal to one on a
neighborhood of \(\bar\Omega'\) and let \(F\) be the cutoff times
\(f\), extended by zero. Approximate \(F\) in \(L^p\) by smooth
\(F_k\) supported in a common compact set, and set \(v_k=\Phi*F_k\).
On each fixed bounded region the local \(L^1\) norms of
\(\Phi,\nabla\Phi\), truncated to cover all relevant differences
of points, and Young's inequality make \((v_k)\) Cauchy in
\(W^{1,p}\) there. Their Hessians are Cauchy in \(L^p\) by the
Riesz bound. The limit \(v=\Phi*F\) is locally \(W^{2,p}\),
with the indicated local bound by \(\|F\|_p\), and
\(-\Delta v=F\). The distribution \(u-v\) is harmonic near
\(\bar\Omega'\). Weyl's lemma and the harmonic derivative
estimates of Sections~\ref{pde:L02}--\ref{pde:L03} bound its
derivatives through order two on \(\Omega'\) by an \(L^1\)
norm on a slightly larger bounded neighborhood. Hölder's
inequality controls this by the corresponding \(L^p\) norm of
\(u\) and \(v\). Adding the estimates proves the result.
\end{proof}

Compare \cite[Chapter 9]{GT} and the Fourier method in
\cite[Section 4.3.1]{Evans}. For the local constant-coefficient
assertion the distribution need not have an initial \(W^{1,p}\)
representative.

\section{Variable-coefficient Schauder theory}\label{pde:E09}

In this section \(L=a^{ij}D_iD_j+b^iD_i+c\) has a symmetric
uniformly positive principal matrix. We give the a priori
estimates before using them to solve the Dirichlet problem.
The flat-boundary estimate also records why simply reflecting
the source oddly can lose Hölder regularity.

\begin{lemma}[Constant coefficients at a flat boundary]
\label{pde:elliptic-flat-schauder}
For the Laplacian in \(B_{2r}^+=B_{2r}(0)\cap\{x_n>0\}\),
let \(u\in C^2(B_{2r}^+)\) be continuous up to the flat face,
vanish there, and satisfy \(\Delta u=f\) with \(f\) Hölder
up to that face. Then the scaled Schauder estimate of the
preceding section holds on \(B_r^+\) including its flat face.
Both this estimate and the interior ball estimate hold for a
constant symmetric positive definite matrix in place of the
identity, with constants controlled by its ellipticity and
upper bounds.
\end{lemma}

\begin{proof}
Scale to \(r=1\). When \(n=1\), integrating the equation twice
bounds the second derivatives and their Hölder seminorm, and a
first difference bounds the first derivative. Suppose \(n\geq2\).
Reflect the source \emph{evenly}, cut it off in \(B_2\) with a
cutoff equal to one on \(B_{3/2}\), and solve the resulting
whole-space Poisson equation by the Newtonian potential with
the sign for \(\Delta\). This gives a \(C^{2,\alpha}\) function
\(w\) on \(B_2\), with its norm bounded by the given source norm.
To correct its trace, choose a compactly supported
\(g\in C^{2,\alpha}(\mathbb R^{n-1})\) equal to \(w(x',0)\)
on the disk of radius \(3/2\), using a cutoff in the disk of
radius \(2\).

We detail the half-space harmonic lift of \(g\).
For \(m=n-1\), write
\[
 P_t(z)=c_m\,\frac{t}{(t^2+|z|^2)^{(m+1)/2}},\qquad t>0,
\]
where \(c_m\) makes the integral one. It is harmonic in
\((z,t)\) for \(t>0\), as is seen by differentiating the
fundamental solution in the normal direction. Scaling gives
\(\int |z|^\alpha P_t(z)\,dz=C_{m,\alpha}t^\alpha\).
For the Fourier phase \(e^{-2\pi i z\cdot\xi}\), its
multiplier in \(z\) is \(e^{-2\pi t|\xi|}\):
harmonicity gives the scalar ordinary differential equation
\(\partial_{tt}\widehat P_t=(2\pi|\xi|)^2\widehat P_t\), the
probability normalization gives boundedness, and the
approximate-identity property gives initial value one.
The zero frequency is constant one. In particular
\(P_t*P_s=P_{t+s}\). For any bounded Hölder function \(b\),
\(P_t*b\) is uniformly Hölder in \(x'\) by integration
against this probability kernel, and
\[
 \|P_{t+s}*b-P_t*b\|_\infty
       \leq \|P_s*b-b\|_\infty\leq C[b]_\alpha s^\alpha .
\]
These estimates include the limit at \(t=0\).

Let \(h(x',t)=P_t*g(x')\).
Its first and second tangential derivatives are the extensions
of the corresponding derivatives of \(g\). Put
\(\Lambda b=\sum_{i=1}^m R_iD_ib\). For compactly supported
\(C^{1,\alpha}\) functions \(b\) this is bounded Hölder
by Lemma~\ref{pde:elliptic-holder-kernel}, since the Riesz
kernels are odd homogeneous kernels of degree \(-m\).
The multiplier identity gives
\(\partial_t h=-P_t*\Lambda g\) and
\(\partial_t D_jh=-P_t*\Lambda D_jg\). These identities
can first be read in \(L^2\) by Plancherel, and then pointwise
for \(t>0\) by continuity. Finally
\(\partial_{tt}h=-P_t*\Delta' g\).
It follows from the kernel estimates that all second
derivatives extend Hölder continuously to the closed
half-space, with bounds controlled by the compact
\(C^{2,\alpha}\) datum. The limiting first derivatives
identify the classical up-to-boundary derivatives by the
fundamental theorem of calculus on one-sided segments.
This constructs the desired lift without a vanishing
assumption on the source at the plane.

Now \(u-w+h\) is harmonic on \(B_{3/2}^+\) with zero flat
boundary value. Its odd reflection is harmonic across the
plane. Indeed, for each small ball centered on the plane,
solve the full-ball harmonic Dirichlet problem using the
continuous odd boundary data of the reflected function.
The solution is odd by uniqueness, vanishes on the plane,
and agrees with the original harmonic function in the
upper half-ball by comparison. The ball Poisson result of
Section~\ref{pde:L05} justifies this reflection argument.
Interior harmonic derivative estimates on the reflected
function, added to the bounds for \(w,h\), prove the flat
estimate.

For a constant positive matrix, a linear change of
coordinates turns its second-order operator into the
Laplacian. A subsequent rotation takes the transformed flat
plane into a coordinate plane. The change and its inverse
have norm controlled by the two ellipticity bounds. Cover
the image of the smaller ball or half-ball by interior
balls and plane-centered half-balls of a fixed comparable
radius in the larger image. The local estimates give the
seminorm for close pairs, and their supremum bounds give it
for pairs farther apart. Scaling back proves the last
assertion.
\end{proof}

\begin{theorem}[Local and global Schauder estimates]
\label{pde:elliptic-variable-estimate}
Suppose the coefficients of \(L\) and \(f\) are locally
\(C^{0,\alpha}\) and \(Lu=f\). For a \(C^{2,\alpha}_{\rm loc}\)
solution the interior \(C^{2,\alpha}\) estimate on
\(\Omega'\Subset\Omega''\) has right-hand side
\(C(\|u\|_{C^0(\Omega'')}+\|f\|_{C^{0,\alpha}(\Omega'')})\)
if the indicated data and coefficient norms are finite.
On a bounded \(C^{2,\alpha}\) domain with coefficients Hölder
on \(\bar\Omega\), if \(u\in C^{2,\alpha}(\bar\Omega)\) has
boundary values of \(G\in C^{2,\alpha}(\bar\Omega)\), then
\[
 \|u\|_{C^{2,\alpha}(\bar\Omega)}
 \leq C\bigl(\|u\|_{C^0(\bar\Omega)}
       +\|f\|_{C^{0,\alpha}(\bar\Omega)}
       +\|G\|_{C^{2,\alpha}(\bar\Omega)}\bigr).
\]
The constants depend on the domains, exponent, ellipticity,
and coefficient norms. A corresponding local estimate holds
up to any \(C^{2,\alpha}\) boundary portion, away from its ends.
\end{theorem}

\begin{proof}
We prove a local estimate on a unit ball or a unit half-ball
with zero flat trace. Write \(d(x)=1-|x|\), the distance to
the spherical boundary, and \(d_{xy}=\min(d(x),d(y))\).
Use the weighted quantities
\[
 M_j=\sup_x d(x)^j|D^j u(x)|\quad(0\leq j\leq2),\qquad
 N_j=\sup_{x\ne y}d_{xy}^{j+\alpha}
       \frac{|D^ju(x)-D^ju(y)|}{|x-y|^\alpha}\quad(0\leq j\leq2).
\]
For every \(\epsilon>0\) there is \(C_\epsilon\) such that
\[
 M_1+M_2+N_0+N_1\leq \epsilon N_2+C_\epsilon M_0.
 \tag{*}\label{pde:elliptic-weighted-interp}
\]
For \(M_2\), take a one-sided second difference at \(x\) with
step \(h d(x)\), \(h\leq1/8\), in each of the directions
\(e_i\) and \((e_i+e_j)/\sqrt2\). Their normal components are
nonnegative, so the segments stay in the half-ball and all
weights on them are comparable to \(d(x)\).
Taylor's formula with integral remainder gives
\(M_2\leq C(h^{-2}M_0+h^\alpha N_2)\); polarization
controls every Hessian entry. A first difference gives
\(M_1\leq C(h^{-1}M_0+hM_2)\). For \(N_j\), \(j=0,1\),
separate pairs with \(|x-y|\leq d_{xy}/2\) and farther
pairs. The segment estimate in the first case uses \(M_{j+1}\)
because the ball and half-ball are convex and
\(d\) along the segment is at least \(d_{xy}\).
The second case uses \(2^{1+\alpha}M_j\).
Choosing \(h\) small for the estimate of \(M_2\) proves
\eqref{pde:elliptic-weighted-interp}.

For a pair \(x,y\) assume \(d(x)=d_{xy}=d\) and fix
\(0<\mu<1/32\). If \(|x-y|\geq\mu d/8\), its
contribution to \(N_2\) is bounded by \(C\mu^{-\alpha}M_2\).
Otherwise use an interior ball of radius comparable to
\(r=\mu d\) containing both points in its smaller ball,
or a half-ball of comparable radius centered at \((x',0)\)
if \(x_n<2r\). For the latter case one can take outer
radius \(8r\) and inner radius \(4r\).
The weights throughout the outer region are comparable to \(d\).
Freeze the principal coefficients at \(x\), so the new
right-hand side is
\(F=f+(a(x)-a)D^2u-bDu-cu\).
The constant-coefficient estimate gives, for the weighted
contribution of the pair,
\[
 C(d/r)^{2+\alpha}
       \bigl(M_0+r^2\|F\|_\infty+r^{2+\alpha}[F]_\alpha\bigr).
\]
Let \(K\) bound the Hölder norms of the coefficients.
On this local region \(\|a-a(x)\|_\infty\leq CKr^\alpha\).
The term in the product Hölder estimate using this small
supremum contributes at most \(CK\mu^\alpha N_2\);
the term using \([a]_\alpha\|D^2u\|_\infty\), and the product
supremum, contribute at most \(C(\mu,K)M_2\).
For example, the Hölder term with the small factor is
bounded by \(C d^{2+\alpha}Kr^\alpha
d^{-2-\alpha}N_2\).
Using the weighted norms for \(bDu\) and \(cu\) gives a
bound by \(C(\mu,K)(M_0+M_1+N_0+N_1)\).
The data term is at most \(C(\mu)\|f\|_{C^{0,\alpha}}\).
Taking the supremum over all pairs now gives
\[
 N_2\leq CK\mu^\alpha N_2
    +C(\mu,K)\bigl(\|f\|_{C^{0,\alpha}}+M_0+
                                     M_1+M_2+N_0+N_1\bigr).
\]
First choose \(\mu\) small and then use
\eqref{pde:elliptic-weighted-interp} with small \(\epsilon\)
to absorb the two terms containing \(N_2\). This proves
the local estimate on a smaller ball or half-ball.
If the weighted norms might initially be infinite, apply
the proof on concentric radii increasing to one, where
they are finite by local \(C^{2,\alpha}\) regularity,
including the flat face in the boundary case.
For radii between \(1/2\) and \(1\) the scaled constants
are uniform. Pass to the limit for each pair.

A \(C^{2,\alpha}\) chart straightening a curved boundary
preserves the ellipticity bounds up to fixed factors, and
the chain rule yields \(C^{0,\alpha}\) transformed
coefficients, including the new first-order terms.
Subtract the boundary lift and apply the flat local
estimate. The transformed source is bounded by
\(C(\|f\|_{C^{0,\alpha}}+\|G\|_{C^{2,\alpha}})\).
Interior and boundary estimates in a finite cover give
the global supremum norms. Every sufficiently close pair
in the compact closure lies together in one smaller patch,
and pairs farther apart are controlled by these supremum
norms, giving the global Hölder seminorm. The same covering
argument on a compact interior region gives the nested
interior estimate.
\end{proof}

\begin{theorem}[Classical Dirichlet solvability]
\label{pde:elliptic-schauder-solvability}
Let \(\Omega\) be bounded with \(C^{2,\alpha}\) boundary.
Suppose \(a,b,c,f\in C^{0,\alpha}(\bar\Omega)\), the
principal matrix is uniformly positive, and \(c\leq0\).
Given \(G\in C^{2,\alpha}(\bar\Omega)\), there is a unique
\(u\in C^{2,\alpha}(\bar\Omega)\) satisfying \(Lu=f\) with
boundary values \(G\). In this setting the estimate above
can omit the \(\|u\|_{C^0}\) term, by adjusting its constant.
The local interior \(C^{2,\alpha}\) assertion and estimate
also hold for a \(C^2\) solution with locally Hölder
coefficients and source, without a sign restriction on \(c\).
\end{theorem}

\begin{proof}
We begin with the Laplacian on a ball. A Hölder right-hand
side on the closed ball extends with bounded Hölder norm
to a compactly supported function, by composition with
the nearest-point projection onto the ball and a cutoff.
Its whole-space Newtonian potential, with the sign
for \(\Delta\), is \(C^{2,\alpha}\) near the closed ball.
The ball Poisson integral then corrects the boundary
values, giving a classical interior solution continuous
on the closure. Subtract the prescribed lift so the
remaining boundary values are zero. To see its up-to-boundary
regularity, invert about another point of the sphere.
In dimension \(n\geq2\) the resulting Kelvin transform
\(v(y)=|y|^{2-n}u(y/|y|^2)\), in translated coordinates,
satisfies
\[
 \Delta v(y)=|y|^{-n-2}(\Delta u)(y/|y|^2).
\]
For this identity \(D(y/|y|^2)=|y|^{-2}
(I-2\hat y\otimes\hat y)\); the displayed matrix in
parentheses is orthogonal, and the first-order terms from
differentiating the inversion cancel those from
\(D(|y|^{2-n})\). The weight is harmonic off the origin.
The inversion maps the sphere minus its pole to a plane.
Near any other boundary point the coordinate change and
weight are smooth with bounded derivatives, so the flat
estimate gives \(C^{2,\alpha}\) there, then compactness
of the sphere gives the global class. For an interval the
same conclusion follows by integrating twice.

On a ball, use \(L_s=(1-s)\Delta+sL\) on the space
\(X=\{v\in C^{2,\alpha}(\bar B):v|_{\partial B}=0\}\)
with values in \(Y=C^{0,\alpha}(\bar B)\).
The estimates are uniform for \(0\leq s\leq1\).
To absorb the supremum term, on a slab containing the
domain choose \(d>\sup x_1\) and
\(\psi=e^{kd}-e^{kx_1}>0\). For large \(k\) there is a
uniform \(\delta>0\) with \(L_s\psi\leq-\delta\).
Indeed \(L_s e^{kx_1}\geq
(\lambda_0 k^2-Bk-C)e^{kx_1}\) for common coefficient
bounds, and \(L_s(e^{kd})\leq0\).
Comparison with multiples of \(\psi\) yields
\(\|v\|_\infty\leq C\|L_s v\|_\infty\) for zero boundary.
Thus \(\|v\|_X\leq C\|L_s v\|_Y\).
Invertibility is open by the convergent geometric operator
series. It is closed along the path: if \(s_j\to s\) and
\(L_{s_j}v_j=f\), the uniform estimate bounds \(v_j\) in \(X\)
and
\(\|v_j-v_k\|_X\leq C|s_j-s_k|
\|(L-\Delta)v_k\|_Y\).
Hence \(v_j\) converges in the Banach space \(X\) and solves
the limiting equation. Injectivity follows from the same
estimate. The Laplacian start therefore gives ball
solvability for \(L\), with arbitrary lifted Hölder boundary
data by subtraction.

A ball solution also exists for any continuous boundary
data \(g\), with \(C^{2,\alpha}\) regularity in the interior.
Approximate \(g\) uniformly by smooth boundary data.
The maximum principle applied to differences makes the
solutions uniformly Cauchy on the closure. The interior
estimate makes them Cauchy in \(C^{2,\alpha}\) on smaller
balls, since their equation differences have zero right-hand
side. If \(g\) is \(C^{2,\alpha}\) only on an open spherical
portion, the approximants can be chosen with uniformly
bounded \(C^{2,\alpha}\) norms on smaller portions.
To do this, multiply \(g\) by a cutoff supported in the
portion and smooth in boundary charts. Smooth the remaining
continuous data uniformly while keeping it zero on a still
smaller portion. Local normal extensions of these smooth
data have uniformly bounded \(C^{2,\alpha}\) norms there.
Subtract those lifts and use the local boundary estimate.
On a smaller closed patch the functions and their first
two derivatives have a uniformly Hölder bound at the
highest order. Arzelà--Ascoli and integration along local
segments identify the limiting derivatives and preserve
the Hölder bound. This proves the claimed local boundary
regularity for continuous ball data.

We next construct a solution on the general domain with
zero data, since subtracting \(G\) reduces to that case.
A continuous function \(v\) on \(\bar\Omega\), with
\(v\leq0\) on \(\partial\Omega\), is a subfunction if on
every ball compactly contained in \(\Omega\) it is dominated
by the solution of \(Lw=f\) with its own boundary values.
Every \(C^2\) subsolution \(Lv\geq f\) with the required trace
is a subfunction. The maximum of finitely many subfunctions
is one, by comparison of the corresponding ball solutions.
Replacing a subfunction on a ball by that solution gives
a larger continuous subfunction. To check comparison on
a test ball meeting the replaced ball, first compare the
original function with the new test-ball solution, then
compare the two classical solutions on the intersection
of the balls. On the interface their inequalities follow
from the original comparison and the ordered test-ball
boundary values. The weak classical maximum principle
applies on the intersection, with no boundary smoothness
needed. This proves the replacement assertion.

A subfunction is dominated by any smooth supersolution
with ordered outer boundary values. In fact add
\(\epsilon\psi\) to the supersolution to make its inequality
strict. If the difference has a positive interior global
maximum \(M\), compare on a small ball about that point
with the strict supersolution plus \(M\). Since \(cM\leq0\),
the ball solution is strictly below this comparison function
at the center: an interior equality would be a zero
maximum of their difference, with strictly positive \(L\),
which is impossible. This contradicts the defining
subfunction comparison. Let \(\epsilon\downarrow0\).

Choose \(M\) so large that \(-M\psi\) is a classical
subsolution and \(M\psi\) a classical supersolution.
The supremum \(U\) of all subfunctions is finite and
between these bounds. On any interior ball select a
countable dense set of points. For each \(j\), take a finite
maximum of subfunctions within \(1/j\) of \(U\) at each of
the first \(j\) points, and include the previous maximum.
Replace these functions on the ball. The replacements
are monotone, bounded, and satisfy \(Lw=f\) on that ball.
By the interior estimates a subsequence converges locally
with its first two derivatives to a \(C^{2,\alpha}_{\rm loc}\)
solution \(w\). Monotonicity gives the same pointwise
limit, and \(w\leq U\), with equality on the dense set.
The supremum \(U\) is lower semicontinuous, so the open
set \(U>w\) must be empty. Therefore \(U\) solves the
equation in every interior ball.

At each boundary point \(p\) the graph and its bounded
second derivatives give an exterior tangent ball
\(B_R(y)\) with closure touching \(\bar\Omega\) only at \(p\).
Indeed a sufficiently small sphere lies below the local
boundary graph by the quadratic bound, and can be chosen
within its chart. Reducing its radius along the same
normal makes its contact strict away from \(p\).
The function
\(\beta=R^{-q}-|x-y|^{-q}\) is nonnegative on the closure,
zero at \(p\), and for large \(q>0\) has \(L\beta\leq-\delta_p<0\).
To check this, write \(z=x-y\) and \(\rho=|z|\).
The derivative terms, divided by \(q\rho^{-q-2}\),
have upper bound
\(\operatorname{tr}a-(q+2)\lambda+\rho|b|\), while
\(c\beta\leq0\). Since \(\Omega\) is bounded and
\(\rho\geq R>0\), the desired strict uniform bound follows.
Multiples \(-C\beta\) and \(C\beta\) are a subfunction
and a supersolution respectively, for \(C\) sufficiently
large compared to \(\|f\|_\infty\).
They bound \(U\), proving the zero boundary limit.

To raise regularity at a curved boundary point, transfer
a neighborhood to one in which a small ball lies inside
the domain and shares an open spherical patch with its
boundary. Here is an explicit chart. After flattening
write the domain locally as \(x_n>0\).
The lower boundary of \(B_R(Re_n)\) is the graph
\(b_R(x')=R-\sqrt{R^2-|x'|^2}\) for \(|x'|<R\).
Add to \(x_n\) a smooth cutoff of \(b_R\), equal to \(b_R\)
for \(|x'|<R/4\), supported in \(|x'|<R/2\), and between
zero and \(b_R\). The resulting graph lies below the
sphere and agrees with it near the contact point.
For small \(R\) the whole ball fits in the chart.
The transformed equation has Hölder coefficients and
\(c\leq0\). On the ball the existing solution is
continuous up to the sphere and has lifted regular
data on the shared patch. By uniqueness it equals the
ball solution, whose local boundary regularity was
just proved. Pulling back gives the required regularity
on a neighborhood of the boundary point in the closure.
In one dimension a local interval with a shared endpoint
serves the same purpose.
A finite cover gives \(U\in C^{2,\alpha}(\bar\Omega)\).
Adding back \(G\) gives the solution. Uniqueness and the
supremum bound follow from classical comparison and the
slab barrier, completing the global estimate.

Finally let a local solution be initially only \(C^2\).
On a sufficiently small ball centered at \(x_0\), take
\(\chi=1-k|x-x_0|^2>0\) with \(L\chi<0\).
Choose \(k\) large relative to the local bound of \(c\)
and then the radius small relative to \(k\) and the
bound of \(b\); the principal negative term
\(-2k\operatorname{tr}a\) dominates the other terms.
For \(u=\chi v\), the equation divided by \(\chi\)
has the same principal coefficients and a nonpositive
zero-order coefficient \(L\chi/\chi\).
Solve that ball problem with the continuous boundary
data of \(v\). Classical comparison identifies \(v\)
with this locally \(C^{2,\alpha}\) solution. Thus
\(u\) has the regularity needed for the already proved
interior estimate with its original coefficients.
\end{proof}

The estimates and existence argument correspond to
\cite[Theorems 6.2, 6.6, and 6.13--6.14]{GT}.
The constant-coefficient flat estimate uses the Poisson
extension to handle the boundary trace of an even-source
Newtonian potential. The global \(C^{2,\alpha}\) result is a
classical nondivergence result with Hölder coefficients
and data, separate from the \(C^{1,1}\) divergence-form
\(H^2\) theorem in Section~\ref{pde:E06}.

\section{Closed forms and the Dirichlet spectrum}\label{pde:E10}

In the spectral sections the Hilbert inner product \((u,v)\) is
linear in \(u\). One may use real spaces for the symmetric problem
or their complexifications. For a dense form domain \(V\subset H\),
define the operator associated with the form \(a\) by
\[
 \begin{split}
 D(A_H)&=\{u\in V:\ \exists f\in H,\ 
                    a(u,v)=(f,v)\ \text{for all }v\in V\},\\
 A_Hu&=f.
 \end{split}
\]
Density makes \(f\) unique.

\begin{theorem}[Closed forms and the Friedrichs extension]
\label{pde:elliptic-friedrichs}
Let \(a\) be a dense nonnegative symmetric form in \(H\),
closed in the sense that its domain \(V\) is complete for
\(\|u\|_V^2=a(u,u)+\|u\|_H^2\).
Then \(A_H\) is nonnegative self-adjoint and
\((I+A_H)^{-1}\) is the solution map of the shifted form.
If \(V\hookrightarrow H\) is compact, this inverse is compact.
A densely defined symmetric operator bounded below has a
self-adjoint extension with the same lower bound, obtained
by closing its shifted quadratic form.
\end{theorem}

\begin{proof}
The shifted form is the Hilbert inner product on \(V\).
For \(f\in H\), Hilbert representation gives a unique \(Sf\in V\)
with \(a(Sf,v)+(Sf,v)=(f,v)\) for all \(v\in V\), and
\(\|Sf\|_V\leq\|f\|_H\). Using two such solutions as tests
shows \(S=S^*\) on \(H\), and taking one solution as its own
test gives positivity. Its kernel is zero by density of \(V\),
so its range is dense. By definition
\(D(A_H)=\operatorname{Ran}S\) and \(A_H=S^{-1}-I\).
The inverse of a bounded injective self-adjoint operator with
dense range is self-adjoint on its range. Indeed, if
\((S^{-1})^*v=w\), test against \(Sx\) for all \(x\in H\)
to get \(v=Sw\), and the converse follows from \(S=S^*\).
Also \((A_Hu,u)=a(u,u)\geq0\). Compactness follows from the
bound of \(S\) into \(V\) and compactness of the inclusion.

For a symmetric \(T\geq mI\) shift to \(T_1=T+(1-m)I\geq I\).
The form \(q(u,v)=(T_1u,v)\) on \(D(T)\) is closable.
Indeed, if \(u_k\) is Cauchy in the form norm and tends
to zero in \(H\), then \(q(u_k,v)=(u_k,T_1v)\to0\) for each
\(v\in D(T)\). Its limit in the form-norm completion is
orthogonal to that dense initial domain, hence is zero.
Thus the completion maps injectively into \(H\), gives a
dense closed form domain, and still satisfies
\(q(u,u)\geq\|u\|_H^2\). For \(u\in D(T)\) the identity
\(q(u,v)=(T_1u,v)\) extends from \(D(T)\) to the completion.
The associated operator therefore extends \(T_1\) and is
bounded below by \(I\). Undoing the shift gives the result.
\end{proof}

\begin{theorem}[Dirichlet spectrum and min--max]
\label{pde:elliptic-dirichlet-spectrum}
Let \(\Omega\subset\mathbb R^n\), \(n\geq1\), be nonempty,
bounded and Lipschitz, and let \(A\) be bounded, real symmetric
and uniformly elliptic. For \(H=L^2(\Omega)\), \(V=H_0^1(\Omega)\)
and \(a(u,v)=\int_\Omega A\nabla u\cdot\overline{\nabla v}\),
the associated \(L\) has compact positive self-adjoint inverse
on \(H\). There is an orthonormal basis \((e_j)_{j\geq1}\)
with \(Le_j=\lambda_j e_j\), where
\(0<\lambda_1\leq\lambda_2\leq\cdots\) and
\(\lambda_j\to\infty\), repeating eigenvalues according to
their finite multiplicities. Furthermore
\[
 \lambda_j=\min_{\substack{W\subset H_0^1(\Omega)\\ \dim W=j}}
                  \max_{0\ne u\in W}\frac{a(u,u)}{\|u\|_2^2}.
\]
\end{theorem}

\begin{proof}
Ellipticity and the slab inequality make \(a\) a Hilbert
inner product equivalent to the usual \(H_0^1\) norm.
The unshifted variational problem supplies \(Sf\in V\) for
each \(f\in H\), with a bounded solution map into \(V\).
As in the shifted proof, \(S\) on \(H\) is positive
self-adjoint and injective with dense range, and
\(S^{-1}=L\) is positive self-adjoint on that range.
Rellich compactness makes \(S\) compact. The inclusion
can also be proved on any bounded open set by zero
extension, uniform \(L^2\) translation estimates, and
bounded common support, as in \cite[V08.01]{P106}.
On its domain \(L\) is the distribution
\(-\operatorname{div}(A\nabla u)\): equality on compact
tests extends to \(H_0^1\) by density.

The compact self-adjoint spectral theorem gives a complete
orthonormal eigenbasis for \(S\), since its kernel is zero.
Its eigenvalues are positive and of finite multiplicity.
The space \(L^2(\Omega)\) is infinite dimensional, by taking
functions supported on disjoint small balls inside \(\Omega\).
Consequently the eigenvalues of \(S\) tend to zero and their
reciprocals are the listed eigenvalues of \(L\).
The domain \(\operatorname{Ran}S\) is exactly the set of
\(u\in H\) with \(\sum_j\lambda_j^2|(u,e_j)|^2<\infty\),
by the diagonal formula for \(S\).

In the form Hilbert space the vectors \(e_j/\sqrt{\lambda_j}\)
are orthonormal and complete: a form-orthogonal \(u\) has
\(a(u,e_j)=\lambda_j(u,e_j)=0\) for every \(j\), so is zero.
Therefore \(a(u,u)=\sum_j\lambda_j|(u,e_j)|^2\) on \(V\).
The span of the first \(j\) eigenvectors has maximum
Rayleigh quotient \(\lambda_j\). Every \(j\)-dimensional
\(W\subset V\) has a nonzero vector orthogonal in \(H\) to
the first \(j-1\) eigenvectors, whose Rayleigh quotient
is at least \(\lambda_j\) by this expansion.
\end{proof}

Compare \cite[Section 6.5.1]{Evans},
\cite[Section 9.3.2]{RR}, and \cite[Chapter 8]{TaylorII}.
The abstract Friedrichs and compact spectral arguments
are also developed in \cite{D16}. The form domain controls
energy, while the operator domain requires the form to
be represented by an \(L^2\) source.

\section{Spectral parameters and Fredholm solvability}\label{pde:E11}

\begin{theorem}[The shifted Fredholm argument]
\label{pde:elliptic-fredholm}
Let \(\Omega\) be bounded Lipschitz, \(H=L^2(\Omega)\), and
\(V=H_0^1(\Omega)\). Let \(A\) be bounded and uniformly
elliptic, possibly nonsymmetric, and let \(b,d,c\) be
bounded. Use the form
\[
 a(u,v)=\int_\Omega \bigl(
 A\nabla u\cdot\overline{\nabla v}
       +(b\cdot\nabla u)\bar v+u\,d\cdot\overline{\nabla v}
       +c u\bar v\bigr).
\]
For complex coefficients assume the principal form has
real part at least \(\lambda_0|\xi|^2\) on complex \(\xi\).
The associated \(L\) is densely defined and closed in \(H\),
and \((L+\mu I)^{-1}\) exists and is compact for sufficiently
large real \(\mu\).
For any \(\lambda\in\mathbb C\), the equation
\((L-\lambda I)u=f\), \(f\in H\), is solvable exactly when
\[
 f\perp\ker(L^*-\bar\lambda I).
\]
Both kernels are finite dimensional, and solutions are unique
modulo \(\ker(L-\lambda I)\). If \(\lambda\) is not an
eigenvalue, there is a unique solution for every \(f\).
\end{theorem}

\begin{proof}
The first-order terms cost arbitrarily small multiples
of \(\|\nabla u\|_2^2\) plus fixed multiples of \(\|u\|_2^2\)
by Cauchy--Schwarz and Young's inequality. The zero-order
term costs another multiple of \(\|u\|_2^2\). Thus
\[
 \operatorname{Re}a(u,u)\geq
       \tfrac12\lambda_0\|\nabla u\|_2^2-C\|u\|_2^2.
\]
Choose \(\mu>C\) so \(a_\mu=a+\mu(\cdot,\cdot)\) is coercive
in the full \(H_0^1\) norm. The real or complex Lax--Milgram
theorem gives a bounded map \(S:H\to V\) with
\(a_\mu(Sf,v)=(f,v)\). On \(H\) it is compact by Rellich
and injective by density of \(V\).
If \(y\perp\operatorname{Ran}S\), the adjoint coercive form
gives \(z\in V\) with \(a_\mu(w,z)=(w,y)\) for \(w\in V\).
Putting \(w=Sf\) gives \((f,z)=(Sf,y)=0\) for all \(f\in H\).
Thus \(z=0\), and then \(y=0\) by density.
It follows that \(D(L)=\operatorname{Ran}S\) is dense and
\(L+\mu I=S^{-1}\) by the form definition.
This inverse is closed, since \(u_j\to u\) and
\((L+\mu I)u_j\to g\) imply
\(u=\lim S((L+\mu I)u_j)=Sg\).

Set \(z=(L+\mu I)u\). The equation becomes
\((I-(\lambda+\mu)S)z=f\).
The identity minus a compact operator has closed range,
Fredholm index zero, and finite-dimensional kernel.
Its range is therefore the orthogonal complement of
\(\ker(I-(\bar\lambda+\mu)S^*)\).
The adjoint resolvent satisfies \(S^*=(L^*+\mu I)^{-1}\):
for \(v\in D(L)\), \(((L+\mu I)v,S^*y)=(v,y)\),
and for \(w\in D(L^*)\) the defining adjoint identity
with \(v=Sx\) gives \(w=S^*(L^*+\mu I)w\).
Thus \(\ker(I-(\bar\lambda+\mu)S^*)=
\ker(L^*-\bar\lambda I)\), including the zero-kernel case
\(\lambda+\mu=0\). The map \(u=Sz\) identifies the two
unadjointed kernels. Index zero makes the range all of
\(H\) when that kernel is zero, and the differences of
solutions are exactly the kernel elements.
\end{proof}

\begin{corollary}[Symmetric spectral compatibility]
\label{pde:elliptic-symmetric-fredholm}
For a coercive symmetric Dirichlet realization on a bounded
Lipschitz domain with compact \(L^2\) resolvent, and real
\(\lambda\), \((L-\lambda I)u=f\) is uniquely solvable for
all \(f\in L^2\) if \(\lambda\) is not an eigenvalue.
At an eigenvalue it is solvable exactly when \(f\) is
orthogonal to its eigenspace, with uniqueness modulo that
eigenspace.
\end{corollary}

\begin{proof}
Self-adjointness identifies the adjoint kernel with the
same eigenspace. For a realization given by a coercive
symmetric form, coercivity gives its inverse at zero.
If compactness is known at another resolvent point,
the resolvent identity makes the inverse at zero compact
as well. Use that inverse in the same factorization.
\end{proof}

The compact alternative used above is the Riesz--Schauder
or Atkinson result in \cite{D16}. Compare
\cite[Section 6.2.3]{Evans}, \cite[Sections 9.2.3 and 9.3.1]{RR},
and \cite[Section 8.2]{GT}. Coercivity belongs to the shifted
operator. Compactness and the adjoint kernel determine
solvability at the spectral parameter.
\chapter{The Heat Equation and Parabolic Energy}
\label{chap:heat}

We use the sign convention
\[
  P=\partial_t-\Delta,\qquad
  \Delta=\sum_{j=1}^n\partial_{x_j}^2,\qquad n\geq 1.
\]
For an elliptic realization denoted by \(B\geq0\) the evolution
equation is \(u'+Bu=f\), so the semigroup generator is \(-B\).
A time-dependent function satisfying \(Pu=0\) is called caloric.
Integrals taking values in a Banach space are Bochner integrals.

\section{The heat kernel and the Cauchy problem}
\label{pde:H01}

\begin{theorem}[Gaussian kernel]
\label{pde:heat-gaussian}
For \(t>0\) define
\[
 K_t(x)=(4\pi t)^{-n/2}\exp\bigl(-|x|^2/(4t)\bigr).
\]
Then \(K_t\geq0\), \(\int_{\mathbb R^n}K_t=1\),
\(\partial_tK_t=\Delta K_t\), and \(K_t*K_s=K_{t+s}\) for \(s,t>0\).
The locally integrable function \(\Gamma(x,t)=K_t(x)\) for \(t>0\),
extended by zero for \(t\leq0\), satisfies \(P\Gamma=\delta_{(0,0)}\)
in distributions. The operators \(S(t)g=K_t*g\), with \(S(0)g=g\),
give a strongly continuous contraction semigroup on
\(L^p(\mathbb R^n)\) for \(1\leq p<\infty\).
They are also contractions on \(L^\infty\).
For bounded continuous \(g\), \(S(t)g\to g\) locally uniformly
as \(t\downarrow0\), and the convergence is uniform on \(\mathbb R^n\)
if \(g\) is bounded and uniformly continuous.
\end{theorem}

\begin{proof}
The one-dimensional Gaussian integral is obtained by squaring it and
integrating in polar coordinates on \(\mathbb R^2\).
Applying it in each coordinate proves the mass identity.
Differentiation gives
\[
 \partial_tK_t(x)
   =\left(-\frac{n}{2t}+\frac{|x|^2}{4t^2}\right)K_t(x)
   =\Delta K_t(x).
\]
For the convolution, complete the square using
\[
 \frac{|x-y|^2}{t}+\frac{|y|^2}{s}
 =\frac{t+s}{ts}\left|y-\frac{s}{t+s}x\right|^2
     +\frac{|x|^2}{t+s}.
\]
The same Gaussian integral gives the normalizing factor of \(K_{t+s}\).
The scaling \(K_t(x)=t^{-n/2}K_1(x/\sqrt t)\) implies that the mass
outside every fixed ball tends to zero. In particular
\(\int K_\varepsilon(x)\varphi(x,\varepsilon)\,dx
\to\varphi(0,0)\) for a compactly supported smooth space-time test.
Integrating by parts in \(t\) from \(\varepsilon\) onwards and in space,
\[
 \int_\varepsilon^\infty\int_{\mathbb R^n}
     K_t(x)(-\partial_t\varphi(x,t)-\Delta\varphi(x,t))\,dx\,dt
   =\int_{\mathbb R^n}K_\varepsilon(x)\varphi(x,\varepsilon)\,dx .
\]
The integral over \(0<t<\varepsilon\) is bounded by \(\varepsilon\)
times the supremum of the differentiated test. This proves the
distributional assertion and local integrability follows from the
mass identity.

Young's inequality proves contraction, including \(p=\infty\).
For \(p<\infty\),
\[
 \|S(t)g-g\|_p
  \leq \int_{\mathbb R^n}K_t(y)\|g(\cdot-y)-g\|_p\,dy .
\]
Translations are norm continuous on \(L^p\), by density of compactly
supported smooth functions and the translation isometry. Split the
integral into a small ball, on which the integrand norm is small, and
its complement, on which it is at most \(2\|g\|_p\). The vanishing
tail mass proves convergence at zero, and contraction and the
semigroup law then give continuity at every time. For bounded
continuous \(g\), on each compact set uniform continuity on its
fixed compact neighborhood gives the same small-ball argument
pointwise and uniformly in that set. Global uniform continuity
permits the argument uniformly in all of space.
\end{proof}

\begin{theorem}[Classical Cauchy solution on the whole space]
\label{pde:heat-cauchy}
Let \(T>0\), \(g\in C_b(\mathbb R^n)\), and
\(f\in C([0,T];C_b^2(\mathbb R^n))\), where the latter space has the
supremum norms of the function and its derivatives through order two.
Then
\[
 u(t,x)=S(t)g(x)+\int_0^t S(t-s)f(s)(x)\,ds
 \qquad (0<t\leq T)
\]
is continuous up to \(t=0\), has \(u(0,x)=g(x)\), and satisfies
\(Pu=f\) classically for \(0<t\leq T\), with a left derivative at \(T\).
It is the unique such solution bounded on \(\mathbb R^n\times[0,T]\).
\end{theorem}

\begin{proof}
The initial-data term is smooth for \(t>0\) by differentiation
under the Gaussian integral and solves the homogeneous equation.
It has the asserted initial values by Theorem~\ref{pde:heat-gaussian}.
Write \(v(t)\) for the integral term. Its absolute value is at most
\(t\sup_s\|f(s)\|_\infty\). Passing the spatial derivatives to \(f\)
is allowed by their uniform bounds. In particular
\(\Delta v(t)=\int_0^t S(t-s)\Delta f(s)\,ds\).
For \(r>0\), \(\partial_r S(r)f(s)=S(r)\Delta f(s)\), by the
kernel equation and integration by parts. The decay of the kernel and
bounded derivatives of \(f\) make the spatial boundary terms tend to
zero on expanding balls, or one can use smooth cutoffs.
This derivative is bounded uniformly by \(\sup_s\|\Delta f(s)\|_\infty\).
It permits differentiation in \(t\) under the integral away from its
upper endpoint, then across the endpoint by dominated convergence.
The endpoint contributes \(f(t)\), since \(S(r)f(s)\to f(t)\) locally
uniformly as \((r,s)\to(0,t)\). Thus \(v_t=\Delta v+f\).
The same bounds and local uniform convergence give continuity of all
derivatives needed in the equation, including the one-sided endpoint.

For uniqueness let \(w\) be the bounded difference of two solutions.
For \(\epsilon>0\) put
\(w_\epsilon=w-\epsilon(1+|x|^2+2nt)\). The subtracted function is
caloric. On the bottom \(w_\epsilon<0\), and for each sufficiently
large \(R\) the same holds on \(\partial B_R\times[0,T]\), by the
boundedness of \(w\). The bounded-cylinder maximum principle
of Section~\ref{pde:H03} gives \(w_\epsilon\leq0\) on that cylinder.
Every fixed point lies in one of these cylinders. Let
\(\epsilon\downarrow0\) and then apply the argument to \(-w\).
\end{proof}

These are the Gaussian representation and uniqueness in a bounded
growth class. The norm statement at the initial time on \(L^p\) is
independent of the classical spatial continuity assumptions.
See \cite[\S\S 2.3.1 and 2.3.3]{Evans},
\cite[Chapter 6, \S 2]{TaylorI}, and the heat semigroup in \cite{D16}.

\section{Smoothing and local estimates for caloric distributions}
\label{pde:H02}

\begin{theorem}[Heat smoothing]
\label{pde:heat-lp-smoothing}
Let \(1\leq p\leq q\leq\infty\), \(g\in L^p(\mathbb R^n)\), let
\(\alpha\) be a spatial multi-index, and let \(j\geq0\) be an integer.
The function \(S(t)g\) has a smooth representative for \(t>0\), and
\[
 \bigl\|\partial_t^j D_x^\alpha S(t)g\bigr\|_{q}
 \leq C(n,\alpha,j,p,q)\,
 t^{-j-|\alpha|/2-(n/2)(1/p-1/q)}\|g\|_p .
\]
Here \(1/\infty=0\). The time derivatives agree with the corresponding
powers of \(\Delta\).
\end{theorem}

\begin{proof}
Repeated differentiation writes
\(\partial_t^jD_x^\alpha K_t(x)\) as
\(t^{-n/2-j-|\alpha|/2}\) times a polynomial in \(x/\sqrt t\)
multiplied by a fixed Gaussian in that variable.
Thus its \(L^r\) norm is at most
\[
 C\,t^{-j-|\alpha|/2-(n/2)(1-1/r)}\qquad(1\leq r\leq\infty).
\]
On a compact range of positive times and of evaluation points, all
differentiated kernels are dominated by a polynomial times a
Gaussian in the integration variable. That majorant belongs to
the conjugate space of \(L^p\), including both endpoints.
It justifies differentiation and continuity pointwise.
Now apply Young's inequality with \(1+1/q=1/p+1/r\).
The equality of time and Laplace derivatives follows by differentiating
the kernel identity. Compare \cite[\S\S 2.3.1 and 2.3.3]{Evans}
and the convolution estimates of \cite[Unit U2]{P102}.
\end{proof}

\begin{theorem}[Interior caloric regularity and scaling]
\label{pde:heat-local}
Let \(u\) be a distribution on an open set \(\mathcal O\subset\mathbb
R^n\times\mathbb R\) and \(Pu=0\) on \(\mathcal O\). Then \(u\) is
represented there by a smooth caloric function. If
\(\overline B_r(x_0)\times[t_0-r^2,t_0]\subset\mathcal O\), then
\[
 |\partial_t^j D_x^\alpha u(x_0,t_0)|
 \leq C(n,\alpha,j)\,r^{-n-2-|\alpha|-2j}
       \int_{t_0-r^2}^{t_0}\int_{B_r(x_0)} |u(y,s)|\,dy\,ds .
\]
The estimate holds in particular at each point of a smaller cylinder
when a larger cylinder is compactly contained in \(\mathcal O\),
by choosing a common admissible radius at all those points.
\end{theorem}

\begin{proof}
Fix a smooth compactly supported cutoff \(\chi\) on \(\mathcal O\)
equal to one near the point to be studied. Extend \(\chi u\) by zero.
The compact support of this factor permits convolution with
\(\Gamma\), and the differentiation rule for convolution gives
\[
 \chi u=(P\Gamma)*(\chi u)=\Gamma*P(\chi u),\qquad
 P(\chi u)=(\chi_t+\Delta\chi)u
                    -2\sum_i\partial_{x_i}\bigl((\partial_{x_i}\chi)u\bigr).
\]
The second identity follows from the distributional product rule
and \(Pu=0\). Its right side is supported where \(\chi\) varies,
away from a neighborhood of the point. The causal kernel is smooth
off the space-time origin, including across \(t=0\) at \(x\ne0\):
all its derivatives there tend to zero from the positive side by
Gaussian decay. In a smaller neighborhood the convolution can
therefore be evaluated by applying the compactly supported
distribution to a smooth translate of this kernel, after multiplying
by a cutoff equal to one near the distribution's support.
Differentiation of those smooth tests proves smoothness of the
convolution there.

For the estimate, translate the point to the origin and use the
parabolic rescaling \(x=rX,\ t=r^2s\). It suffices to estimate at
the center for \(r=1\). Take a product cutoff supported in
\(B_1(0)\) and at times later than \(-1\), equal to one on
\(B_{1/2}(0)\times[-1/4,0]\) and in a neighborhood above time zero.
The compact containment in the hypothesis permits the latter
extension, possibly with a short positive-time cutoff. Terms at
positive times make no contribution by causality, and the
positive-time cutoff may be disregarded when bounding the
convolution at time zero. At past times, derivatives of the cutoff
are supported where \(|y|\geq1/2\) or \(-s\geq1/4\). Move the
derivatives in the divergence term to the kernel. The derivatives
of the kernel on these compact separated sets have finite bounds,
and the fixed spatial and past-time cutoffs have bounded derivatives.
This gives the estimate with \(r=1\). Rescaling the derivative and
the space-time integral proves the displayed powers of \(r\).
See \cite[\S 2.3.3]{Evans} and the compact-support convolution rules
of \cite[Unit U4]{P102}.
\end{proof}


\section{The classical parabolic maximum principle}
\label{pde:H03}

For a bounded nonempty open set \(\Omega\) and \(T>0\), the parabolic
boundary of \(\Omega\times(0,T]\) is
\[
 \partial_p Q_T=(\bar\Omega\times\{0\})
                 \cup(\partial\Omega\times[0,T]).
\]
In particular the top over interior spatial points is a time at
which the comparison can be made.

\begin{theorem}[Classical maximum and comparison]
\label{pde:heat-classical-maximum}
Let \(a^{ij},b^i,c\) be continuous real coefficients on
\(\Omega\times(0,T]\), with \(a\) symmetric uniformly positive
definite and \(c\leq0\). If
\(u\in C^{2,1}(\Omega\times(0,T])\cap C(\bar\Omega\times[0,T])\) and
\[
 Lu:=u_t-\sum_{i,j}a^{ij}u_{ij}-\sum_i b^i u_i-cu\leq0 ,
\]
then
\[
 \max_{\bar\Omega\times[0,T]}u
     \leq \max\bigl(0,\max_{\partial_p Q_T}u\bigr).
\]
When \(c=0\) the outer maximum with zero can be dropped.
If \(Lu\leq Lv\) and \(u\leq v\) on the parabolic boundary, then
\(u\leq v\) everywhere. In particular classical solutions with
the same source, initial values, and lateral values are unique.
\end{theorem}

\begin{proof}
Let \(M=\max(0,\max_{\partial_p Q_T}u)\), or take the boundary
maximum alone when \(c=0\). The function \(w=u-M-\epsilon t\)
satisfies \(Lw\leq -\epsilon+cM+c\epsilon t\leq-\epsilon\).
It is nonpositive on the parabolic boundary. If it were positive
somewhere, compactness would give a positive maximum at a spatial
interior point at positive time. Its spatial gradient there is
zero and its Hessian is negative semidefinite. At an interior
time the time derivative is zero, and at the top its left
derivative is nonnegative. Thus \(Lw\geq0\) at the point, since
\(c\leq0\) and \(w>0\). This is a contradiction. Letting
\(\epsilon\downarrow0\) gives the bound. Apply it to \(u-v\),
and then to the opposite difference for uniqueness.
\end{proof}

The proof uses values of the coefficients only at the candidate
maximum. In particular continuity up to the lateral boundary is
required of \(u\), while the pointwise differential inequality is
imposed only at interior spatial points. See
\cite[\S 7.1.4]{Evans}, \cite[\S 4.4]{RR}, and
\cite[Chapter II, \S 1]{Lieberman}.

\section{Heat balls, strong maximum, and Harnack}
\label{pde:H04}

\begin{definition}[Heat ball]
For \(r>0\) and \(z_0=(x_0,t_0)\), set
\[
 E_r(z_0)=\{(y,s):s<t_0,\ K_{t_0-s}(x_0-y)>r^{-n}\}.
\]
Its top point is \(z_0\). A closure taken in space-time
includes that point.
\end{definition}

\begin{theorem}[Heat-ball mean-value formula]
\label{pde:heat-ball}
Suppose \(u\) is caloric on a space-time neighborhood of
\(\overline{E_r(x_0,t_0)}\), or is classical at the points of
this closure at times \(s<t_0\) and continuous up to the top point.
Then
\[
 u(x_0,t_0)=\frac{1}{4r^n}\iint_{E_r(x_0,t_0)}
       u(y,s)\frac{|x_0-y|^2}{(t_0-s)^2}\,dy\,ds .
\]
In particular the weight with the prefactor shown has total mass one.
\end{theorem}

\begin{proof}
Write \(\tau=t_0-s\) and \(G(y,s)=K_\tau(x_0-y)\).
It satisfies \(G_s+\Delta_y G=0\) for \(s<t_0\).
Fix \(c\geq r^{-n}\) and let \(B_c(s)=\{y:G(y,s)>c\}\).
For \(0<\tau<\tau_c=c^{-2/n}/(4\pi)\), this slice is a ball
of radius
\[
 \rho_c(\tau)=\bigl(4\tau\log((4\pi\tau)^{-n/2}/c)\bigr)^{1/2}.
\]
For \(\tau\geq\tau_c\) the slice is empty.
Since \(G=c\) on the moving sphere, the boundary-motion term in
differentiating \(\int_{B_c(s)}u(G-c)\) vanishes.
Using \(u_s=\Delta_y u\) and the adjoint equation gives
\[
 \frac{d}{ds}\int_{B_c(s)}u(G-c)\,dy
 =\int_{\partial B_c(s)}
       \bigl((G-c)\partial_\nu u-u\,\partial_\nu G\bigr)\,d\sigma
 =\int_{\partial B_c(s)}u\,|\nabla_y G|\,d\sigma .
\]
The identity follows from differentiating the radial integral and
then applying the divergence theorem to
\((G-c)\nabla u-u\nabla G\).
At the earlier limiting time \(t_0-\tau_c\) the spatial volume
tends to zero and the integrand is bounded.
At the top \(s\uparrow t_0\) the radius tends to zero, whereas
\(\rho_c(\tau)/\sqrt\tau\to\infty\). Hence the Gaussian mass of
the slice tends to one. Continuity at \(z_0\) shows that the
\(uG\) term tends to \(u(z_0)\), and the \(cu\) term tends to zero.
Integrating the derivative identity with both time endpoints
truncated and then taking these limits gives
\[
 u(z_0)=\int_{t_0-\tau_c}^{t_0}
           \int_{\partial B_c(s)}u(y,s)|\nabla_y G(y,s)|\,d\sigma\,ds .
\]
For \(u=1\) this gives a positive integral equal to one.
For a bounded \(u\) on the given heat ball it therefore also
justifies absolute convergence of the flux integral.

Multiply by \(c^{-2}\) and integrate over \(c\geq r^{-n}\).
The left side becomes \(r^n u(z_0)\).
On each spatial ray at a fixed \(s\) make the change of variable
\(c=G(x_0+\rho\omega,s)\), with
\(|dc/d\rho|=|\nabla_y G|\) and sphere element
\(\rho^{n-1}d\omega\). This also applies in dimension one, where
the two rays are counted. Changing integration order is justified
by the constant-function computation and boundedness of \(u\).
The right side becomes
\[
  \iint_{E_r(z_0)}u\,\frac{|\nabla_y G|^2}{G^2}\,dy\,ds
   =\frac14\iint_{E_r(z_0)}u\,\frac{|x_0-y|^2}{(t_0-s)^2}\,dy\,ds .
\]
This proves the formula and its normalization.
\end{proof}

\begin{corollary}[Strong maximum principle for the heat equation]
\label{pde:heat-strong-maximum}
Let \(\Omega\) be connected, and let \(u\) be caloric on
\(\Omega\times(0,T]\), continuous there with the one-sided
convention at \(T\). If \(u\leq M\) and
\(u(x_0,t_0)=M\) at \(x_0\in\Omega,\ 0<t_0\leq T\),
then \(u=M\) throughout \(\Omega\times(0,t_0]\).
The minimum statement follows by changing the sign.
\end{corollary}

\begin{proof}
At a maximum point, apply Theorem~\ref{pde:heat-ball} to \(M-u\)
in a sufficiently small heat ball. It is nonnegative and has
weighted integral zero. The weight is positive except on a
space-time line of measure zero, so continuity gives \(u=M\)
on the entire heat ball.

Consider any straight space-time segment inside the cylinder
starting at such a point and ending at an earlier positive time.
For backward time difference \(\tau\) the spatial displacement
along the segment is \(O(\tau)\), whereas
\(\rho_c(\tau)/\tau\to\infty\). Thus the backward part of the
segment sufficiently near the point lies in a small heat ball
and consists of maximum points. Take the infimum of times down
to which the segment consists of maximum points. Continuity
gives equality at that infimum, and another heat ball extends
it further unless the segment's endpoint has been reached.

A connected open subset of \(\mathbb R^n\) is polygonally
connected, since the set of points reachable by polygonal
paths is both relatively open and relatively closed.
Given \(x\in\Omega\) and \(0<t<t_0\), distribute the elapsed
time among the finitely many segments of a polygonal path from
\(x_0\) to \(x\), with strictly decreasing times at the vertices.
The segment argument gives \(u(x,t)=M\). Let \(t\uparrow t_0\)
for the final time slice.
\end{proof}

The heat-ball formula and its backward propagation agree with
the conclusions of \cite[\S\S 2.3.2--2.3.3]{Evans}.
For Harnack below we give a kernel proof for constant coefficients.

\begin{theorem}[Local parabolic Harnack inequality]
\label{pde:heat-harnack}
Let \(u\) be nonnegative and caloric on \(\Omega\times(0,T)\).
If \(V\) is nonempty and connected, with
\(\bar V\) compact and contained in \(\Omega\), and
\(0<t_1<t_2<T\), then
\[
 \sup_{x\in V}u(x,t_1)\leq C\inf_{y\in V}u(y,t_2).
\]
The constant depends on \(n,\Omega,V,t_1,t_2\), independently of \(u\).
\end{theorem}

\begin{proof}
First consider an open cube \(Q\) whose closure lies in
\(\Omega\). By translation let \(Q=(0,\ell)^n\).
The Dirichlet heat kernel on this cube is
\[
 k_Q(\tau;x,y)=\left(\frac{2}{\ell}\right)^n
 \sum_{\mathbf m\in\mathbb N_{\geq1}^n}
   e^{-\pi^2|\mathbf m|^2\tau/\ell^2}
   \prod_{j=1}^n\sin(m_j\pi x_j/\ell)\sin(m_j\pi y_j/\ell).
\]
The series and its differentiated series converge uniformly on
the closed cubes when \(\tau\) stays in a compact positive
interval. The sum is caloric in \((x,\tau)\) and has zero
lateral boundary values. For \(g\in C_c^\infty(Q)\),
\(v(x,\tau)=\int_Q k_Q(\tau;x,y)g(y)\,dy\) has \(v(x,0)=g(x)\).
Indeed the sine coefficients of \(g\) decay faster than any
power by repeated integration by parts. The odd \(2\ell\)-periodic
extension in each coordinate is smooth. Its absolutely
convergent Fourier series equals the extension by completeness
of the periodic Fourier basis, or by convergence of its
product Fejér means. Thus the convergence at zero is uniform.
The periodic Fourier facts used here are given in
\cite[Unit U1]{P102}.

By Theorem~\ref{pde:heat-classical-maximum}, \(g\geq0\)
implies \(v\geq0\). Testing against nonnegative smooth \(g\)
supported near each point and using continuity of the kernel
shows \(k_Q\geq0\).
For fixed \(y\in Q\), \(k_Q(\tau;x,y)\) is strictly positive
at every \(x\in Q\) and every \(\tau>0\).
Otherwise Corollary~\ref{pde:heat-strong-maximum}, applied to the
negative of the kernel on a cylinder starting at any smaller
positive time, would force zero on preceding slices.
But its integral against the positive first sine mode
\(\phi_1(x)=\prod_j\sin(\pi x_j/\ell)\) equals
\(e^{-n\pi^2\tau/\ell^2}\phi_1(y)>0\) by orthogonality.
This proves strict positivity, and compactness then gives a
positive lower bound when \(x,y\) stay in a compact inner cube
and \(\tau\) stays in a compact positive interval.

Fix nested cubes \(Q^*\Subset Q'\Subset Q\), a time increment
\(\delta>0\), and \(r>0\) with \(r^2<t_1\) and
\(B_r(x)\subset Q'\) for all \(x\in Q^*\).
For \(t\geq t_1,\ t+\delta\leq t_2\) and
\(s\in[t-r^2,t]\), comparison on \(Q\times[s,t+\delta]\) gives
\[
 u(y,t+\delta)\geq \int_Q k_Q(t+\delta-s;y,z)u(z,s)\,dz .
\]
To see this without imposing a compatibility condition at the
initial edge, first use initial data
\(g_m=\eta_m u(\cdot,s)\) with \(0\leq\eta_m\leq1\) smooth,
compactly supported in \(Q\), and \(\eta_m\to1\) pointwise.
The solution with those data has zero lateral values and
initial values bounded by \(u\). Apply comparison and pass to
the limit by boundedness of \(u(\cdot,s)\) on \(\bar Q\).
For \(y\in Q^*,z\in Q'\) and
\(\tau=t+\delta-s\in[\delta,\delta+r^2]\), the kernel lower
bound is some \(m>0\). Thus for \(x,y\in Q^*\),
\[
 \int_{t-r^2}^t\int_{B_r(x)}u(z,s)\,dz\,ds
 \leq r^2 m^{-1}u(y,t+\delta).
\]
The center estimate of Theorem~\ref{pde:heat-local} gives
\(u(x,t)\leq C(n)r^{-n}m^{-1}u(y,t+\delta)\).

Finally cover \(V\) by finitely many inner cubes \(Q_j^*\)
with larger cubes \(Q_j\) compactly contained in \(\Omega\).
Such a cover exists by compact containment, and each inner
cube can be chosen to meet \(V\).
Join two indices when their inner cubes intersect in \(V\).
This finite graph is connected, since a separation of the
vertices would separate \(V\) into two relatively open sets.
If the cover has \(N\) cubes, take
\(\delta=(t_2-t_1)/N\).
For \(x,y\in V\), follow a path of cubes between a cube
containing \(x\) and one containing \(y\). Choose intermediate
points in their successive overlaps. At most \(N\) one-cube
inequalities are needed, and any unused time steps can repeat
\(y\) in its final cube. A common bound on the finitely many
one-cube constants gives \(u(x,t_1)\leq C u(y,t_2)\).
Taking the supremum and infimum proves the result.
\end{proof}

This is the constant-coefficient interior Harnack theorem.
The time order in the conclusion is essential to the positivity
argument. For the broader classical formulation see
\cite[\S 7.1.4, Theorem 10]{Evans}.


\section{Weak evolution on a Hilbert triple}
\label{pde:H05}

Let \(V\hookrightarrow H\) be a dense continuous injection of real
Hilbert spaces. Identify \(H\) with a subspace of \(V'\) through
\(\langle h,v\rangle=(h,v)_H\).
The energy space on \((0,T)\) consists of
\(u\in L^2(0,T;V)\) with distributional time derivative
\(u'\in L^2(0,T;V')\). We always use its unique continuous
\(H\)-representative from Section~\ref{pde:C06}.

\begin{theorem}[Coercive and G{\aa}rding parabolic forms]
\label{pde:parabolic-form}
For \(0<t<T\), let \(a_t:V\times V\to\mathbb R\) be bilinear,
write \(\langle A(t)v,w\rangle=a_t(v,w)\), and assume that
\(t\mapsto A(t)v\) is strongly measurable in \(V'\) for each \(v\).
Suppose, with a common exceptional null set independent of \(v,w\),
\[
 |a_t(v,w)|\leq M\|v\|_V\|w\|_V,\qquad
 a_t(v,v)\geq\alpha\|v\|_V^2-c_0\|v\|_H^2,\qquad
 \alpha>0,\quad c_0\geq0 .
\]
For \(u_0\in H\) and \(f\in L^2(0,T;V')\) there exists a unique
\(u\) in the energy space such that
\[
 u'+A(t)u=f\quad\text{in }V'\text{ almost everywhere},
 \qquad u(0)=u_0 .
\]
For all \(0\leq t\leq T\) it satisfies
\[
 \frac12\|u(t)\|_H^2+\int_0^t a_s(u(s),u(s))\,ds
 =\frac12\|u_0\|_H^2+\int_0^t\langle f(s),u(s)\rangle\,ds .
\]
In particular
\[
 \sup_{0\leq t\leq T}\|u(t)\|_H^2+
 \|u\|_{L^2(0,T;V)}^2+\|u'\|_{L^2(0,T;V')}^2
 \leq C\bigl(\|u_0\|_H^2+\|f\|_{L^2(0,T;V')}^2\bigr),
\]
where \(C\) depends only on the form bounds and \(T\).
The autonomous bounded coercive case is obtained with \(c_0=0\)
and \(a_t=a\). In a separable triple, scalar measurability of
\(a_t(v,w)\) implies the stated strong measurability.
\end{theorem}

\begin{proof}
The map \(v(\cdot)\mapsto A(\cdot)v(\cdot)\) is bounded from
\(L^2(0,T;V)\) to \(L^2(0,T;V')\). For simple functions this follows
directly from strong measurability and the uniform bound, and
approximation by simple functions gives the assertion in general.
When \(V\) is separable, identify \(V'\) with \(V\) by its Riesz map.
The finite orthogonal expansions of the representative of \(A(t)v\)
in a countable orthonormal basis have measurable coefficients.
Their pointwise norm limit is strongly measurable.
This justifies the last assertion.

For a finite-dimensional subspace \(F\subset V\), let \(P_F\)
be its orthogonal projection in \(H\) and solve
\[
 (u_F',v)_H+a_t(u_F,v)=\langle f,v\rangle\quad(v\in F),
 \qquad u_F(0)=P_Fu_0 .
\]
In coordinates the mass matrix \(((v_i,v_j)_H)\) is invertible,
since the injection into \(H\) is one-to-one. The resulting system
has bounded measurable coefficient matrix and integrable forcing
on the time interval. Its integral equation has an absolutely
continuous solution: start with the time primitive of the forcing
and the initial vector. Iterating the integral operator with the
bounded matrix gives successive terms bounded in supremum norm
by a fixed bound times \(C^kT^k/k!\). Thus the series converges
uniformly and satisfies the integral equation. The same estimates
on the difference prove uniqueness.

Test the finite-dimensional equation with \(u_F(t)\), use the
G{\aa}rding bound, and apply
\(2ab\leq\alpha a^2+\alpha^{-1}b^2\) to the forcing. This gives
\[
 \frac{d}{dt}\|u_F\|_H^2+\alpha\|u_F\|_V^2
       \leq 2c_0\|u_F\|_H^2+\alpha^{-1}\|f\|_{V'}^2.
\]
Multiplication by \(e^{-2c_0t}\) and integration bound the
\(L^\infty(H)\) term uniformly in \(F\). Integrating the same
inequality gives a uniform \(L^2(V)\) bound.

Direct all finite-dimensional subspaces by inclusion.
The bounded ball in the Hilbert space \(L^2(0,T;V)\) is
weakly compact, so a subnet converges weakly to a function \(u\).
For every fixed \(v\in V\), the subnet eventually contains
\(v\) in its test space. Multiply the tested equation by
a smooth \(\psi:[0,T]\to\mathbb R\) vanishing near \(T\),
and integrate by parts:
\[
 -\int_0^T (u_F,v)_H\psi'\,dt
 +\int_0^T a_t(u_F,v)\psi\,dt
 =\int_0^T\langle f,v\rangle\psi\,dt+(u_0,v)_H\psi(0).
\]
Here \((P_Fu_0,v)_H=(u_0,v)_H\) whenever \(v\in F\).
The terms on the left are continuous linear functionals on
\(L^2(0,T;V)\), the form term by the bounded map proved above.
Hence this identity passes to the weak limit.
Taking \(\psi\) compactly supported inside the time interval shows
\(u'=f-A(t)u\) as a \(V'\)-valued distribution. Indeed the
vector-valued distribution paired with any \(v\in V\) vanishes,
and these pairings separate \(V'\).
Its right side is Bochner \(L^2(V')\).
The Hilbert-triple trace result now gives a continuous
\(H\)-representative. In the identity with \(\psi(0)\ne0\),
integrating the tested limit equation in time with that
representative replaces the boundary term by \((u(0),v)_H\).
Density of \(V\) identifies \(u(0)=u_0\).

The time energy identity from Section~\ref{pde:C06} justifies
pairing the equation with \(u(t)\) and integrating to obtain the
asserted equality. Applying to it the same differential inequality
and integrating factor as above gives the first two norm bounds.
The equation and the uniform bound on \(A(t)\) give the derivative
bound. Apply the norm bounds to the difference of two solutions
with identical data to get uniqueness.
\end{proof}

Using a net in the spatial approximation lets the weak existence
argument work for a general Hilbert triple. The time derivative
bound follows from the limiting equation. No uniform estimate on
the finite-dimensional \(H\)-projections as operators on \(V\)
is needed. See \cite[\S 7.1.2]{Evans},
\cite[\S\S 11.1.1--11.1.2]{RR}, and the parabolic forms in \cite{D16}.

\section{Weak Dirichlet heat flow and weak comparison}
\label{pde:H06}

\begin{theorem}[The energy solution of the Dirichlet heat problem]
\label{pde:heat-dirichlet-weak}
Let \(\Omega\subset\mathbb R^n\) be bounded Lipschitz,
\(u_0\in L^2(\Omega)\), and
\(f\in L^2(0,T;H^{-1}(\Omega))\), with
\(H^{-1}=(H^1_0)'\).
There is a unique
\[
 u\in L^2(0,T;H^1_0(\Omega))\cap C([0,T];L^2(\Omega)),\qquad
 u'\in L^2(0,T;H^{-1}(\Omega)),
\]
satisfying \(u(0)=u_0\) and the weak heat equation
\[
 \langle u'(t),v\rangle+
      \int_\Omega\nabla u(t)\cdot\nabla v
          =\langle f(t),v\rangle
 \qquad(v\in H^1_0(\Omega))
\]
in \(H^{-1}\) almost everywhere. For every \(t\in[0,T]\),
\[
 \frac12\|u(t)\|_2^2+\int_0^t\|\nabla u(s)\|_2^2\,ds
 =\frac12\|u_0\|_2^2+\int_0^t\langle f(s),u(s)\rangle\,ds .
\]
In particular the energy norms are bounded by the corresponding
norms of the data. The lateral condition means zero Sobolev trace
for almost every time, and the initial value is in \(L^2\).
\end{theorem}

\begin{proof}
Take \(H=L^2(\Omega)\), \(V=H^1_0(\Omega)\), and
\(a(v,w)=\int_\Omega\nabla v\cdot\nabla w\).
The Poincar\'e inequality of Section~\ref{pde:E01} makes the
gradient norm equivalent to the \(H^1\) norm on \(V\).
The embedding is dense since it contains \(C_c^\infty(\Omega)\).
Theorem~\ref{pde:parabolic-form} gives the solution and the
displayed equality. More explicitly, the dual norm for the
gradient norm on \(V\) gives
\[
 |\langle f,u\rangle|
     \leq \tfrac12\|\nabla u\|_2^2
                       +\tfrac12\|f\|_{V'}^2.
\]
Using the equivalent \(H^1\) norm changes this bound by the
Poincar\'e constant. The equality and this bound give the
supremum and integrated energy estimates. The derivative estimate
comes from the equation in \(V'\).
The trace characterization of \(H^1_0\) gives the asserted
lateral interpretation.
\end{proof}

We next compare weak subsolutions without requiring their
boundary traces individually to vanish. The positive part of
their difference is the relevant zero-trace test.

\begin{lemma}[Time chain rule for an ordered trace]
\label{pde:heat-positive-chain}
Let \(\Omega\) be bounded Lipschitz and
\(w\in L^2(0,T;H^1(\Omega))\cap C([0,T];L^2(\Omega))\), with
distributional derivative \(w'\in L^2(0,T;H^{-1}(\Omega))\).
If its Sobolev boundary trace is nonpositive almost everywhere
in time and boundary measure, then
\[
 w_+\in L^2(0,T;H^1_0(\Omega))\cap C([0,T];L^2(\Omega)),
 \qquad
 \frac12\|w_+(t)\|_2^2-\frac12\|w_+(s)\|_2^2
     =\int_s^t\langle w'(r),w_+(r)\rangle\,dr
\]
for \(0\leq s\leq t\leq T\). Here \(w_+=\max(w,0)\).
\end{lemma}

\begin{proof}
First recall the spatial chain rule
\(\nabla g_+=\boldsymbol{1}_{\{g>0\}}\nabla g\) for \(g\in H^1\).
For a smooth scalar composition with bounded derivative,
the weak chain rule follows by approximating \(g\) in \(H^1\)
by smooth functions, passing to a pointwise subsequence, and
using dominated convergence for the derivative factors multiplied
by the limiting gradient. To see that \(\nabla g=0\) almost
everywhere on \(\{g=0\}\), take smooth scalar functions with
derivative equal to one for \(|s|\leq\epsilon\), equal to zero
for \(|s|\geq2\epsilon\), and between zero and one everywhere.
Choose their value at zero to be zero. The compositions converge
to zero in \(L^2\), since their absolute values are at most
\(2\epsilon\). Their gradients converge in \(L^2\) to
\(\boldsymbol{1}_{\{g=0\}}\nabla g\).
The distributional derivative of the limit is zero.
Now approximate the positive part by smooth uniformly
Lipschitz scalar functions whose derivatives converge to
\(\boldsymbol{1}_{\{s>0\}}\) at every nonzero \(s\).
The chain rule, the zero-gradient fact, and dominated convergence
give the claimed positive-part formula.

The positive-part map is continuous on \(H^1\), as is seen by
writing the difference of two gradients as the sum of
an indicator times the difference of gradients and
the difference of indicators times the limiting gradient.
For the latter term pass to a pointwise convergent subsequence
and use again the zero-gradient fact. The same proof on the
product measure \(dt\,dx\) proves continuity on \(L^2_tH^1\).
The trace commutes with the positive part. For smooth
\(g\) this follows by the same smooth scalar approximations
and continuity of the trace. For general \(g\), use density
of smooth restrictions on the Lipschitz domain, continuity
of the trace into \(L^2(\partial\Omega)\), and the just proved
continuity in \(H^1\). The trace and density statements are
the Sobolev results of \cite[Unit U5]{P102}.
Thus \(w_+(t)\) has zero trace almost everywhere, and the
trace kernel is \(H^1_0(\Omega)\).
The scalar positive part is \(1\)-Lipschitz, which gives
continuity of \(w_+\) in time with values in \(L^2\).

On a compact subinterval of \((0,T)\) convolve \(w\) in time
with a smooth nonnegative unit-mass kernel of sufficiently
small support. Denote the result by \(w_\epsilon\).
It is smooth in time with values in \(H^1\),
its derivative converges in \(L^2(H^{-1})\) to \(w'\),
and \(w_\epsilon\to w\) in \(L^2(H^1)\) and uniformly in
\(L^2(\Omega)\) on that subinterval.
The trace of \(w_\epsilon\) remains nonpositive, since the
trace commutes with the Bochner integral and the kernel is
nonnegative. Hence \((w_\epsilon)_+\) is a zero-trace test
and converges to \(w_+\) in \(L^2(H^1_0)\).
For the smooth-in-time function use the \(L^2\) chain rule
for \(g\mapsto\frac12\|g_+\|_2^2\), whose derivative is
the \(L^2\) pairing with \(g_+\).
Indeed the scalar function \(F(a)=\frac12(a_+)^2\) has
\(1\)-Lipschitz derivative \(a_+\), so its linearization error
is bounded by \(\frac12|b|^2\) for increment \(b\).
Integrating gives the derivative on \(L^2\).
This gives the displayed identity for \(w_\epsilon\), with
the \(L^2\) pairing agreeing with the \(H^{-1},H^1_0\) pairing.
Pass to the limit by the strong convergences and Cauchy--Schwarz.
The identity extends to \(s=0\) or \(t=T\) by time continuity
and integrability of the dual product.
\end{proof}

\begin{theorem}[Comparison for divergence-form weak heat flow]
\label{pde:heat-weak-comparison}
Let \(A(t,x)\) be a bounded measurable real matrix field,
uniformly elliptic in the sense that
\(\xi\cdot A(t,x)\xi\geq\nu|\xi|^2\) almost everywhere,
\(\nu>0\). Suppose \(u,v\) are in the class of
Lemma~\ref{pde:heat-positive-chain} before imposing a trace sign,
and their difference \(w=u-v\) satisfies
\[
 \langle w'(t),\eta\rangle+
       \int_\Omega A(t,x)\nabla w(t)\cdot\nabla\eta\,dx\leq0
 \qquad(\eta\in H^1_0(\Omega),\ \eta\geq0)
\]
for almost every \(t\). If \(u(0)\leq v(0)\) almost everywhere
and the lateral traces are ordered in the same direction, then
\(u(t)\leq v(t)\) almost everywhere for every \(t\in[0,T]\).
In particular the hypothesis on the difference holds for
solutions with ordered sources in the positive dual cone.
\end{theorem}

\begin{proof}
The trace of \(w\) is nonpositive, so
Lemma~\ref{pde:heat-positive-chain} supplies \(w_+\) as an
admissible test. The inequality is interpreted as a dual-cone
inequality almost everywhere. Equivalently, if it is known
separately for each fixed test, take a countable dense subset
of the nonnegative cone in the separable \(H^1_0(\Omega)\)
and then use continuity to choose a common exceptional set.
For this test,
\[
 \int_\Omega A\nabla w\cdot\nabla w_+
    =\int_\Omega A\nabla w_+\cdot\nabla w_+
    \geq\nu\|\nabla w_+\|_2^2,
\]
by the spatial chain rule. Integrate from zero to \(t\) and
use the time chain rule:
\[
 \frac12\|w_+(t)\|_2^2+
    \nu\int_0^t\|\nabla w_+(s)\|_2^2\,ds
 \leq \frac12\|w_+(0)\|_2^2=0 .
\]
This proves the comparison at every time.
\end{proof}

The spatial locality of the principal part is used when the
positive part replaces the difference in the form. An abstract
coercive form by itself gives the energy estimate but does not
supply this replacement. See \cite[\S\S 7.1.1--7.1.2]{Evans},
\cite[\S\S 10.1--10.2]{Brezis}, and
\cite[Chapter IX, \S\S 3--4]{Lieberman}.


\section{The self-adjoint heat semigroup and its weak solution}
\label{pde:H07}

\begin{theorem}[The operator of a symmetric coercive form]
\label{pde:heat-form-operator}
Let \(V\hookrightarrow H\) be a dense continuous injection of real
Hilbert spaces. Let \(a\) be symmetric, bounded, and coercive on \(V\).
Write \(B=A_H\) for the operator with
\[
 D(B)=\{v\in V:\exists h\in H,\ a(v,w)=(h,w)_H\text{ for all }w\in V\},
 \qquad Bv=h .
\]
Then \(B\) is a densely defined nonnegative self-adjoint operator on
\(H\). Its form domain is \(D(B^{1/2})=V\), with
\(a(v,v)=\|B^{1/2}v\|_H^2\).
If \(V\hookrightarrow H\) is compact, the resolvent is compact.
\end{theorem}

\begin{proof}
Density of \(V\) makes \(h\) unique. The Lax--Milgram theorem
of \cite[Chapter V03]{P106}, applied to \(a(v,w)+(v,w)_H\),
gives for each \(h\in H\) a unique solution \(Kh\in V\)
to
\[
  a(Kh,w)+(Kh,w)_H=(h,w)_H\qquad(w\in V).
\]
In particular \(K=(I+B)^{-1}\), and \(K:H\to V\) is bounded
by coercivity and the continuous injection.
Symmetry and testing two such equations against the two
solutions give \((Kh,g)_H=(h,Kg)_H\).
If \(Kh=0\), then \(h\) is orthogonal to the dense space \(V\),
so \(h=0\). Its range \(D(B)\) is dense in \(H\) because the
orthogonal complement of the range of the bounded self-adjoint
\(K\) is its kernel. The definition gives that \(B\) is
symmetric and nonnegative. For \(w\in D(B^*)\) set
\(z=K(w+B^*w)\). Then
\[
 ((I+B)v,w-z)_H=0\qquad(v\in D(B)),
\]
by symmetry of \(B\) and the definition of \(z\).
Since \(I+B\) is onto, \(w=z\in D(B)\). This proves
self-adjointness.

Apply the self-adjoint spectral theorem, as stated and proved
in \cite{D16}. Over a real Hilbert space one can first
complexify \(B\). Its resolvent construction above complexifies,
and uniqueness of the spectral measure on the real axis shows
that its projections commute with conjugation.
Thus the projections and real functions of \(B\) restrict
to the real space.

On \(D(B)\), \(a(v,v)=(Bv,v)_H=\|B^{1/2}v\|_H^2\).
The form defines an inner product on \(V\) with norm
equivalent to the given one. If a vector \(w\in V\) is
\(a\)-orthogonal to \(D(B)\), solve
\(a(v,\cdot)=(h,\cdot)_H\) by Lax--Milgram for each \(h\in H\).
This yields \(v\in D(B)\), \(Bv=h\), and hence \((h,w)_H=0\).
It follows that \(D(B)\) is dense in \(V\).
On the other hand for \(v\in D(B^{1/2})\) its spectral
cutoffs to \([0,N]\) belong to \(D(B)\) and converge in the
graph norm of \(B^{1/2}\), by square integrability of the
spectral weight. The two norm completions of \(D(B)\)
inside \(H\) therefore agree, and the quadratic equality
extends. Finally \(K\) is compact when the injection
\(V\hookrightarrow H\) is compact.
\end{proof}

In the Dirichlet setting take \(V=H^1_0(\Omega)\),
\(H=L^2(\Omega)\), and the symmetric elliptic form from
Section~\ref{pde:E10}. Poincar\'e gives coercivity, and
Rellich compactness gives the compact injection used above.

\begin{theorem}[Heat semigroup and spectral smoothing]
\label{pde:heat-dirichlet-semigroup}
Let \(B\) be as in Theorem~\ref{pde:heat-form-operator}.
Then \(S_B(t)=e^{-tB}\) is a strongly continuous
self-adjoint contraction semigroup on \(H\) with generator
\(-B\). For every integer \(m\geq1\) and \(t>0\),
\[
 S_B(t)H\subset D(B^m),\qquad
 \|B^m S_B(t)\|_{\mathcal L(H)}\leq (m/(et))^m .
\]
When the resolvent is compact, an orthonormal eigenbasis
\((\phi_j)\) of \(B\), with eigenvalues \(\lambda_j\),
gives the norm-convergent expansion
\[
 S_B(t)g=\sum_j e^{-t\lambda_j}(g,\phi_j)_H\,\phi_j
 \qquad (g\in H).
\]
The sum has finitely many terms if \(H\) is finite-dimensional.
\end{theorem}

\begin{proof}
Define \(S_B(t)\) by the bounded spectral function
\(\lambda\mapsto e^{-t\lambda}\) on \([0,\infty)\).
Multiplication gives the semigroup law. The function is real
and bounded by one, so the operators are self-adjoint
contractions. Dominated convergence in the scalar spectral
measure of each vector gives strong continuity.
For \(g\in D(B)\) the difference quotient at zero converges
to \(-Bg\) by
\(|(e^{-t\lambda}-1)/t|\leq\lambda\) and dominated convergence.
Conversely if the difference quotient has a norm limit,
Fatou's lemma gives integrability of \(\lambda^2\)
against that vector's spectral measure. Thus its generator
domain is exactly \(D(B)\). The bound for \(B^m S_B(t)\)
follows by maximizing \(\lambda^m e^{-t\lambda}\) over
\(\lambda\geq0\).
For the expansion, the compact self-adjoint \(K=(I+B)^{-1}\)
has an orthonormal eigenbasis since it is injective.
Each eigenvector with eigenvalue \(\kappa>0\) lies in the
range of \(K\) and has eigenvalue \(\kappa^{-1}-1\) for \(B\).
The spectral formula on those basis vectors and Parseval's
identity give the expansion.
\end{proof}

\begin{proposition}[Agreement of mild and energy solutions]
\label{pde:heat-mild-energy}
Let \(B\) and \(a\) be as above, \(u_0\in H\), and
\(f\in L^1(0,T;H)\). The mild solution in \(H\) is
\[
 U(t)=S_B(t)u_0+\int_0^t S_B(t-s)f(s)\,ds .
\]
Whenever there is an energy solution of
\(u'+Au=f\) in \(V'\) with \(u(0)=u_0\), it equals \(U\) at
every time. In particular this applies to \(f\in L^2(0,T;H)\).
For arbitrary \(f\in L^2(0,T;V')\), the energy solution is
given by Theorem~\ref{pde:parabolic-form} in its own spaces.
\end{proposition}

\begin{proof}
The formula in \(H\) is the variation-of-constants theorem of
Section~\ref{pde:C07}, using generator \(-B\).
Let \(u\) be an energy solution. For \(N\geq1\) let
\(P_N\) be the spectral projection of \(B\) on \([0,N]\).
It maps \(H\) boundedly into \(V\), since coercivity and the
form-domain identity give
\(\alpha\|P_N h\|_V^2\leq N\|h\|_H^2\).
It thus acts from \(V'\) to \(H\) by
\[
 (P_N\ell,h)_H=\langle \ell,P_Nh\rangle \qquad(h\in H).
\]
This agrees with the original projection on \(H\).
The derivative of \(P_Nu\) in \(H\)-valued distributions is
\(P_Nu'\in L^2(0,T;H)\).
For \(h\in H\), \(P_Nh\in D(B)\), so symmetry gives
\[
 a(u,P_Nh)=a(P_Nh,u)=(BP_Nh,u)_H=(h,BP_Nu)_H .
\]
The boundedness of \(B\) on \(P_NH\) and the projected equation
give \((P_Nu)'+BP_Nu=P_N f\) there.
Solving this bounded-operator equation by multiplying with
\(e^{tB}\) on \(P_NH\) gives the projected formula for every
time, using its initial value \(P_Nu_0\).
Let \(N\to\infty\). The projections converge strongly to the
identity on \(H\), and dominated convergence using
\(\|f\|_{L^1(H)}\) passes the projected integral to the
integral displayed in the proposition. This proves \(u=U\).
Existence for \(f\in L^2(H)\) follows from its continuous
inclusion in \(L^2(V')\) and the energy theorem.
\end{proof}

The two source spaces serve different purposes.
The displayed Bochner convolution uses \(H\)-valued data.
For data given only in \(V'\) the pairing and estimate in
Theorem~\ref{pde:parabolic-form} remain meaningful.
For the classical bounded-domain model see
\cite[Chapter 6, \S 1]{TaylorI}, \cite[\S 10.1]{Brezis}, and
\cite[\S 7.4]{Evans}. The underlying spectral and closed-form
results can also be compared with \cite{D16}.

\section{Improved energy regularity}
\label{pde:H08}

\begin{theorem}[An additional energy derivative]
\label{pde:heat-improved}
Under the symmetric time-independent form assumptions of
Theorem~\ref{pde:heat-form-operator}, suppose \(u_0\in V\)
and \(f\in L^2(0,T;H)\). The energy solution satisfies
\[
 u'\in L^2(0,T;H),\qquad
 Bu\in L^2(0,T;H),\qquad
 u\in L^\infty(0,T;V).
\]
Here \(u(t)\in D(B)\) for almost every \(t\), and
\(u'+Bu=f\) holds in \(H\). Quantitatively,
\[
 \|u'\|_{L^2(H)}^2+\|Bu\|_{L^2(H)}^2+
      \sup_{0\leq t\leq T}\|u(t)\|_V^2
 \leq C\bigl(\|u_0\|_V^2+\|f\|_{L^2(H)}^2\bigr),
\]
with \(C\) depending on the form bounds.
If \(B\) is the Dirichlet operator of Section~\ref{pde:E06}
on a bounded \(C^{1,1}\) domain with symmetric uniformly elliptic
Lipschitz coefficients, then additionally
\(u\in L^2(0,T;H^2(\Omega))\), with the bound from elliptic
regularity applied to \(Bu\).
\end{theorem}

\begin{proof}
Use the spectral projections \(P_N\) from the preceding proof
and write \(u_N=P_Nu\).
The projected equation holds in \(H\), with \(B\) bounded on
\(P_NH\), and \(u_N'\in L^2(H)\).
Take the inner product with \(u_N'\). Because \(B\) is
self-adjoint on that spectral subspace,
\[
 \|u_N'\|_H^2+\tfrac12\frac{d}{dt}\|B^{1/2}u_N\|_H^2
        =(P_Nf,u_N')_H
        \leq \tfrac12\|P_N f\|_H^2+\tfrac12\|u_N'\|_H^2.
\]
Integration yields for every \(t\)
\[
 \int_0^t\|u_N'\|_H^2\,ds+\|B^{1/2}u_N(t)\|_H^2
 \leq \|B^{1/2}u_0\|_H^2+\int_0^t\|f\|_H^2\,ds .
\]
The equation bounds \(\|Bu_N\|_{L^2(H)}\) by the same data.
By monotone convergence of the scalar spectral integrals
with weights \(\lambda\) and \(\lambda^2\), the uniform
bounds imply \(u(t)\in V\) for all \(t\),
\(u(t)\in D(B)\) for almost every \(t\), and the stated
integrability bound for \(Bu\).
Spectral convergence gives strong convergence
\(Bu_N\to Bu\) in \(L^2(H)\), and the limit is strongly
measurable there. Likewise \(u_N\to u\) in \(L^2(H)\)
and \(P_Nf\to f\) in \(L^2(H)\).
Pass in the projected distributional derivative equation to
get \(u'=f-Bu\in L^2(H)\). Use the form-domain identity and
coercivity to obtain the \(V\) bound.

For the last assertion, for almost every \(t\) the weak
Dirichlet elliptic equation for \(u(t)\) has \(L^2\)
right-hand side \(Bu(t)\).
The \(H^2\) estimate of Section~\ref{pde:E06} bounds
\(\|u(t)\|_{H^2}\) by \(C(\|Bu(t)\|_2+\|u(t)\|_{H^1})\).
The resulting map from the graph domain to \(H^2\) is
continuous, so \(u\) is strongly measurable in \(H^2\)
as a graph-domain-valued limit of spectral truncations.
Integrate the squared estimate in time.
\end{proof}

This last implication is a spatial elliptic regularity step.
The energy estimate alone controls the operator domain.
Compare \cite[\S 7.1.3, Theorem 5]{Evans} and
\cite[\S 10.1]{Brezis}.
\chapter{Characteristics, Waves, and Hyperbolic Systems}
\label{chap:pde-hyperbolic}

In this chapter \(n\geq 1\). All time derivatives at endpoints are
one-sided. The wave speed is denoted by \(c>0\). For complex-valued
fields the energy identities use the real part of the Hilbert pairing.

\section{Characteristics and transport}\label{pde:W01}

\begin{theorem}[Transport by a nonautonomous flow]\label{thm:pde-transport-flow}
Let \(b\in C^1([0,T]\times\mathbb R^n;\mathbb R^n)\) be bounded with
bounded spatial derivative. Denote the flow by \(X(t,s,y)\), where
\(X(s,s,y)=y\). For each \(t,s\), \(X(t,s,\cdot)\) is a \(C^1\)
diffeomorphism with inverse \(X(s,t,\cdot)\). For \(g\in C^1\),
the unique classical solution of
\[
 u_t+b\cdot\nabla u=0,\qquad u(0)=g
\]
is \(u(t,x)=g(X(0,t,x))\). If \(q,f\) are continuous with
continuous first spatial derivatives, the equation
\(u_t+b\cdot\nabla u+qu=f\) has the solution
\begin{align*}
 u(t,x)={}&
 e^{-\int_0^t q(r,X(r,t,x))\,dr}g(X(0,t,x))\\
 &+\int_0^t e^{-\int_s^t q(r,X(r,t,x))\,dr}
                     f(s,X(s,t,x))\,ds .
\end{align*}
For a local \(C^1\) vector field, the formulas apply on the actual
flow domain where the entire characteristic reaches the initial data.
\end{theorem}

\begin{proof}
The bounded derivative gives uniform spatial Lipschitz continuity.
On any finite time interval the estimate
\(\lvert X(t,s,y)-y\rvert\leq \|b\|_\infty |t-s|\) prevents escape
before the ends of the time slab. Local existence, uniqueness,
continuation and \(C^1\) dependence on \(t,s,y\) then give the stated
flow by the nonautonomous ODE results in
\cite[Chapters 3, 4 and 6]{D13}. Uniqueness gives
\(X(t,r,X(r,s,y))=X(t,s,y)\) on the common flow domain. In particular
\(X(s,t,\cdot)\) is the inverse. Along \(X(t,0,y)\) the equation is
the scalar ODE \(v'(t)+q(t,X(t,0,y))v(t)=f(t,X(t,0,y))\).
The integrating factor gives exactly the displayed formula.
The spatial derivatives of its integrands are continuous and
locally bounded on compact time intervals, so the formula is
\(C^1\) in space. In flow coordinates the time derivative follows
from the scalar ODE. The \(C^1\) change of coordinates then gives
the equation and a continuous time derivative in the original
coordinates. Uniqueness also follows by restricting any solution
to the same characteristics. The argument is local whenever the
flow and data are local.
\end{proof}

\begin{proposition}[Continuity equation and transverse data]\label{prop:pde-transport-jacobian}
Under the flow hypotheses, set \(J(t,y)=\det D_yX(t,0,y)\).
Then
\[
 J(t,y)=\exp\left(\int_0^t
             \operatorname{div}b(r,X(r,0,y))\,dr\right).
\]
A locally integrable initial density \(g\) transported by the
flow has density
\[
 \rho(t,X(t,0,y))=\frac{g(y)}{J(t,y)} .
\]
This density satisfies
\(\rho_t+\operatorname{div}(b\rho)=0\) with initial density \(g\)
in the distributional sense. Locally, prescribed \(C^1\) data on
a \(C^1\) hypersurface transverse to the space--time vector
\((1,b)\) determine a unique \(C^1\) transport solution near
the initial point, also with the zero-order term and source above.
\end{proposition}

\begin{proof}
The variational equation for \(D_yX\) is \(Y'=D_xb(t,X)Y\),
\(Y(0)=I\). Liouville's determinant formula in
\cite[Chapter 7]{D13} gives the claimed positive Jacobian. Thus
change of variables gives, for each compactly supported test,
\[
 \int\rho(t,x)\varphi(t,x)\,dx
     =\int g(y)\varphi(t,X(t,0,y))\,dy .
\]
Integrating the derivative of the right-hand side in time gives
\[
 \int_0^T\!\!\int \rho(\varphi_t+b\cdot\nabla\varphi)\,dx\,dt
   =-\int g(y)\varphi(0,y)\,dy
\]
when the test vanishes near \(T\). The bounded flow displacement
puts the pulled-back supports in a common compact set, justifying
this calculation for locally integrable \(g\). This weak statement
uses only the stated \(C^1\) regularity of \(b\).

For the last assertion parametrize the initial hypersurface locally
by \(y\mapsto \sigma(y)\). The elapsed-time flow of \((1,b)\)
from \(\sigma(y)\) defines a \(C^1\) map of \((\tau,y)\).
At \(\tau=0\) its differential has the tangent vectors of
the hypersurface and the transverse vector as columns. These
form a basis. The inverse function theorem gives flow coordinates
near the point. Solving the scalar ODE in the elapsed-time variable
gives the result there, and the same ODE proves uniqueness.
\end{proof}

See \cite[Sections 2.1 and 3.2]{Evans},
\cite[Chapter 2]{RR}, and \cite[Chapter 1, Section 5]{John}.

\section{Classical wave formulas in every dimension}\label{pde:W02}

For \(n\geq 2\) write \(M_n h(x,r)\) for the normalized spherical
mean of \(h\) over \(x+rS^{n-1}\), using unit-sphere probability
measure. It extends evenly in \(r\) with \(M_n h(x,0)=h(x)\).
Put \(D_r=r^{-1}\partial_r\) for \(r>0\).

\begin{theorem}[Spherical means and descent]\label{thm:pde-classical-wave}
For smooth data \(g,h\) the equation \(u_{tt}-c^2\Delta u=0\)
with \(u(0)=g,\ u_t(0)=h\) has the classical solution
\[
 u(t)=\partial_t S_n(t)g+S_n(t)h .
\]
The velocity operators \(S_n\) have the formulas
\begin{align*}
 S_1(t)h(x)&=\frac1{2c}\int_{x-ct}^{x+ct}h(y)\,dy,\\
 S_n(t)h(x)&=\left.
     \frac{D_r^{(n-3)/2}\{r^{n-2}M_n h(x,r)\}}
          {c(n-2)!!}\right|_{r=ct}
                 &&(n\geq3\text{ odd}),\\
 S_n(t)h(x)&=\left.
     \frac{D_r^{(n-2)/2}\{r^{n-1}\widetilde M_n h(x,r)\}}
          {c(n-1)!!}\right|_{r=ct}
                 &&(n\geq2\text{ even}),
\end{align*}
where \(\widetilde M_n\) is the mean from dimension \(n+1\) of
the lift independent of the last coordinate. Explicitly,
\[
 \widetilde M_n h(x,r)=\frac{2}{|S^n|}
       \int_{|z|<1}\frac{h(x+rz)}{\sqrt{1-|z|^2}}\,dz .
\]
In particular \(S_3(t)h=t M_3h(x,ct)\), and
\[
 S_2(t)h(x)=\frac{1}{2\pi c}
       \int_{|y-x|<ct}\frac{h(y)}{\sqrt{c^2t^2-|y-x|^2}}\,dy .
\]
For smooth forcing \(f(t,x)\), add \(\int_0^t S_n(t-s)f(s)\,ds\).
In odd dimensions \(n\geq3\) the homogeneous value at positive time
depends only on the data and their relevant derivatives on the
sphere \(|y-x|=ct\).
\end{theorem}

\begin{proof}
In one dimension, differentiating the integral gives
\(\partial_tS_1(t)g=(g(x+ct)+g(x-ct))/2\).
A second differentiation, or the factorization into the two
characteristic derivatives \(\partial_t\pm c\partial_x\), verifies
the equation. The initial velocity and displacement follow
by differentiating at \(t=0\).

For \(n\geq2\) the divergence theorem on the ball gives
\[
 r^{n-1}\partial_r M_n h(x,r)
   =|S^{n-1}|^{-1}\int_{B(x,r)}\Delta h(y)\,dy .
\]
Differentiate in \(r\) to obtain
\(\Delta_x M_n h=\partial_r^2 M_n h+(n-1)r^{-1}\partial_rM_n h\).
In odd dimension write \(n=2m+3\), \(F=M_n h\), and \(G=r^{2m+1}F\).
Direct differentiation gives
\[
 \Delta_xG=(\partial_r^2-2m r^{-1}\partial_r)G,\qquad
 D_r(\partial_r^2-2q r^{-1}\partial_r)
    =(\partial_r^2-2(q-1)r^{-1}\partial_r)D_r .
\]
For the first equality the term proportional to \(F\) cancels.
For the second, applying both sides to a smooth \(v\) gives
\(v'''/r-2qv''/r^2+2qv'/r^3\). Iterating shows
\(\Delta_xD_r^mG=\partial_r^2D_r^mG\). Substituting \(r=ct\)
therefore proves the wave equation for \(S_n(t)h\).

The apparent singularity in this formula is removable at zero.
If \(F\) is smooth and even, \(D_rF\) extends smoothly and evenly
since \(F'(r)/r=\int_0^1 F''(\theta r)\,d\theta\).
The identity
\(D_r(r^k F)=kr^{k-2}F+r^k D_rF\) then shows that \(D_r^mG\)
is smooth and odd. Its linear term is \((2m+1)!!\,rh(x)\),
with remainder \(O(r^3)\) locally uniformly with all needed
\(x\)-derivatives. Thus \(S_n(0)h=0\), \(S_n'(0)h=h\),
and \(S_n''(0)h=0\). These three identities also hold for \(n=1\).
They give the two initial data for
\(\partial_t S_n(t)g+S_n(t)h\), which solves the wave equation
by commutation with the time derivative.

For even \(n\) lift \(h\) to \(\mathbb R^{n+1}\) as a function
independent of the last coordinate. The odd-dimensional formula is
also independent of that coordinate, so its equation reduces to
the \(n\)-dimensional equation and its initial data are retained.
To calculate the mean, each hemisphere of the unit sphere is
the graph \(z\mapsto (z,\pm\sqrt{1-|z|^2})\) over the unit ball.
Its area density is \((1-|z|^2)^{-1/2}\), and the equator has
surface measure zero. This gives the displayed \(\widetilde M_n\)
and, after the change \(y=x+ct\,z\), the two-dimensional constant.

All differentiations in the source integral are justified on compact
space--time sets by the formulas with a fixed unit sphere or unit
ball and the smooth zero-time extension just verified. In two time
derivatives the only nonzero endpoint term is
\(S_n'(0)f(t)=f(t)\), so the integral solves the forced equation
with zero initial data. Uniqueness for smooth solutions follows
from the local cone estimate proved in Section~\ref{pde:W04}.
For odd \(n\geq3\), the radial and time derivatives of the means
involve only finitely many derivatives at points on the indicated
sphere, proving the Huygens assertion.
\end{proof}

The formulas and their normalization are compared with
\cite[Sections 2.4.1--2.4.2]{Evans},
\cite[Chapter 5, Section 1]{John}, and
\cite[Sections 2.1 and 9.2]{Strauss}.

\section{Whole-space finite-energy waves}\label{pde:W03}

\begin{theorem}[Forced wave Cauchy problem in the energy space]\label{thm:pde-wave-energy}
Let \(g\in H^1(\mathbb R^n)\), \(h\in L^2(\mathbb R^n)\) and
\(f\in L^1_{\mathrm{loc}}([0,\infty);L^2(\mathbb R^n))\).
There is a unique distributional solution of
\(u_{tt}-c^2\Delta u=f\) in
\(C([0,\infty);H^1)\cap C^1([0,\infty);L^2)\)
with \(u(0)=g,\ u_t(0)=h\). Put
\[
 G(t)=\bigl(\|u_t(t)\|_2^2+c^2\|\nabla u(t)\|_2^2\bigr)^{1/2},
 \qquad E(t)=G(t)^2/2 .
\]
For \(0\leq r\leq t\),
\begin{align*}
 E(t)-E(r)&=\int_r^t \operatorname{Re}(f(s),u_t(s))_{L^2}\,ds,\\
 G(t)&\leq G(r)+\int_r^t\|f(s)\|_2\,ds,\\
 \|u(t)\|_2&\leq \|u(r)\|_2+\int_r^t G(s)\,ds .
\end{align*}
In particular the homogeneous energy is conserved, and the
solution depends continuously on all three data on every bounded
time interval.
\end{theorem}

\begin{proof}
Use the transform with phase \(e^{-2\pi i x\cdot\xi}\) from
\cite[Unit U3]{P102} and set \(\omega=2\pi c|\xi|\).
Let \(C(t)\) and \(S(t)\) be the multipliers
\(\cos(t\omega)\) and \(\sin(t\omega)/\omega\), with the latter
equal to \(t\) at zero. Plancherel and the multiplier bounds
\[
 |S(t)|\leq |t|,\qquad |\omega S(t)|\leq 1,\qquad
 |\cos(t\omega)|\leq 1
\]
give strongly continuous operator families \(C(t):H^1\to H^1\)
and \(S(t):L^2\to H^1\), as well as their time derivatives from
the same source spaces to \(L^2\).
The strong continuity and derivative formulas follow from
dominated convergence for each fixed input, using the bound on
\(\omega\widehat g\) for \(C'(t)g\). The Bochner formula
\[
 u(t)=C(t)g+S(t)h+\int_0^t S(t-s)f(s)\,ds
\]
is continuous in \(H^1\). Its derivative in \(L^2\) is
\(C'(t)g+C(t)h+\int_0^t C(t-s)f(s)\,ds\), since \(S(0)=0\).
One can see this by differentiating the integral first for continuous
\(f\), and using the bounds to pass to \(L^1(L^2)\) uniformly
for both the function and derivative. The two initial values follow.

To verify the equation and uniqueness, let \(Q_N\) be Fourier
projection onto \(|\xi|\leq N\). For any competing energy
solution the distributional equation also holds as a time
distribution with values in \(H^{-1}\), by density of smooth
compactly supported spatial tests in \(H^1\). The map
\(Q_N:H^{-1}\to L^2\) is bounded, and on its range \(\Delta\)
is bounded in \(L^2\). The truncated formula solves the bounded
operator equation with the corresponding data, by differentiation
in distributions (or first with continuous forcing and then in
\(L^1\)). Passing to the limit gives the PDE. For any proposed
energy solution, its projection has distributional second derivative
\(c^2\Delta Q_Nu+Q_N f\in L^1_{\mathrm{loc}}(L^2)\). The first
derivative is therefore locally absolutely continuous and it
solves the same bounded operator equation with the same data.
The difference with the projected formula has zero data, and
the integral form of this bounded first-order system for the pair
\((Q_Nu,Q_Nu_t)\) gives uniqueness by Gronwall. Let \(N\) tend to
infinity.

For the projected solutions differentiation of the energy gives
the identity by Plancherel and cancellation of the Laplacian term.
At each fixed time the energies converge. The source pairings
converge in \(L^1\) on finite intervals because \(u_t\) is bounded
there in \(L^2\) and \(\|Q_Nf\|_2\leq\|f\|_2\).
This proves the energy identity and absolute continuity of \(E\).
For \(\delta>0\) the derivative of \(\sqrt{2E+\delta}\) is
bounded above by \(\|f\|_2\) almost everywhere. Integration and
\(\delta\downarrow0\) prove the estimate on \(G\).
Finally integrate \(u_t\) in \(L^2\) to control the displacement.
Applying these estimates to the difference of two solutions
proves continuous dependence.
\end{proof}

See \cite[Sections 2.4 and 4.3.1]{Evans},
\cite[Chapter 6, Section 2]{TaylorI}, and the homogeneous abstract
spectral wave formula in \cite{D16}.

\section{Cone energy, domain of dependence, and finite speed}\label{pde:W04}

\begin{theorem}[Localized wave energy]\label{thm:pde-wave-cone}
Let \(u\) solve \(u_{tt}-c^2\Delta u=f\) as in
Theorem~\ref{thm:pde-wave-energy}. Fix \(x_*\in\mathbb R^n\),
\(R>0\), and write \(B_s=B(x_*,R-cs)\) for \(0\leq s<R/c\).
For
\[
 E_{B_s}(s)=\frac12\int_{B_s}(|u_t(s)|^2+c^2|\nabla u(s)|^2)\,dx
\]
one has
\[
 \sqrt{2E_{B_t}(t)}\leq \sqrt{2E_{B_0}(0)}
       +\int_0^t\|\mathbf 1_{B_s}f(s)\|_2\,ds
                 \qquad (0\leq t<R/c).
\]
If the two initial data vanish almost everywhere on \(B_0\)
and \(f\) vanishes almost everywhere in the backward cone, then
\(u(t)\) vanishes almost everywhere on every \(B_t\).
The same statement for differences gives the domain of dependence.
For \(f=0\) and both data supported in \(\overline B(x_0,R_0)\),
the support of \(u(t)\) lies in \(\overline B(x_0,R_0+ct)\).
\end{theorem}

\begin{proof}
For smooth solutions the density
\(e=(|u_t|^2+c^2|\nabla u|^2)/2\) satisfies
\[
 e_t=\operatorname{div}
          (c^2\operatorname{Re}(\overline{u_t}\nabla u))
        +\operatorname{Re}(\overline{u_t}f).
\]
The moving boundary has outward normal speed \(-c\).
Consequently the derivative of its energy is the interior
forcing term plus the integral of
\(c^2\operatorname{Re}(\overline{u_t}\partial_\nu u)-ce\).
The boundary expression is nonpositive by
\(2c|u_t||\partial_\nu u|\leq |u_t|^2+c^2|\nabla u|^2\).
The forcing term is bounded by
\(\|\mathbf 1_{B_s}f(s)\|_2\sqrt{2E_{B_s}(s)}\).
Regularizing the square root as in the preceding proof gives
the inequality.

On any finite time interval, approximate the data in
\(H^1\times L^2\) and the forcing in \(L^1(L^2)\) by smooth
compactly supported data and sources. Their solutions are
smooth by applying the Fourier bounds to all spatial derivatives
and using the equation for further time derivatives. Spatial
mollifiers may be taken compactly
supported, so compact spatial supports are enlarged by only
the mollifier radius. The energy bounds give uniform convergence
of the solutions and time derivatives in their energy spaces,
so the restricted energies at the initial and terminal times
converge. The restricted forcing norms converge in \(L^1\)
by the \(L^1(L^2)\) approximation. This proves the inequality
for energy solutions without imposing any trace on a moving
sphere for their weak gradients.

With zero data and source as in the statement, the inequality
gives \(u_t(s)=0\) on \(B_s\) for each such \(s\). For fixed \(t\)
the ball \(B_t\) is contained in \(B_s\) for all \(s\leq t\).
Integrating the time derivative in \(L^2(B_t)\) gives
\(u(t)|_{B_t}=u(0)|_{B_t}=0\). This last step uses the initial
displacement itself, beyond its gradient. If \(x\) lies outside
\(\overline B(x_0,R_0+ct)\), choose \(R>ct\) such that
\(B(x,R)\) is disjoint from \(\overline B(x_0,R_0)\).
The result just proved gives zero on a neighborhood of \(x\)
at time \(t\). A countable base of balls gives the essential
support assertion. For smooth solutions the same calculation
is local and applies without global energy assumptions.
\end{proof}

This is the energy argument of \cite[Section 2.4.3, Theorem 6]{Evans}
and \cite[Sections 2.2 and 9.1]{Strauss}.

\section{The Dirichlet wave in the form domain}\label{pde:W05}

\begin{theorem}[Finite-energy Dirichlet wave]\label{thm:pde-dirichlet-wave}
Let \(\Omega\subset\mathbb R^n\) be nonempty, bounded and Lipschitz.
Use \(H=L^2(\Omega)\), \(V=H^1_0(\Omega)\), and a bounded symmetric
coercive form \(a\) on \(V\), Hermitian over \(\mathbb C\).
For \(g\in V\), \(h\in H\), and \(f\in L^1(0,T;H)\), there is a
unique
\[
 u\in C([0,T];V)\cap C^1([0,T];H)
\]
with the prescribed initial values and
\(\langle u_{tt},v\rangle+a(u,v)=(f,v)_H\) for every \(v\in V\)
as a time distribution. In particular
\(u_{tt}\in L^1(0,T;V')\).
For \(E_a(t)=\tfrac12(\|u_t(t)\|_H^2+a(u(t),u(t)))\) the energy
identity is
\[
 E_a(t)-E_a(r)
   =\int_r^t\operatorname{Re}(f(s),u_t(s))_H\,ds .
\]
For the Dirichlet elliptic realization of Section~\ref{pde:E10},
this is the weak wave equation with homogeneous Dirichlet boundary
data. In the formula \(u=\cos(t\sqrt{A_H})g+
\sin(t\sqrt{A_H})h/\sqrt{A_H}\) plus the corresponding source
convolution, the displacement data belong to the form domain \(V\).
\end{theorem}

\begin{proof}
The coercive form theorem gives the bounded solution operator
\(R:H\to V\) characterized by \(a(Rw,v)=(w,v)_H\).
The inclusion \(V\to H\) is compact by Rellich, so on \(H\) the
operator \(R\) is compact. Symmetry gives self-adjointness, and
\((Rw,w)_H=a(Rw,Rw)\geq0\). If \(Rw=0\), then \(w\) is orthogonal
to the dense subspace \(V\), hence \(w=0\). As in
Section~\ref{pde:E10}, the compact spectral theorem gives an
orthonormal basis \(\phi_k\) of \(H\) with \(R\phi_k=\lambda_k^{-1}\phi_k\)
and \(\lambda_k>0\). This proof applies to the abstract form here
as well as to the Dirichlet form there. In particular
\(a(\phi_k,v)=\lambda_k(\phi_k,v)_H\).
The vectors \(\phi_k/\sqrt{\lambda_k}\) form an orthonormal
basis of \(V\) for its form norm: they are orthonormal and a vector
\(a\)-orthogonal to all of them has all its \(H\)-coefficients
zero. Thus \(a(v,v)=\sum_k\lambda_k|(v,\phi_k)_H|^2\).
The realization is \(A_H=R^{-1}\) on \(\operatorname{Ran}R\).
In coefficients this domain requires
\(\sum_k\lambda_k^2|(v,\phi_k)_H|^2<\infty\), while the form
norm needs only the first power of \(\lambda_k\).

Project the data onto the first \(N\) basis vectors. In each mode
solve \(u_k''+\lambda_k u_k=f_k\) by the sine and cosine
formula, with the source convolution in the sine operator.
The finite sums give the stated energy identity by differentiation
of absolutely continuous coefficients. For two such sums, the
energy norm of the difference at time \(t\) is bounded by its
initial energy norm plus the \(L^1(H)\) norm of the difference of
the projected sources. This follows by regularizing the square
root of the energy as in Theorem~\ref{thm:pde-wave-energy}.
The initial tails tend to zero by the two Parseval identities.
The source tails tend to zero by pointwise strong convergence of
the projections and dominated convergence in \(L^1(H)\). The
finite sums therefore converge in \(C(V)\) and their derivatives
in \(C(H)\), giving a \(C^1(H)\) limit with the correct initial
values. The equation passes in time distributions tested against
\(V\), and \(u_{tt}=f-a(u,\cdot)\) gives the asserted integrability
in \(V'\). The limit has zero trace in space because its values
lie in \(V\). The energy identity passes by the same uniform
convergence and \(L^1\) control. Finally, testing the equation
of any competing solution with \(\phi_k\) gives the scalar
equation and both initial values, hence the same coefficients
for every \(k\). Completeness gives uniqueness.
\end{proof}

See \cite[Section 7.2]{Evans}, \cite[Chapter 6, Section 1]{TaylorI},
\cite[Section 10.3]{Brezis}, and \cite[Section 12.3.2]{RR}.
The abstract coercive form and spectral constructions are in
\cite{P106,D16}.

\section{The first-order symmetric energy estimate}\label{pde:W06}

\begin{theorem}[An \(L^2\) estimate with linear growth in the principal part]
\label{thm:pde-symmetric-energy}
Suppose \(A^j=A^{j*}\) are \(C^1\) matrix fields on
\([0,T]\times\mathbb R^n\), with bounded first derivatives, and
\(B\) is bounded. For a smooth solution of
\[
 u_t+\sum_{j=1}^n A^j(t,x)\partial_j u+B(t,x)u=f(t,x)
\]
with \(u(0)\in L^2\) and \(f\in L^1(0,T;L^2)\), one has
\begin{align*}
 \|u(t)\|_2&\leq e^{Ct}
          \left(\|u(0)\|_2+\int_0^t\|f(s)\|_2\,ds\right),\\
 C&=\tfrac12\sum_j\|\partial_j A^j\|_\infty+
                      \left\|\tfrac12(B+B^*)\right\|_\infty .
\end{align*}
The matrix norms are operator norms. In particular the estimate
applies to differences, yielding uniqueness and stability.
\end{theorem}

\begin{proof}
For \(e=|u|^2/2\), take the real inner product with \(u\).
The spatial terms satisfy
\[
 \operatorname{Re}(u^*A^j\partial_j u)
    =\tfrac12\partial_j(u^*A^j u)
        -\tfrac12u^*(\partial_j A^j)u .
\]
Thus the flux is \(\tfrac12u^*A^j u\), and the remaining terms in
\(e_t+\operatorname{div}(\text{flux})\) are bounded above by
\(2Ce+|u||f|\). The coefficients \(A^j\) themselves can grow
linearly in \(x\). In fact bounded spatial derivatives and
continuity of \(A^j(t,0)\) give an \(L>0\) with
\[
 \left\|\sum_j\nu_jA^j(s,x)\right\|\leq L(1+|x|)
       \qquad(|\nu|=1,\ 0\leq s\leq T).
\]
Fix a terminal time \(t\) and a radius \(R>0\). On \([0,t]\)
integrate over the ball of radius
\(r(s)=(R+1)e^{L(t-s)}-1\) about zero.
Its outward normal speed is \(-L(1+r(s))\).
In the derivative of the ball energy the net boundary term
\(-u^*(\sum_j\nu_jA^j)u/2+r'(s)e\) is nonpositive.
The integral of the interior forcing term is bounded by
\(\|f(s)\|_2\sqrt{2E(s)}\), where \(E(s)\) is the ball energy.
Hence \(E'(s)\leq 2CE(s)+\|f(s)\|_2\sqrt{2E(s)}\).
The regularized square root gives the integrated estimate on the
moving ball by Gronwall. The initial ball norm is at most
\(\|u(0)\|_2\). Letting \(R\) grow to infinity yields the full
norm estimate. The argument needed no global flux integrability
and proves finiteness of \(\|u(t)\|_2\). Applying it to the
difference gives the claims.
\end{proof}

This is the symmetric flux computation in
\cite[Section 7.3.2]{Evans},
\cite[Chapter 5, Section 3]{John}, and
\cite[Chapter 6, Section 5]{TaylorI}. The moving ball keeps the
linear-growth case within that computation.

\section{Whole-space existence for symmetric systems}\label{pde:W07}

\begin{theorem}[Strong Sobolev well-posedness]\label{thm:pde-symmetric-hs}
Let \(s\) be an integer with \(s>n/2+1\).
In the system of Theorem~\ref{thm:pde-symmetric-energy}, let
\(A^j=A^{j*}\) and \(B\) be smooth, with their spatial derivatives
through order \(s\), including order zero, bounded uniformly on
\([0,T]\times\mathbb R^n\).
For \(g\in H^s(\mathbb R^n)\), \(f\in C([0,T];H^s)\), there is
a unique solution
\[
 u\in C([0,T];H^s)\cap C^1([0,T];H^{s-1}),\qquad u(0)=g .
\]
There is a constant \(C_s\) determined by the stated coefficient
bounds, \(s\), and \(n\), with
\[
 \|u(t)\|_{H^s}
 \leq e^{C_s t}\left(\|g\|_{H^s}
                         +\int_0^t\|f(r)\|_{H^s}\,dr\right).
\]
Here the squared norm is the sum of the squared \(L^2\) norms
of spatial derivatives of orders at most \(s\).
\end{theorem}

\begin{proof}
Put \(P(t)=-\sum_j A^j(t)\partial_j-B(t)\).
For each integer \(0\leq k\leq s\) multiplication by the
coefficients is bounded on \(H^k\), by Leibniz' rule and the
uniform derivative bounds. These multiplier maps act strongly
continuously on each fixed \(H^k\) vector: in each product of a
coefficient derivative and an \(L^2\) derivative of the vector,
apply dominated convergence to the former pointwise and use its
uniform bound. Thus \(P(t)\) is a uniformly bounded, strongly
continuous family from \(H^k\) to \(H^{k-1}\) when \(1\leq k\leq s\).

For \(v\in H^{k+1}\), integrate by parts in the principal term
of \((P(t)v,v)_{H^k}\). The term with \(A^j\partial_jD^\alpha v\)
is controlled by \(\partial_j A^j\) as in the \(L^2\) calculation.
Every commutator term has a positive-order derivative on \(A^j\)
and at most \(|\alpha|\) derivatives on \(v\). The \(B\)-terms
satisfy the same bound by the multiplier estimate. Density of
compactly supported smooth functions in \(H^{k+1}\) justifies
integration by parts. This gives
\[
 \big|\operatorname{Re}(P(t)v,v)_{H^k}\big|
                 \leq C_k\|v\|_{H^k}^2 .
\]
Let \(J_\epsilon\) be convolution with an even compactly supported
smooth approximate identity of mass one. It commutes with
derivatives, is a self-adjoint contraction on \(H^k\), and maps
\(H^k\) boundedly to \(H^{k+1}\) for each fixed \(\epsilon>0\).
Solve the integral ODE on \(H^s\)
\[
 u_\epsilon'=J_\epsilon P(t)J_\epsilon u_\epsilon+
                 J_\epsilon f,\qquad u_\epsilon(0)=J_\epsilon g .
\]
For fixed \(\epsilon\) its coefficient is bounded and strongly
continuous in time on fixed vectors, so Picard iteration on
short time intervals gives a unique \(C^1(H^s)\) solution and
iteration of the intervals covers \([0,T]\).
Pairing against \(u_\epsilon\) in \(H^s\) moves the outside
\(J_\epsilon\) to the second argument. The preceding energy
bound and Gronwall give the displayed estimate uniformly.

The approximation estimate
\[
 \|(J_\epsilon-I)w\|_{H^{k-1}}\leq C\epsilon\|w\|_{H^k}
                    \qquad(1\leq k\leq s)
\]
follows by writing a translation difference as the integral of
its directional derivative, first for smooth functions and then
by density. Write \(P_\epsilon=J_\epsilon P J_\epsilon\).
The decomposition
\[
 P_\epsilon-P_\delta
   =(J_\epsilon-J_\delta)P J_\epsilon
       +J_\delta P(J_\epsilon-J_\delta)
\]
and the multiplier bounds give
\(\|P_\epsilon-P_\delta\|_{H^s\to H^{s-2}}
   \leq C(\epsilon+\delta)\).
Indeed the first term uses the small difference after \(P\)
has lost one derivative, while the second uses it before \(P\).
The difference \(u_\epsilon-u_\delta\) satisfies the regularized
equation with coefficient \(P_\epsilon\) and source
\((P_\epsilon-P_\delta)u_\delta+(J_\epsilon-J_\delta)f\).
The \(H^{s-2}\) energy estimate, the uniform \(H^s\) bound, and
the small initial difference make this family Cauchy in
\(C(H^{s-2})\). Here \(s\geq2\), so this lower norm is a usual
nonnegative-order Sobolev norm. Interpolation,
\(\|w\|_{H^{s-1}}\leq C\|w\|_{H^{s-2}}^{1/2}\|w\|_{H^s}^{1/2}\)
by Fourier Cauchy--Schwarz and norm equivalence, gives convergence
in \(C(H^{s-1})\).

The limit takes values in \(H^s\) and is uniformly bounded there:
at each time extract a weak \(H^s\) subsequence and identify its
limit by the strong lower norm convergence. This also gives
the energy estimate by weak lower semicontinuity. Weak \(H^s\)
continuity follows by applying the same subsequence argument
along any sequence of times tending to a fixed time. The estimate
for \(P_\epsilon-P\) in \(H^s\to H^{s-2}\), obtained by the same
decomposition with the identity in place of \(J_\delta\), lets
the integral equation pass in \(H^{s-2}\). Thus the limit is
\(C^1(H^{s-2})\) and solves the PDE there. This intermediate
class has \(u\in C(H^1)\) and \(u_t\in C(L^2)\).
With bounded coefficients, the whole-space \(L^2\) identity is
legitimate for this class by spatial \(H^1\) approximation in
the integrated spatial terms. Its energy bound gives uniqueness
even in the intermediate class.

To obtain strong \(H^s\) continuity, fix \(\tau\in[0,T]\).
Repeat the construction with datum \(u(\tau)\) both forward
and backward in time. Backward time replaces \(P\) and \(f\)
by their negatives at reversed times. The absolute energy
bound and all regularity estimates still apply. Uniqueness
identifies both restarted limits with \(u\) on their intervals.
Their weak lower semicontinuity estimates give
\[
 \|u(t)\|_{H^s}\leq e^{C_s|t-\tau|}
 \left(\|u(\tau)\|_{H^s}
           +\int_{\min(t,\tau)}^{\max(t,\tau)}
                            \|f(r)\|_{H^s}\,dr\right).
\]
The resulting upper limit inequality and weak continuity imply
strong continuity in the Hilbert norm. Finally \(P(t)u(t)+f(t)\)
is continuous in \(H^{s-1}\), by the strong continuity of the
multiplier action and the strong \(H^s\) continuity of \(u\).
The integral equation therefore holds in \(H^{s-1}\) and gives
the asserted strong time derivative. The estimate for differences
also gives continuous dependence.
\end{proof}

Compare the whole-space weak construction in
\cite[Section 7.3.2]{Evans}, the high-order energy argument in
\cite[Chapter 5, Section 3]{John}, and the Friedrichs commutator
in \cite[Section 12.3.1]{RR}.
\chapter{Classical Linear Integral Equations}
\label{chap:pde-integral}

The compact Fredholm equation and the triangular Volterra equation
use different mechanisms. In the first the obstruction is a
finite-dimensional kernel. In the second the geometry of an
ordered simplex produces a factorial bound for every power.

\section{Continuous kernels and the Fredholm alternative}\label{pde:I01}

\begin{theorem}[Compact-kernel Fredholm bridge]\label{thm:pde-fredholm-kernel}
Let \(K\) be a compact Hausdorff space, \(\mu\) a finite positive
Borel measure of full support, and \(k\in C(K\times K;\mathbb C)\).
The operator
\[
 Tf(x)=\int_K k(x,y)f(y)\,d\mu(y),\qquad f\in C(K;\mathbb C),
\]
is compact on \(C(K;\mathbb C)\), with
\(\|T\|\leq\mu(K)\|k\|_\infty\).
For \(\lambda\ne0\) let \(F_\lambda=\lambda I-T\) and write
\(F_\lambda'\) for its complex-linear Banach transpose.
Its range is closed, and
\[
 \dim\ker F_\lambda
     =\operatorname{codim}\operatorname{Ran}F_\lambda
     =\dim\ker F_\lambda'<\infty .
\]
The equation \(F_\lambda f=g\) is solvable exactly when
\(\ell(g)=0\) for every \(\ell\in\ker F_\lambda'\).
When solvable its solutions differ by \(\ker F_\lambda\).
In particular either \(F_\lambda\) is invertible, or its kernel
and transpose kernel both have positive dimension.
\end{theorem}

\begin{proof}
The integral is defined for each \(x\) by Borel measurability and
boundedness of the integrand. The algebra of finite sums of
functions \(a(x)b(y)\), with \(a,b\in C(K;\mathbb C)\), contains
the constants, is invariant under conjugation, and separates
points of the compact Hausdorff product. By Stone--Weierstrass
there are such kernels \(k_m\) converging uniformly to \(k\)
\cite[Unit U1]{P102}. Each corresponding \(T_m\) has range
in the span of the finitely many \(a\)'s, with coefficient
functionals \(f\mapsto\int_K b(y)f(y)\,d\mu(y)\).
Thus \(T_m f\) is continuous. The integral estimate gives
\[
 \sup_x|Tf(x)-T_m f(x)|\leq
       \mu(K)\|k-k_m\|_\infty\|f\|_\infty .
\]
It follows first that \(Tf\) is continuous and then that \(T\)
is the operator-norm limit of finite-rank operators. The image
of the unit ball is totally bounded, by choosing one approximant
and a finite net in its bounded finite-dimensional image.
Completeness of \(C(K;\mathbb C)\) gives relative compactness.
This argument uses finiteness of \(\mu\) and Borel integrability,
without needing measure regularity.

Apply the compact Riesz--Schauder alternative to
\(I-\lambda^{-1}T\), as in \cite[Part II, compact operators]{D16}.
Its range is closed and its kernel and cokernel have equal finite
dimension. The kernel of \(F_\lambda'\) is the annihilator of
the range of \(F_\lambda\), hence the continuous dual of its
finite-dimensional quotient. This gives the dimension equality.
If \(g\) is outside the range, its nonzero quotient class has a
continuous functional nonzero on it, and lifting that functional
separates \(g\) from the range. This proves necessity and
sufficiency of the stated compatibility condition. The affine
solution description is immediate by subtracting two solutions.
In the zero-kernel case the range has codimension zero, and the
bounded inverse theorem yields the inverse. The empty compact
space, if allowed, gives the same statement on the zero space.
\end{proof}

The continuous-kernel compactness theorem in
\cite[Section 2.5, Theorem 2.27]{Kress} is stated for the classical
Euclidean compact set setting. The argument above also covers a
general compact Hausdorff space. The transpose convention is
\(F_\lambda'=\lambda I-T'\) on the complex-linear dual.
The Hilbert adjoint convention would conjugate \(\lambda\).

\section{The regular Volterra resolvent}\label{pde:I02}

\begin{theorem}[Triangular kernel and factorial convergence]
\label{thm:pde-volterra}
Let \(a\leq b\), let \(k\) be continuous on
\(\Delta=\{(t,s):a\leq s\leq t\leq b\}\), and let
\(g\in C([a,b])\), over either scalar field. For every scalar
\(\lambda\) the equation
\[
 u(t)=g(t)+\lambda\int_a^t k(t,s)u(s)\,ds
\]
has a unique continuous solution. Put \(M=\|k\|_\infty\) and
\[
 Vu(t)=\int_a^t k(t,s)u(s)\,ds .
\]
The solution operator is
\[
 (I-\lambda V)^{-1}=\sum_{m=0}^\infty\lambda^m V^m,
 \qquad \|(I-\lambda V)^{-1}\|\leq e^{|\lambda|M(b-a)} .
\]
For \(k_1=k\) and
\(k_{m+1}(t,s)=\int_s^t k(t,r)k_m(r,s)\,dr\),
the continuous resolvent kernel
\(R_\lambda(t,s)=\sum_{m\geq1}\lambda^{m-1}k_m(t,s)\)
gives
\(u(t)=g(t)+\lambda\int_a^t R_\lambda(t,s)g(s)\,ds\).
\end{theorem}

\begin{proof}
If \(a=b\), every indicated integral is zero and \(u=g\).
The kernel recurrence then has \(k_m=0\) for \(m\geq2\),
so the formulas still hold. Suppose \(a<b\).
Continuity of \(Vu\) follows by putting \(s=a+\theta(t-a)\)
in the integral, giving an integral over a fixed compact
parameter interval. Iterated integration gives
\[
 |V^m u(t)|\leq M^m\|u\|_\infty\,\frac{(t-a)^m}{m!}.
\]
The domain is \(a\leq s_m\leq\cdots\leq s_1\leq t\).
Its volume follows by successively integrating powers, or by
partitioning the \(m\)-dimensional cube up to null boundaries
according to the orders of its coordinates. Hence the operator
series converges in norm with the stated exponential bound.
Multiplying a partial sum on either side by \(I-\lambda V\)
leaves \(I-\lambda^{m+1}V^{m+1}\), whose last term tends to
zero. This proves a two-sided inverse and uniqueness. The
successive approximations \(u_{j+1}=g+\lambda V u_j\) converge
from any continuous starting value because their error at step
\(j\) is \((\lambda V)^j(u_0-u)\).

By the change \(r=s+\theta(t-s)\), continuity of \(k_m\) and \(k\)
implies continuity of \(k_{m+1}\) on the closed triangle.
The recurrence gives
\[
 |k_m(t,s)|\leq M^m\,\frac{(t-s)^{m-1}}{(m-1)!}
\]
by induction. Fubini shows
\(V^m u(t)=\int_a^t k_m(t,s)u(s)\,ds\): one inserts the
formula for \(V^m\) into \(V(V^m u)\) and reverses the
integration on \(a\leq s\leq r\leq t\).
The exponential majorant now gives absolute uniform convergence
of \(R_\lambda\) on \(\Delta\), including \(\lambda=0\) with
zeroth power equal to one. Interchanging the uniformly convergent
sum with the integral proves the resolvent representation.
Equivalently the recurrence gives
\(R_\lambda(t,s)=k(t,s)+\lambda\int_s^t
k(t,r)R_\lambda(r,s)\,dr\) by the same uniform convergence.
\end{proof}

See \cite[Section 3.3, Theorem 3.10]{Kress},
\cite[Section 1.2, Lemma 1.2.1 and Theorem 1.2.3]{Brunner}, and
the operator-valued extension in \cite{D16}.

\backmatter
\begingroup
\small
\begin{thebibliography}{99}
\bibitem{Evans} L.~C. Evans, \emph{Partial Differential Equations}, 2nd ed.,
American Mathematical Society, 2010.
\bibitem{AG} D.~H. Armitage and S.~J. Gardiner, \emph{Classical Potential
Theory}, Springer, 2001.
\bibitem{GT} D. Gilbarg and N.~S. Trudinger, \emph{Elliptic Partial
Differential Equations of Second Order}, 2nd ed., corrected reprint, Springer, 1998.
\bibitem{Grisvard} P. Grisvard, \emph{Elliptic Problems in Nonsmooth
Domains}, Pitman, 1985.
\bibitem{TaylorI} M.~E. Taylor, \emph{Partial Differential Equations I:
Basic Theory}, 3rd ed., Springer.
\bibitem{TaylorII} M.~E. Taylor, \emph{Partial Differential Equations II:
Qualitative Studies of Linear Equations}, 3rd ed., Springer.
\bibitem{RR} M. Renardy and R.~C. Rogers, \emph{An Introduction to Partial
Differential Equations}, 2nd ed., Springer, 2004.
\bibitem{Brezis} H. Brezis, \emph{Functional Analysis, Sobolev Spaces and
Partial Differential Equations}, Springer, 2011.
\bibitem{Lieberman} G.~M. Lieberman, \emph{Second Order Parabolic
Differential Equations}, World Scientific, 1996, revised printing.
\bibitem{EngelNagel} K.-J. Engel and R. Nagel, \emph{One-Parameter
Semigroups for Linear Evolution Equations}, Springer, 2000.
\bibitem{John} F. John, \emph{Partial Differential Equations}, 3rd ed.,
Springer, 1978.
\bibitem{Strauss} W.~A. Strauss, \emph{Partial Differential Equations:
An Introduction}, 2nd ed., Wiley, 2008.
\bibitem{Kress} R. Kress, \emph{Linear Integral Equations}, 3rd ed.,
Springer, 2014.
\bibitem{Brunner} H. Brunner, \emph{Volterra Integral Equations:
An Introduction to Theory and Applications}, Cambridge University Press, 2017.
\bibitem{P102} MathGenome, \emph{Fourier Analysis, Distributions, Sobolev
Spaces, and Approximation}, unified notes P1-02, 2026.
\bibitem{P106} MathGenome, \emph{Variational Methods, Convex Analysis,
and Optimality}, unified notes P1-06, 2026.
\bibitem{D12} MathGenome, \emph{Complex Analysis}, unified notes D12, 2026.
\bibitem{D13} MathGenome, \emph{Ordinary Differential Equations and
Dynamical Systems}, unified notes D13, 2026.
\bibitem{D16} MathGenome, \emph{Expanded Functional Analysis}, corrected
consolidated notes D16, version a132ac18, 2026.
\end{thebibliography}
\endgroup
\end{document}
